% Proof appendix

\section{Guide to the logic}
\label{sec:appendix-guide}
This appendix isolates the minimal interfaces and implications needed to reach the
HLS classification step from the standing interface packages.  The guide below
follows the proof dependence rather than the raw appendix order: preparatory
subprograms are listed where they first enter the argument.
\begin{enumerate}[label=\textnormal{(\roman*)},leftmargin=*,itemsep=0.4em]
\item \textbf{Standing package bundle.} The standing assumptions are packaged as
  \(\mathbf{2T}_{\mathrm{Main}}\)
  \textnormal{(Definition~\ref{def:2T-Main})}, namely
  \[
  \mathbf{2T}_{\mathrm{LS}}
  + \mathbf{2T}_{\mathrm{RG}}
  + \mathbf{2T}_{\mathrm{Peierls}}
  + \mathbf{2T}_{\mathrm{EBC}}
  + \mathbf{2T}_{\mathrm{char}}^{\mathrm{cal}}
  + \mathbf{2T}_{\mathrm{char}}^{\mathrm{act}}
  + \mathbf{2T}_{\mathrm{sel}}
  + \mathbf{2T}_{\mathrm{sel}}^{\mathrm{delay}}
  + \mathbf{2T}_{\mathrm{conf}}.
  \]
  The record interface \(\mathbf{2T}_{\mathrm{rec}}\) and detector interface
  \(\mathbf{2T}_{\mathrm{det}}\) enter only through the LS summand, while the
  derived BRST package will be added below to form
  \(\mathbf{2T}_{\mathrm{Main}}^{\mathrm{spine}}\).
\item \textbf{Subprogram~1a (derived BRST).} Starting from the Lorentzian SK path
  integral, Subprogram~1a constructs the two-time BRST/bicomplex differentials
  \(Q_t,Q_u\) and their Ward identities on SK--admissible states, thereby
  producing the interface \(\mathbf{2T}_{\mathrm{BRST}}\) on the sufficiently
  local class \(\Fsloc\).  At that point the appendix spine works with
  \[
  \mathbf{2T}_{\mathrm{Main}}^{\mathrm{spine}}
  :=
  \mathbf{2T}_{\mathrm{Main}}+\mathbf{2T}_{\mathrm{BRST}}.
  \]
\item \textbf{Subprogram~2a (Results from the 2T Limited Scope Theorem).}  We import the LS
  parity cocycle, the nonsplitting \(\mathbb Z_2\)-cover, the detector witness,
  and the record-process-chain bicomplex identities.  This isolates the
  abstract two-time obstruction data that later feed the parity-transport step.
\item \textbf{Subprogram~2b (Cosmology theory Peierls bracket and swap invariance).}  We record
  the abstract Peierls package on sufficiently local boundary observables, its
  swap-equivariant Green operators, and the resulting Poisson automorphism
  property of the boundary swap involution.
\item \textbf{Subprogram~2c (Parity transport and physical \(\mathbb Z_2\) grading).}  Combining
  the LS import with the swap-equivariant boundary structures, we identify the
  physical parity involution \(\Gamma\) on first descendants and transport the
  LS square-loop sign to the boundary ambiguity sector.
\item \textbf{Subprogram~3a (selected-line nonvanishing, canonicalization, and characteristic activity).}
  This subsection packages the finite-cutoff canonical odd representative
  together with the same-tick obstruction lane, the delayed selected-line lane,
  the resulting nonvanishing statements, and the characteristic-algebra/activity
  consequences that follow once the odd descendant class survives.
\item \textbf{Subprogram~3b (RG monotone and Hamiltonized odd derivation).}
  The RG package \(\mathbf{2T}_{\mathrm{RG}}\) supplies the monotone and its
  anomalous descent source, while the Peierls and characterization packages are
  used to Hamiltonize the canonical odd representative and establish the local
  physically odd derivation \(\delta_{\Lambda_k}\).
\item \textbf{Subprogram~1b (local net, GNS reconstruction, and implementation template).}
  Subprogram~1b supplies the vacuum Wightman/AQFT representation
  \((\mathcal H,\pi,U,\Omega)\), the net \(O\mapsto\mathfrak A(O)\), and the
  dense implementation core on the Minkowski boundary.  These Hilbert-space
  inputs first enter when the argument passes from the local odd derivation to
  the AQFT/HLS step.
\item \textbf{Subprogram~3c (AQFT net and HLS on the Minkowski boundary).}
  Using \(\mathbf{2T}_{\mathrm{EBC}}\) together with the Subprogram~1b vacuum
  representation, Subprogram~3c extends \(\delta\) to the polynomial core,
  proves implementability of the odd derivation, verifies the HLS interface
  \(\mathbf{2T}_{\mathrm{HLS}}\), and identifies the finite-dimensional
  Weyl-spinor odd charge module on the UV vacuum representation.
\item \textbf{Subprogram~4a (dilatation weight \(\tfrac12\) of the odd charges on the common UV sector).}
  Combining the HLS odd module from Subprogram~3c, the HLS completeness
  statement, and the conformal common-sector package
  \(\mathbf{2T}_{\mathrm{conf}}\), Subprogram~4a places the odd charges on the
  common UV sector and forces their dilatation weight to be \(\tfrac12\).
\item \textbf{Subprogram~4b (same-\(\mathcal H\) conformal closure of the HLS odd sector).}
  With the weight-\(\tfrac12\) normalization from Subprogram~4a in hand,
  Subprogram~4b takes the finite conformal closure of the HLS odd seed,
  proves same-\(\mathcal H\) implementability of that odd conformal closure,
  recovers the degree-zero internal algebra, and identifies the resulting
  operator superalgebra with the standard \(4\)D superconformal algebra.
\item \textbf{Subprogram~4c (optional maximality strengthening).}
  This final section is not part of the base Main theorem.  It imports the
  superconformal closure from Subprogram~4b and studies what additional
  assumptions would force maximal supersymmetry, i.e.\ the optional
  \(N=4\) conclusion at the UV fixed point.
\end{enumerate}

\section{Master packages and interface assumptions}
\label{sec:setup}

\subsection{Master package block}

\subsubsection{Input Interfaces}

\begin{definition}[Record two--time interface $\mathbf{2T}_{\mathrm{rec}}$]
\label{def:2T-rec}
A \emph{record two--time interface} consists of:
\begin{enumerate}[label=\textbf{(R\arabic*)},leftmargin=*,itemsep=0.5em]
\item \textbf{SK and records.} Sets $\SKstates$ (SK states) and $\Sys$ (record states),
  a map $\diag:\Sys\to\SKstates$, and an involution $\swap:\SKstates\to\SKstates$ such that
  $\swap^2=\mathrm{id}_{\SKstates}$ and $\swap\circ\diag=\diag$.
\item \textbf{Two arrows.} A group $T$ and the additive monoid $(\mathbb R_{\ge 0},+,0)$,
  with actions $\mathrm{act}_T:T\times\SKstates\to\SKstates$ and
  $\mathrm{act}_U:\mathbb R_{\ge 0}\times\SKstates\to\SKstates$
  satisfying the action axioms and \emph{weak commutativity} on $\SKstates$:
  \[
  \mathrm{act}_U(u,\mathrm{act}_T(t,x))=\mathrm{act}_T(t,\mathrm{act}_U(u,x)).
  \]
\item \textbf{Strict RG diagonalization.} For each $u>0$ a map
  $\Pi_u:\SKstates\to\Sys$ and a record RG action $R_u:\Sys\to\Sys$ such that
  \[
  \mathrm{act}_U(u,x)=\diag(\Pi_u(x)),\qquad
  \Pi_u(\diag r)=R_u(r)\quad(u>0),
  \]
  with $R_0=\mathrm{id}$ and $R_{u_2+u_1}=R_{u_2}\circ R_{u_1}$.
\item \textbf{Chosen RG sector (reachability).} A section
  $\mathrm{sec}_\Pi:\Sys\to\SKstates$ with $\Pi_1(\mathrm{sec}_\Pi(r))=r$.
\item \textbf{No mixed anomaly at the boundary of squares.}
  A bisimplicial set $\mathcal C^{\mathrm{SK}}_{m,n}$ with horizontal and vertical face maps
  $d_i^H,d_j^V$ satisfying simplicial relations in each direction and mixed
  commutation $d_i^H d_j^V=d_j^V d_i^H$~\cite{GoerssJardine1999Simplicial}.
\item \textbf{Record reflection.} Diagonal cancellation:
  \[
  \diag(r)=\diag(r') \Rightarrow r=r'.
  \]
\end{enumerate}
\end{definition}

\begin{remark}
The discrete RG-step case is recovered by restricting $u\in\mathbb N\subset\mathbb R_{\ge 0}$.
\end{remark}

\begin{definition}[Detector interface $\mathbf{2T}_{\mathrm{det}}$]
\label{def:2T-det}
A \emph{detector interface} is any additional structure that implies the
$\mathbb Z_2$--action on the difference object $\Diff=\mathrm{Cofib}(\diag)$ is
nontrivial, i.e.\ $\flip\neq \mathrm{id}_{\Diff}$ or equivalently there exists
$y\in\Diff$ with $\flip(y)\neq y$.
\end{definition}

\begin{definition}[Limited Scope package $\mathbf{2T}_{\mathrm{LS}}$]
\label{def:2T-LS}
Define the \emph{2T Limited Scope package} by
\[
\mathbf{2T}_{\mathrm{LS}} := \mathbf{2T}_{\mathrm{rec}} + \mathbf{2T}_{\mathrm{det}}.
\]
\end{definition}

\begin{definition}[Two independent $\mathbb Z_2$ gradings]
\label{def:two-gradings}
We distinguish:
\begin{enumerate}[label=\textnormal{(\alph*)},leftmargin=*,itemsep=0.2em]
\item \emph{Cohomological/BRST parity} $|\cdot|_{\mathrm{ch}}\in\mathbb Z_2$ on the
  cochain algebra $\Fsloc=(\Fsloc)_{\bar 0}\oplus(\Fsloc)_{\bar 1}$, with
  $|Q_t|_{\mathrm{ch}}=|Q_u|_{\mathrm{ch}}=1$.
  Accordingly, $H^0(\Fsloc,Q_t)$ and $H^1(\Fsloc,Q_t)$ denote even/odd cohomology.
\item \emph{Physical parity} $|\cdot|_{\Gamma}\in\mathbb Z_2$ induced by the involution
  $\Gamma$ (SK swap), giving a decomposition into $\Gamma$--even/odd parts
  $\Fsloc={\Fsloc}^{+}\oplus{\Fsloc}^{-}$.
\end{enumerate}
Oddness in the AQFT/HLS steps always refers to $|\cdot|_{\Gamma}$.
\end{definition}

\begin{definition}[RG--monotone interface $\mathbf{2T}_{\mathrm{RG}}$]
\label{def:2T-RG}
An \emph{RG--monotone interface} consists of:
\begin{enumerate}[label=\textbf{(M\arabic*)},leftmargin=*,itemsep=0.5em]
\item \textbf{Monotone and representative.} A real monotone $M(u)$ near the UV
  ceiling $u_0$, a family of SK--admissible states $\rho_u$, and a quasi--local
  family $\mathcal O^{(0)}(u)\in\Fsloc$ with
  \[
  M(u)=\langle \mathcal O^{(0)}(u)\rangle_{\rho_u},
  \qquad
  Q_t\mathcal O^{(0)}(u)=0,
  \qquad
  [\mathcal O^{(0)}(u_0)]\neq 0\in H^0(\Fsloc,Q_t).
  \]
\item \textbf{RG Ward identity (odd form, refined).} On the monotone representative
  $\mathcal O^{(0)}(u)$, decompose the $u$--differential as
  \[
  Q_u = Q_u^{\mathrm{phys}} + Q_u^{\mathrm{scheme}},
  \]
  where $Q_u^{\mathrm{scheme}}$ encodes scheme/reparametrisation effects
  (local counterterm and coupling redefinition ambiguities). We assume that
  $Q_u^{\mathrm{phys}}$ and $Q_u^{\mathrm{scheme}}$ are odd derivations
  on $\Fsloc$.
  We additionally assume $Q_u^{\mathrm{phys}}$ and $Q_u^{\mathrm{scheme}}$ are
  \emph{presheaf derivations}: for every double cone $O$ they restrict to odd
  derivations
  \[
  Q_u^{\mathrm{phys}}:\Fsloc(O)\to\Fsloc(O),\qquad
  Q_u^{\mathrm{scheme}}:\Fsloc(O)\to\Fsloc(O).
  \]
  Finally, we assume they anticommute with $Q_t$ on this class (so
  $Q_u^{\mathrm{phys}}\mathcal O^{(0)}$ is $Q_t$--closed whenever
  $Q_t\mathcal O^{(0)}=0$), and that for $u$ near $u_0$,
  \[
  \frac{d}{du}M(u)
  = -\big\langle (Q_u^{\mathrm{phys}}\mathcal O^{(0)})(u)\big\rangle_{\rho_u},
  \qquad
  \big\langle (Q_u^{\mathrm{scheme}}\mathcal O^{(0)})(u)\big\rangle_{\rho_u}=0,
  \]
  up to expectation values of $Q_t$--exact terms (which vanish for
  SK--admissible states by $\mathbf{2T}_{\mathrm{BRST}}$-\textbf{(B2)}).
\par
  \textbf{Cohomological triviality of scheme terms on the monotone sector.}
  For every $X\in{\Fsloc}_{\mathrm{mon}}$ with $Q_tX=0$, we assume
  \[
  [\,Q_u^{\mathrm{scheme}}X\,]=0\quad\text{in}\quad H^1(\Fsloc,Q_t).
  \]
\par
  (Optional packaging.) Equivalently, one may encode the same information by an
  even derivation $L_{\rm RG}$ on $\Fsloc$ (Callan--Symanzik plus anomaly/contact
  terms) with
  $
  \tfrac{d}{du}M(u)=\langle (L_{\rm RG}\mathcal O^{(0)})(u)\rangle_{\rho_u}.
  $
\item \textbf{Anomalous first descent (no Ward trivialization).}
  There exist sufficiently local raw first-descendant potentials
  \(\Xi^{(1)}(u)\in\Fsloc\)
  and an anomaly insertion $\mathcal A^{(1)}(u)\in(\Fsloc)_{\bar 1}$ such that
  \begin{equation}
  \label{eq:anom-descent}
  Q_u^{\mathrm{phys}}\mathcal O^{(0)}(u)\ +\ Q_t\Xi^{(1)}(u)\ =\ \mathcal A^{(1)}(u),
  \qquad
  Q_t\mathcal A^{(1)}(u)=0.
  \end{equation}
  We interpret \eqref{eq:anom-descent} as an identity in the locality quotient
  (Definition~\ref{def:locality-quotient}); in particular, $\mathcal A^{(1)}(u)$
  represents the $Q_t$--cohomology class
  $\big[\,Q_u^{\mathrm{phys}}\mathcal O^{(0)}(u)\,\big]\in H^1(\Fsloc,Q_t)$
  (Definition~\ref{def:A1-class}).
  For SK--admissible states, $\langle Q_t X\rangle_{\rho_u}=0$ implies
  $\langle (Q_u^{\mathrm{phys}}\mathcal O^{(0)})(u)\rangle_{\rho_u}
  =\langle \mathcal A^{(1)}(u)\rangle_{\rho_u}$ and hence
  $\tfrac{d}{du}M(u)=-\langle \mathcal A^{(1)}(u)\rangle_{\rho_u}$.
  This replaces an anomaly--free descent
  $Q_u^{\mathrm{phys}}\mathcal O^{(0)}+Q_t\Xi^{(1)}=0$, which would force
  $\tfrac{d}{du}M(u)=0$ by the $Q_t$ Ward identity.
  We furthermore require that the chosen first-descendant potential at the UV ceiling
  is not a c-number on the vacuum sector (i.e.\ not constant on the connected vacuum
  component of the reduced phase space).
\item \textbf{Nontrivial UV slope.} $\partial_u M(u)\rvert_{u_0}<0$ away from fixed
  points / exactly marginal loci.
\end{enumerate}
\end{definition}

\begin{definition}[Monotone sector of sufficiently local observables]
\label{def:Fsloc-mon}
Fix the UV ceiling $u_0$ and choose $\varepsilon>0$.  For each double cone $O$,
define ${\Fsloc}_{\mathrm{mon}}(O)\subset\Fsloc(O)$ to be the smallest
$\mathbb R$--subspace such that:
\begin{enumerate}[label=\textnormal{(\roman*)},leftmargin=*,itemsep=0.2em]
\item for all $u\in(u_0-\varepsilon,u_0+\varepsilon)$, the chosen monotone representatives
  $\mathcal O^{(0)}(u)$ restricted to $O$ lie in ${\Fsloc}_{\mathrm{mon}}(O)$;
\item the associated anomaly insertions $\mathcal A^{(1)}(u)$ and chosen raw first-descendant
  potentials \(\Xi^{(1)}(u)\) (whenever defined) restricted to $O$ lie in ${\Fsloc}_{\mathrm{mon}}(O)$;
\item ${\Fsloc}_{\mathrm{mon}}(O)$ is $\Gamma$--invariant and $Q_t$--stable:
  $\Gamma({\Fsloc}_{\mathrm{mon}}(O))\subset{\Fsloc}_{\mathrm{mon}}(O)$ and
  $Q_t({\Fsloc}_{\mathrm{mon}}(O))\subset{\Fsloc}_{\mathrm{mon}}(O)$.
\end{enumerate}
Here $\Gamma$ denotes the swap-induced involution on $\Fsloc$ introduced in
Section~\ref{sec:parity-bridge}.
We write ${\Fsloc}_{\mathrm{mon}}:=\bigcup_O {\Fsloc}_{\mathrm{mon}}(O)$.

\medskip
\end{definition}

\begin{definition}[Locality quotient relation]
\label{def:locality-quotient}
Fix a double cone $O$.  On $\Fsloc(O)$ define an equivalence relation $\simeq$
generated by the following elementary moves:
\begin{enumerate}[label=\textnormal{(\alph*)},leftmargin=*,itemsep=0.2em]
\item $F \mapsto F+Q_t X$ for any $X\in\Fsloc(O)$;
\item $F \mapsto F+(Q_u^{\mathrm{scheme}}Y)$ for any $Y\in{\Fsloc}_{\mathrm{mon}}(O)$;
\item $F \mapsto F+\partial_\mu K^\mu$ for any sufficiently local current $K^\mu$
  supported in $O$ (so the term vanishes after smearing in double cones);
\item $F \mapsto F+\Delta_{\mathrm{ct}}$ where $\Delta_{\mathrm{ct}}\in\Fsloc(O)$
  is a local counterterm coboundary (supported in $O$) representing a scheme change.
\end{enumerate}
Write $[F]_{\mathrm{loc}}$ for the $\simeq$--class of $F$.
\end{definition}

\begin{definition}[Anomaly insertion as the $Q_t$--cohomology class]
\label{def:A1-class}
Assume $Q_t\mathcal O^{(0)}(u)=0$ and that $\{Q_t,Q_u^{\mathrm{phys}}\}\mathcal O^{(0)}(u)=0$
for $u$ near $u_0$ (equivalently, $Q_t(Q_u^{\mathrm{phys}}\mathcal O^{(0)}(u))=0$).
Define the anomaly class
\[
[\mathcal A^{(1)}(u)]\ :=\ \big[\,Q_u^{\mathrm{phys}}\mathcal O^{(0)}(u)\,\big]\ \in\ H^1(\Fsloc,Q_t),
\]
and choose a sufficiently local representative $\mathcal A^{(1)}(u)\in(\Fsloc)_{\bar 1}$
of this class.  Scheme/reparametrisation contributions are absorbed into the locality quotient
$\simeq$ (Definition~\ref{def:locality-quotient}) and are assumed cohomologically
trivial on ${\Fsloc}_{\mathrm{mon}}$ by Definition~\ref{def:2T-RG}-\textbf{(M2)}.
\end{definition}

\begin{definition}[Peierls interface $\mathbf{2T}_{\mathrm{Peierls}}$]
\label{def:2T-cosmo-P}
A \emph{Peierls interface} consists of the following properties:
\begin{enumerate}[label=\textbf{(P\arabic*)},leftmargin=*,itemsep=0.5em]
\item \textbf{Renormalized SK differentiability.} The renormalized SK generating
  functional $W[J_+,J_-]$ exists and is twice differentiable near $J_a=0$ on a
  specified space of boundary sources (including the gravitational sources).
\item \textbf{Causal Green structure.} The gauge--fixed linearized renormalized
  bulk operator is Green--hyperbolic with the chosen AdS boundary conditions, so
  advanced/retarded propagators exist. The induced boundary retarded kernel
  \[
  G_{\mathrm{ret}}(x,y) := \frac{\delta^2 W}{\delta J_a(x)\,\delta J_r(y)}\Big\rvert_{J_a=0}
  \]
  has retarded support, and $G_{\mathrm{adv}}$ is defined analogously.
\item \textbf{Renormalized symplectic form.} The renormalized covariant
  symplectic current $\omega_{\mathrm{ren}}$ exists and yields a finite conserved
  symplectic form $\Omega_{\mathrm{ren}}$ compatible with the causal propagator
  $\Delta:=G_{\mathrm{ret}}-G_{\mathrm{adv}}$.
\item \textbf{Gauge reduction and independence on invariants.} Pure gauge
  directions are exactly the kernel of $\Omega_{\mathrm{ren}}$, so
  $\Omega_{\mathrm{ren}}$ descends to a symplectic form on the reduced phase
  space. The induced Peierls bracket on gauge--invariant functionals is
  independent of the gauge--fixing choice (equivalently, it is canonical on the
  reduced phase space).
\item \textbf{Closed locality class.} There exists a specified algebra
  $\Fsloc$ of gauge--invariant ``sufficiently local'' boundary functionals,
  together with a locality presheaf $O\mapsto \Fsloc(O)$ (closed under smearing
  with smooth compact support kernels), such that the Peierls bracket is
  well-defined and satisfies Leibniz/Jacobi on $\Fsloc$ (and restricts to each
  $\Fsloc(O)$).
\item \textbf{Pointer-real swap equivariance.} There exists a swap-induced involution
  $\Gamma_*$ on the chosen linearized field, test-source, and boundary-data spaces such that
  \[
  P_{\mathrm{gf,ren}}\Gamma_*=\Gamma_*P_{\mathrm{gf,ren}},
  \qquad
  B_\Lambda\Gamma_*=\Gamma_*B_\Lambda,
  \]
  the source insertion and boundary extraction maps are $\Gamma_*$--equivariant,
  and the pairing used in the Peierls formula is $\Gamma_*$--invariant.
\end{enumerate}
\end{definition}

\begin{remark}[Meaning of $\mathbf{2T}_{\mathrm{Peierls}}$]
The package \(\mathbf{2T}_{\mathrm{Peierls}}\) means that the gauge-fixed
linearized bulk problem admits the causal Green
structure needed for advanced and retarded propagators, with admissible
boundary and junction conditions; that the renormalized symplectic current
yields a finite reduced symplectic form with no extra degeneracies beyond pure
gauge; that the resulting Peierls bracket is well-defined on the
sufficiently local boundary class; and that the pointer-real swap involution
acts equivariantly on the linearized renormalized IBVP and its Peierls pairing.
These are the healthiness and symmetry properties encoded abstractly by
\textbf{(P2)}--\textbf{(P6)}.
\end{remark}

\begin{definition}[Energy bounds and covariance interface $\mathbf{2T}_{\mathrm{EBC}}$ (EB1--EB3 + Poincar\'e covariance)]
\label{def:2T-EBC}
An \emph{energy bounds and covariance interface} $\mathbf{2T}_{\mathrm{EBC}}$
consists of the following data and properties on a Minkowski UV chart (uniformly
in the UV cutoff):
\begin{enumerate}[label=\textnormal{(0\roman*)},leftmargin=*,itemsep=0.4em]
\item a vacuum state $\varpi$ on the local Borchers--Uhlmann net generated by
  $\Fsloc(O)$, with vacuum GNS triple $(\mathcal H,\pi,\Omega)$;
\item a strongly continuous positive-energy unitary representation $U_g$ of the
  Poincar\'e group implementing an induced action $\alpha_g$ on $\Fsloc$ and
  fixing the vacuum vector, $U_g\Omega=\Omega$, with time-translation generator
  $H:=P^0$.
\end{enumerate}
In this vacuum representation, the following estimates and covariance identities
hold:
\begin{enumerate}[label=\textnormal{(\roman*)},leftmargin=*,itemsep=0.4em]
\item \textbf{EB1 (local generators).} For every double cone
  $O\subset\mathbb R^{1,3}$ and every $F\in\Fsloc(O)$, there exist $s\ge 0$ and
  $C_{F,O}>0$ such that on the vacuum core $\mathcal D_0$ (Definition~\ref{def:A-fin})
  one has
  \[
  \|\pi(\Phi(F))\psi\|\ \le\ C_{F,O}\,\|(1+H)^s\psi\|
  \qquad(\psi\in\mathcal D_0).
  \]
\item \textbf{EB2 (Sobolev continuity).} The causal operator $\Delta$ extends to a
  bounded map between the Sobolev spaces controlling functional derivatives on
  each $O$.
\item \textbf{EB3 ($\delta$ energy bound; dependence on the generator).}
  Fix a first-descendant potential $\mathcal J\in\Fsloc(O_{\mathcal J})$ and
  define $\delta_{\mathcal J}(F):=\{\mathcal J,F\}_{\mathrm P}$ on $\Fsloc$.
  (We extend $\delta_{\mathcal J}$ to the Borchers--Uhlmann core by
  $\delta_{\mathcal J}(\Phi(F)):=\Phi(\delta_{\mathcal J}F)$ and the graded Leibniz rule.)
  Then for every double cone $O$ there exist $s\ge 0$ and a constant
  $C_{O,\mathcal J}>0$ such that for all
  $A\in\mathfrak A_{\mathrm{fin}}\cap\mathfrak A_{\mathrm{BU}}(O)$ one has
  \[
  \|\pi(\delta_{\mathcal J}(A))\Omega\|\ \le\ C_{O,\mathcal J}\,
  \|(1+H)^s\,\pi(A)\Omega\|.
  \]
\item \textbf{Covariance (Peierls bracket).} In the Minkowski chart, the Peierls
  bracket is Poincar\'e covariant:
  \[
  \alpha_g(\{F,G\}_{\mathrm P})\ =\ \{\alpha_gF,\alpha_gG\}_{\mathrm P}
  \qquad(F,G\in\Fsloc),
  \]
  where $\alpha_g$ is the induced action on $\Fsloc$ by pushforward of test
  kernels and fields.
\end{enumerate}
\end{definition}

\begin{remark}[Uniformity in EB3]
In the proof of EB3, the bound is \emph{uniform in the cutoff} and
\emph{uniform in the tested observable} once $F\in\Fsloc(O)$ is measured by the
same family of locality seminorms (then transported to energy norms via EB1/EB2).
However, the constant depends on the chosen Hamiltonian generator through finitely
many Sobolev/seminorm controls on its first functional derivative. Accordingly,
one may either (i) fix $\mathcal J$ once and for all and write $C_{O,\mathcal J}$,
or (ii) restrict to an admissible class of $\mathcal J$ with uniform seminorm bounds,
in which case one may choose $C_O$ uniformly within that class.
\end{remark}

\begin{remark}
The interface $\mathbf{2T}_{\mathrm{EBC}}$ is formulated on a Minkowski UV chart
with the cutoff-uniform energy bounds, Poincar\'e covariance, and package theorem
needed below.
\end{remark}

\begin{definition}[EB3-admissible first-descendant potentials]
\label{def:EB3-adm-J}
Fix $\sigma>\frac52$ and a double cone $O_{\mathcal J}$. A first-descendant potential
$\mathcal J\in\Fsloc(O_{\mathcal J})$ is called \emph{EB3-admissible} if its first
functional derivative exists and satisfies a cutoff-uniform Sobolev bound
\[
\mathcal J^{(1)}\in H^\sigma_0(O_{\mathcal J}),
\qquad
\|\mathcal J^{(1)}\|_{H^\sigma(O_{\mathcal J})}\ \le\ C^{(\sigma)}_{\mathcal J,O_{\mathcal J}},
\]
where $C^{(\sigma)}_{\mathcal J,O_{\mathcal J}}$ depends only on the defining locality
seminorms/kernels of $\mathcal J$ and is independent of the UV cutoff for all
sufficiently large cutoff.
\end{definition}

\begin{lemma}[Locality $\Rightarrow$ EB3-admissibility of $\mathcal J$]
\label{lem:J-EB3-adm}
Assume $\mathbf{2T}_{\mathrm{Peierls}}$-\textbf{(P5)} and let
$\mathcal J\in\Fsloc(O_{\mathcal J})$. Then $\mathcal J$ is EB3-admissible in the
sense of Definition~\ref{def:EB3-adm-J}.
\end{lemma}

\begin{proof}
By $\mathbf{2T}_{\mathrm{Peierls}}$-\textbf{(P5)}, the sufficiently local class
$\Fsloc(O_{\mathcal J})$ is closed under the functional-derivative operations
used in the Peierls construction (in particular, first functional derivatives
exist and are supported in the same double cone).
Thus the first functional derivative $\mathcal J^{(1)}$ exists, is supported in
$O_{\mathcal J}$, and is controlled by finitely many locality seminorms of
$\mathcal J$ uniformly in the UV cutoff (this is the locality input used in
the EB3 estimate).
Translating these locality seminorm bounds into Sobolev norms yields, for the
fixed $\sigma>\frac52$ of Definition~\ref{def:EB3-adm-J},
\[
\mathcal J^{(1)}\in H^\sigma(O_{\mathcal J}),\qquad
\|\mathcal J^{(1)}\|_{H^\sigma(O_{\mathcal J})}\le C^{(\sigma)}_{\mathcal J,O_{\mathcal J}},
\]
with $C^{(\sigma)}_{\mathcal J,O_{\mathcal J}}$ independent of the cutoff.
Since $\mathcal J^{(1)}$ is supported in $O_{\mathcal J}$, it lies in
$H^\sigma_0(O_{\mathcal J})$, proving EB3-admissibility.
\end{proof}

\begin{definition}[Weak polarized characterization interface \(\mathbf{2T}_{\mathrm{char}}^{\mathrm{weak}}\)]
\label{def:2T-char-weak}
A \emph{weak polarized characterization interface} on the ledger-step skeleton
\(\{\Lambda_k\}_{k\in K_{\mathrm{acc}}}\) consists, for each ledger step \(k\)
\textup{(with \(K_{\mathrm{acc}}\) denoting the retained ledger-step index set)},
of the following data:
\begin{enumerate}[label=\textbf{(C\arabic*)},leftmargin=*,itemsep=0.5em]
\item \textbf{Odd generator space.} A real vector space
  \[
  E^-_{\Lambda_k}:=\{F\in\Fsloc:\Gamma_{\Lambda_k}F=-F\}
  \]
  of physically odd sufficiently local generators.
\item \textbf{Descendant ambiguity subspace.} A linear subspace
  \[
  \mathcal A^-_{\Lambda_k}
  \subset E^-_{\Lambda_k}.
  \]
\item \textbf{Admissible odd descendant class.} An affine odd class
  \[
  [\Xi^{(1)}_{\Lambda_k}]
  =
  J_{0,\Lambda_k}+\mathcal A^-_{\Lambda_k}
  \subset E^-_{\Lambda_k},
  \]
  where \(J_{0,\Lambda_k}\in E^-_{\Lambda_k}\) is any chosen representative.
\item \textbf{Orientation functional.} A continuous linear functional
  \[
  \chi_{\Lambda_k}:E^-_{\Lambda_k}\to\mathbb R
  \]
  such that
  \[
  \chi_{\Lambda_k}(A)=0
  \qquad
  \forall\,A\in\mathcal A^-_{\Lambda_k},
  \]
  and \(\chi_{\Lambda_k}\) is constant and nonzero on the affine class
  \([\Xi^{(1)}_{\Lambda_k}]\).
\item \textbf{Orientation sign.} A sign
  \[
  \varepsilon_{\Lambda_k}\in\{\pm 1\}
  \]
  such that
  \[
  \chi_{\Lambda_k}(J)=\varepsilon_{\Lambda_k}
  \qquad
  \forall\,J\in[\Xi^{(1)}_{\Lambda_k}].
  \]
\end{enumerate}
\end{definition}

\begin{definition}[Calibration supplement \(\mathbf{2T}_{\mathrm{char}}^{\mathrm{cal}}\)]
\label{def:2T-char-cal}
A \emph{calibration supplement} over
\(\mathbf{2T}_{\mathrm{char}}^{\mathrm{weak}}\) consists, for each ledger step \(k\),
of a symmetric positive definite bilinear form
  \[
  \mathsf B_{\Lambda_k}:E^-_{\Lambda_k}\times E^-_{\Lambda_k}\to\mathbb R
  \]
  whose induced norm Hilbertizes \(E^-_{\Lambda_k}\), and for which
  \(\mathcal A^-_{\Lambda_k}\) is closed in the induced Hilbert topology. In particular, the affine
  class \([\Xi^{(1)}_{\Lambda_k}]\) is a closed affine subspace of the associated
  Hilbert completion.
\end{definition}

\begin{definition}[Characterization-activity supplement \(\mathbf{2T}_{\mathrm{char}}^{\mathrm{act}}\)]
\label{def:2T-char-act}
A \emph{characterization-activity supplement} over
\(\mathbf{2T}_{\mathrm{Peierls}}+\mathbf{2T}_{\mathrm{char}}^{\mathrm{weak}}\) consists,
for each ledger step \(k\), of a linear subspace
\[
\mathcal P^+_{\Lambda_k}
\subset
E^+_{\Lambda_k}
:=
\{A\in\Fsloc:\Gamma_{\Lambda_k}A=A\}
\]
of physically even sufficiently local probes such that for every
\(F\in E^-_{\Lambda_k}\),
\[
\{F,A\}_{\mathrm P,\Lambda_k}=0
\qquad
\forall\,A\in\mathcal P^+_{\Lambda_k}
\]
implies
\[
\chi_{\Lambda_k}(F)=0.
\]
\end{definition}

\begin{remark}[Operational reading of \(\mathbf{2T}_{\mathrm{char}}^{\mathrm{act}}\)]
\label{rem:2T-char-act-reading}
Definition~\ref{def:2T-char-act} says that the characterization functional
\(\chi_{\Lambda_k}\) is not merely an orientation on the odd ambiguity quotient:
it is detected by a sufficiently local even probe sector.
Equivalently, if \(F\in E^-_{\Lambda_k}\) satisfies
\(\chi_{\Lambda_k}(F)\neq 0\), then there exists
\(A\in\mathcal P^+_{\Lambda_k}\subset\Fsloc\) with
\[
\{F,A\}_{\mathrm P,\Lambda_k}\neq 0.
\]
This is a targeted activity statement for the distinguished odd direction, not a
global Poisson-separation axiom on the full reduced phase space.
\end{remark}

\begin{definition}[Characterization-algebra supplement \(\mathbf{2T}_{\mathrm{alg}}\)]
\label{def:2T-alg}
A \emph{characterization-algebra supplement} consists, for each
ledger step \(k\), of a distinguished subalgebra
\[
\mathcal O^{\mathrm{char}}_{\Lambda_k}\subset \Fsloc
\]
attached to the derived odd generator.
\end{definition}

\begin{definition}[Aggregate polarized characterization interface \(\mathbf{2T}_{\mathrm{char}}\)]
\label{def:2T-char}
Define the aggregate characterization package by
\[
\mathbf{2T}_{\mathrm{char}}
:=
\mathbf{2T}_{\mathrm{char}}^{\mathrm{weak}}
+ \mathbf{2T}_{\mathrm{char}}^{\mathrm{cal}}
+ \mathbf{2T}_{\mathrm{alg}}.
\]
\end{definition}

\begin{remark}[Scope of \(\mathbf{2T}_{\mathrm{char}}^{\mathrm{weak}}\)]
\label{rem:char-cosmo-supplement-separate}
Definition~\ref{def:2T-char-weak} contains only the weak characterization data
for the descendant class. Additional ambiguity-control or reduction statements may be
formulated separately, but they are not part of
\(\mathbf{2T}_{\mathrm{char}}^{\mathrm{weak}}\) itself.
\end{remark}

\begin{remark}[Role of the weak characterization interface]
\label{rem:char-cosmo-realization}
The weak interface \(\mathbf{2T}_{\mathrm{char}}^{\mathrm{weak}}\) specifies the
abstract descendant class
\[
[\Xi^{(1)}_{\Lambda_k}]
=
J_{0,\Lambda_k}+\mathcal A^-_{\Lambda_k}.
\]
It does not identify a preferred odd generator representative. The calibration
form \(\mathsf B_{\Lambda_k}\) may be any admissible Hilbertizing form on
\(E^-_{\Lambda_k}\) compatible with this weak interface. At the level of the
base theorem,
orientation is introduced canonically only after nonvanishing, via
Proposition~\ref{prop:char-activity-from-nonzero-class}.
The characterization algebra \(\mathcal O^{\mathrm{char}}_{\Lambda_k}\) is isolated
separately as the supplement \(\mathbf{2T}_{\mathrm{alg}}\).
Stronger identifications among the orientation functional, the characterization
algebra, and active-generator data belong to an optional stronger supplement, not to
the base interface.
\end{remark}

\begin{theorem}[Every weak characterization interface admits a calibration supplement]
\label{thm:char-weak-to-cal}
Assume the weak polarized characterization interface
\(\mathbf{2T}_{\mathrm{char}}^{\mathrm{weak}}\)
\textnormal{(Definition~\ref{def:2T-char-weak})}. Then the calibration supplement
\(\mathbf{2T}_{\mathrm{char}}^{\mathrm{cal}}\)
\textnormal{(Definition~\ref{def:2T-char-cal})} holds.
\end{theorem}

\begin{proof}
Fix a ledger step \(k\). Since \(\mathcal A^-_{\Lambda_k}\subset E^-_{\Lambda_k}\) is a
linear subspace, choose an algebraic complement \(C_{\Lambda_k}\) so that
\[
E^-_{\Lambda_k}=\mathcal A^-_{\Lambda_k}\oplus C_{\Lambda_k}.
\]
Choose Hamel bases of \(\mathcal A^-_{\Lambda_k}\) and \(C_{\Lambda_k}\), declare the union
of those bases orthonormal, and extend bilinearly. This gives a symmetric positive
definite bilinear form \(\mathsf B_{\Lambda_k}\) on \(E^-_{\Lambda_k}\).

By construction, the induced norm identifies \(E^-_{\Lambda_k}\) with a pre-Hilbert
space whose Hilbert completion contains \(\mathcal A^-_{\Lambda_k}\) as a closed subspace,
because \(\mathcal A^-_{\Lambda_k}\) is orthogonal to \(C_{\Lambda_k}\).
Therefore the affine class
\(
[\Xi^{(1)}_{\Lambda_k}]=J_{0,\Lambda_k}+\mathcal A^-_{\Lambda_k}
\)
is a closed affine subspace in the induced Hilbert topology. This is exactly the
statement of Definition~\ref{def:2T-char-cal}.
\end{proof}

\begin{definition}[Selected-line supplement package \(\mathbf{2T}_{\mathrm{sel}}\)]
\label{def:2T-sel}
Fix a ledger step \(k\), and let \(\mathcal V^1\) be the first-descendant ambiguity group of
Definition~\ref{def:desc-ambiguity}. The supplementary package
\[
\mathbf{2T}_{\mathrm{sel}}
\]
consists of:
\begin{enumerate}[label=\textnormal{(\roman*)},leftmargin=*,itemsep=0.3em]
\item a nonzero selected odd line
  \[
  L^{\mathrm{sel}}_{\Lambda_k}
  =
  \mathbb R\,\ell_{\Lambda_k}
  \subset \mathcal V^1;
  \]
\item a selected representative
  \[
  J^{\mathrm{sel}}_{\Lambda_k}\in[\Xi^{(1)}_{\Lambda_k}]\cap E^-_{\Lambda_k};
  \]
\item the selected-line compatibility implication
  \begin{equation}
  q_{\Lambda_k}(J^{\mathrm{sel}}_{\Lambda_k})=0
  \quad\Longrightarrow\quad
  [J^{\mathrm{sel}}_{\Lambda_k}]_{\mathcal V^1}=0.
  \label{eq:selected-line-compatibility}
  \end{equation}
\end{enumerate}
\end{definition}

\begin{remark}[Same-tick versus delayed role of \(\mathbf{2T}_{\mathrm{sel}}\)]
\label{rem:Main-sel-two-lanes}
The package \(\mathbf{2T}_{\mathrm{sel}}\) supplies the same-tick selected-line
obstruction contract used in the square-loop no-extra-collapse argument.
It governs the same-tick selected quotient.
The corresponding delayed passage to the successor step is carried by
the separate supplement \(\mathbf{2T}_{\mathrm{sel}}^{\mathrm{delay}}\).
\end{remark}

\begin{remark}[Why \(\mathbf{2T}_{\mathrm{sel}}\) is separate]
\label{rem:Main-sel-separate}
The packages \(\mathbf{2T}_{\mathrm{sel}}\) and
\(\mathbf{2T}_{\mathrm{sel}}^{\mathrm{delay}}\) govern different
selected-line steps and are therefore kept as separate summands of
\(\mathbf{2T}_{\mathrm{Main}}\). The former supplies the
same-tick obstruction input used in
Theorem~\ref{thm:Main-faithful-LS-desc}; the latter supplies the delayed
lift to the successor step together with its nondegeneracy
clause.
\end{remark}

\begin{definition}[Delayed selected-line supplement package \(\mathbf{2T}_{\mathrm{sel}}^{\mathrm{delay}}\)]
\label{def:2T-sel-delay}
Fix consecutive steps \(k\) and \(k+1\).
The delayed selected-line supplement package consists of:
\begin{enumerate}[label=\textnormal{(\roman*)},leftmargin=*,itemsep=0.3em]
\item a delayed selected line
  \[
  \widehat L^{\mathrm{sel,delay}}_k
  =
  \mathbb R\,\widehat\ell^{\mathrm{delay}}_k
  \subset
  \tau^{\mathrm{comp,delay}}_k;
  \]
\item a delayed quotient line
  \[
  \overline L^{\mathrm{Main}}_{\Lambda_{k+1}}
  =
  \mathbb R\,\overline\ell^{\mathrm{Main}}_{\Lambda_{k+1}}
  \subset
  \mathfrak D^-_{\Lambda_{k+1}};
  \]
\item a linear delayed split-to-descendant map
  \[
  \mathfrak L^{\mathrm{split}}_{k\to k+1}:
  \widehat L^{\mathrm{sel,delay}}_k
  \longrightarrow
  \mathfrak D^-_{\Lambda_{k+1}}
  \]
  satisfying
  \[
  \mathfrak L^{\mathrm{split}}_{k\to k+1}
  (\widehat\ell^{\mathrm{delay}}_k)
  =
  \overline\ell^{\mathrm{Main}}_{\Lambda_{k+1}};
  \]
\item a descendant representative
  \[
  J^{\mathrm{split}}_{\Lambda_{k+1}}
  \in
  [\Xi^{(1)}_{\Lambda_{k+1}}]\cap E^-_{\Lambda_{k+1}}
  \]
  with
  \[
  q_{\Lambda_{k+1}}(J^{\mathrm{split}}_{\Lambda_{k+1}})
  =
  \overline\ell^{\mathrm{Main}}_{\Lambda_{k+1}};
  \]
\item dominant activity / nondegeneracy clause on the delayed image:
  \[
  \overline\rho^{\mathrm{dom}}_{\Lambda_{k+1}}
  (\overline\ell^{\mathrm{Main}}_{\Lambda_{k+1}})
  =
  c_{\Lambda_{k+1}}
  \neq 0.
  \]
\end{enumerate}
\end{definition}

\begin{definition}[UV conformal common-sector interface $\mathbf{2T}_{\mathrm{conf}}$]
\label{def:2T-conf}
A \emph{UV conformal common-sector interface} on the same vacuum representation
\((\mathcal H,\pi,\Omega)\) as \(\mathbf{2T}_{\mathrm{EBC}}\) consists of:
\begin{enumerate}[label=\textbf{(CF\arabic*)},leftmargin=*,itemsep=0.4em]
\item a dense UV-local conformal core \(\mathcal D_{\mathrm{conf}}\subset\mathcal H\)
  and UV-local subalgebras
  \[
  \mathfrak A_{\mathrm{conf}}(O)\subset \mathfrak A(O)
  \]
  for double cones \(O\subset\mathbb R^{1,3}\), with
  \(\pi(\mathfrak A_{\mathrm{conf}}(O))\Omega\subset\mathcal D_{\mathrm{conf}}\);
\item a strongly continuous unitary representation
  \(U_{\mathrm{conf}}(g)\) of the connected conformal group
  \(\mathrm{Conf}_0(\mathbb R^{1,3})\) on \(\mathcal H\), fixing the vacuum,
  extending the Poincar\'e representation of
  \(\mathbf{2T}_{\mathrm{EBC}}\), and acting geometrically on the UV-local core:
  \[
  U_{\mathrm{conf}}(g)\,\pi(\mathfrak A_{\mathrm{conf}}(O))\,U_{\mathrm{conf}}(g)^{-1}
  \;=\;
  \pi(\mathfrak A_{\mathrm{conf}}(gO));
  \]
\item a common Stone core
  \(\mathcal D_{\mathrm{Stone}}\subset\mathcal D_{\mathrm{conf}}\) for the selfadjoint
  generators \(P_\mu,M_{\mu\nu},D,K_\mu\) of the translation, Lorentz,
  dilatation, and special-conformal one-parameter subgroups, satisfying on
  \(\mathcal D_{\mathrm{Stone}}\) the standard conformal commutation relations,
  in particular
  \[
  [D,P_\mu]=iP_\mu,\qquad
  [D,K_\mu]=-\,iK_\mu,\qquad
  [K_\mu,P_\nu]=2i(\eta_{\mu\nu}D-M_{\mu\nu}),
  \]
  together with the usual Lorentz commutators of \(M_{\mu\nu}\) with
  \(P_\rho\), \(K_\rho\), and \(M_{\rho\sigma}\);
\item on the compactly smeared scalar UV class,
  the infinitesimal dilatation subgroup in \textnormal{(CF2)} agrees with the
  physical radial/RG derivation.
\end{enumerate}
\end{definition}

\begin{remark}[Why $\mathbf{2T}_{\mathrm{conf}}$ is isolated]
\label{rem:2T-conf-separate}
The package \(\mathbf{2T}_{\mathrm{conf}}\) isolates the UV
conformal/common-sector hypothesis. It requires that the conformal
generators \(P_\mu,M_{\mu\nu},D,K_\mu\) and the odd charges act on the same
vacuum representation and common core, which is the structural input needed for
the Lorentz finite-type and superconformal arguments.
\end{remark}

\begin{remark}[Physical meaning of $\mathbf{2T}_{\mathrm{conf}}$]
\label{rem:2T-conf-source-role}
The package \(\mathbf{2T}_{\mathrm{conf}}\) provides an abstract
same-representation conformal interface on the UV vacuum sector. Its content is
that the UV-local
net carries a conformal enhancement whose dilatation subgroup agrees with the
physical radial/RG action on the scalar UV class.
\end{remark}

\begin{definition}[Main package $\mathbf{2T}_{\mathrm{Main}}$]
\label{def:2T-Main}
Define the \emph{2T Main theorem package} by
\[
\mathbf{2T}_{\mathrm{Main}}
:=
\mathbf{2T}_{\mathrm{LS}}
+ \mathbf{2T}_{\mathrm{RG}}
+ \mathbf{2T}_{\mathrm{Peierls}}
+ \mathbf{2T}_{\mathrm{EBC}}
+ \mathbf{2T}_{\mathrm{char}}^{\mathrm{cal}}
+ \mathbf{2T}_{\mathrm{char}}^{\mathrm{act}}
+ \mathbf{2T}_{\mathrm{sel}}
+ \mathbf{2T}_{\mathrm{sel}}^{\mathrm{delay}}
+ \mathbf{2T}_{\mathrm{conf}}.
\]
\end{definition}

\subsubsection{Spine Interface}

\begin{definition}[Local two--time BRST interface $\mathbf{2T}_{\mathrm{BRST}}$]
\label{def:2T-BRST}
A \emph{local two--time BRST interface} consists of:
\begin{enumerate}[label=\textbf{(B\arabic*)},leftmargin=*,itemsep=0.5em]
\item A net (or presheaf) of sufficiently local observable cochains
  $\Fsloc(O)$ on the UV boundary, equipped with two (cohomological) differentials
  $Q_t$ (physical time) and $Q_u$ (RG/radial) such that on each $O$:
  \[
  Q_t^2=0,\qquad Q_u^2=0,\qquad Q_t Q_u + Q_u Q_t = 0.
  \]
\item A class of \emph{SK--admissible} states $\rho$ such that expectation values
  of $Q_t$--exact operators vanish:
  \[
  \langle Q_t X\rangle_\rho =0.
  \]
\end{enumerate}
\end{definition}

\begin{definition}[Main spine package]
\label{def:2T-Main-spine}
Define the \emph{Main spine package} by
\[
\mathbf{2T}_{\mathrm{Main}}^{\mathrm{spine}}
:= \mathbf{2T}_{\mathrm{Main}} + \mathbf{2T}_{\mathrm{BRST}}.
\]
\end{definition}

\begin{remark}
The package $\mathbf{2T}_{\mathrm{BRST}}$ is obtained in Subprogram~1a rather
than taken as an independent standing assumption. Thus
$\mathbf{2T}_{\mathrm{Main}}^{\mathrm{spine}}$ abbreviates the hypothesis
$\mathbf{2T}_{\mathrm{Main}}$ together with the derived bicomplex
$(\Fsloc,Q_t,Q_u)$ on which $\mathbf{2T}_{\mathrm{RG}}$ is interpreted.
\end{remark}

\begin{remark}[Derived BRST bicomplex]
Definition~\ref{def:2T-BRST} specifies the local two--time BRST/bicomplex
structure used in the descent arguments. The package
\(\mathbf{2T}_{\mathrm{BRST}}\) is not taken as an independent standing
assumption: Subprogram~1a derives the bicomplex and the \(Q_t\) Ward identity
from the bulk SK construction.
\end{remark}

\subsubsection{Internal Interfaces}

\begin{definition}[AQFT interface $\mathbf{2T}_{\mathrm{AQFT}}$]
\label{def:2T-AQFT}
An \emph{AQFT interface} on the Minkowski UV boundary consists of:
\begin{enumerate}[label=\textbf{(Q\arabic*)},leftmargin=*,itemsep=0.5em]
\item A Haag--Kastler net $O\mapsto\mathfrak A(O)$ of local *-algebras
  (dense Borchers--Uhlmann core) on $\mathbb R^{1,3}$ with isotony and locality.
\item A Poincar\'e action $\alpha_{(a,\Lambda)}$ by *-automorphisms with strong
  continuity on the vacuum representation.
\item A vacuum state $\varpi$ on $\mathfrak A$ that is Poincar\'e invariant and
  yields a GNS triple $(\mathcal H,\pi,\Omega)$.
\item Spectrum condition: the translation generators $P^\mu$ satisfy
  $\mathrm{spec}(P)\subset\overline V_+$.
\item \textbf{(H4) OS/positivity input.} The chosen renormalized counterterm scheme yields
  reflection positivity / temperedness on the boundary theory so that $\varpi$ is
  positive and the GNS construction applies.
\end{enumerate}
\end{definition}

\begin{remark}
We use the standard Haag--Kastler/AQFT setup \cite{HaagKastler1964} and treat
the BU algebra as a dense core for the local algebras.
\end{remark}

\begin{definition}[HLS interface $\mathbf{2T}_{\mathrm{HLS}}$]
\label{def:2T-HLS}
An \emph{HLS interface} consists of a nontrivial local derivation $\delta$ on
$\mathfrak A_{\mathrm{loc}}$, odd with respect to the physical $\Gamma$--grading, such that:
\begin{enumerate}[label=\textbf{(H\arabic*)},leftmargin=*,itemsep=0.4em]
\item $\delta$ is a graded *-derivation (for the $\Gamma$--grading), local, and densely defined.
\item \textbf{Physical oddness.} $\delta$ is odd with respect to the physical $\Gamma$--grading:
  \[
  \Gamma\circ\delta\circ\Gamma^{-1}=-\delta.
  \]
\item \textbf{Covariance (translations + Lorentz finite type).}
  $\delta$ commutes with translations,
  \[
  \alpha_a\circ\delta\circ\alpha_a^{-1}=\delta\qquad(a\in\mathbb R^{1,3}),
  \]
  and its Lorentz orbit spans a finite-dimensional module: writing
  $\delta_\Lambda:=\alpha_\Lambda\circ\delta\circ\alpha_\Lambda^{-1}$ for
  $\Lambda\in SO^+(1,3)$, the real linear span
  $\mathscr D:=\mathrm{span}_{\mathbb R}\{\delta_\Lambda\}$ is finite-dimensional.
\item $\varpi\circ\delta=0$ on $\mathfrak A^{+}$ (vacuum invariance on the
  $\Gamma$--even subalgebra), and $\delta$ is implementable by closable operators
  on a common invariant core.
\end{enumerate}
\end{definition}

\begin{remark}[Derived AQFT and HLS input]
Definitions~\ref{def:2T-AQFT} and~\ref{def:2T-HLS} record the hypotheses needed
to apply the HLS classification step. They are not taken as independent parts of
\(\mathbf{2T}_{\mathrm{Main}}\): Subprogram~1b derives the Wightman/AQFT
representation of the UV boundary theory, and Subprogram~3c verifies the HLS
admissibility properties for the constructed odd derivation~\(\delta\).
\end{remark}

\subsection{Boundary interface data}
At the base level one assumes the weak interface data
\[
\bigl(E^-_{\Lambda_k},\mathcal A^-_{\Lambda_k},[\Xi^{(1)}_{\Lambda_k}],\chi_{\Lambda_k},\varepsilon_{\Lambda_k}\bigr)
\]
and in particular the transported odd descendant class
\(
[\Xi^{(1)}_{\Lambda_k}]
=
J_{0,\Lambda_k}+\mathcal A^-_{\Lambda_k}
\)
and the ambiguity subspace \(\mathcal A^-_{\Lambda_k}\).
The full calibration supplement \(\mathbf{2T}_{\mathrm{char}}^{\mathrm{cal}}\) is then
derived abstractly from the weak interface by
Theorem~\ref{thm:char-weak-to-cal}. The characterization algebra
\(\mathcal O^{\mathrm{char}}_{\Lambda_k}\) is likewise derived from the canonical
representative by Theorem~\ref{thm:char-cal-to-alg}.
A stronger interface may also supply a normalized dominant-response functional
\(
\nu_{\Lambda_k}^{-1}\rho^{\mathrm{dom}}_{\Lambda_k}
\)
to witness activity directly, but the base package does not require that datum.
A stronger layer may identify this data with a rank-one active generator
\(J^{\mathrm{act}}_{\Lambda_k}\), a response functional \(\rho_{\Lambda_k}\), and a one-line
stiffness \(\widetilde\alpha(\Lambda_k)\).
Throughout Subprograms~3b and~3c we work on the chosen Minkowski boundary chart on
which the Peierls and EBC interfaces are assumed.

% ============================================================
\section{Subprogram 1a: Two--time SK bicomplex and Ward identities}
\label{sec:subprogram1a}

This subprogram provides the \emph{local two--time descent machinery} used later
in Subprograms~3a and~3b.  Starting from the Lorentzian Schwinger--Keldysh (SK) path
integral on a two--time base $(t,u)$ \cite{Schwinger1961,Keldysh1964}, we record
two odd differentials $Q_t,Q_u$ implementing infinitesimal reparametrisations
along $t$ (physical time) and $u$ (radial/RG time), and their Ward identities on
SK--admissible states.

\subsection{Two odd differentials from reparametrisations in $t$ and $u$}
\label{subsec:two-BRST}

\paragraph{Definition of $Q_t,Q_u$.}
On local bulk SK fields (doubled on the two SK branches) introduce two
Grassmann--odd scalar ghosts $c_t,c_u$ and define $Q_t,Q_u$ as Lie derivatives
along the odd vector fields $c_t\partial_t$ and $c_u\partial_u$:
\[
Q_t(\Psi) := \mathcal L_{c_t\partial_t}\Psi,\qquad
Q_u(\Psi) := \mathcal L_{c_u\partial_u}\Psi,
\]
extended by the graded Leibniz rule to local functionals.  We assume the base
coordinates commute, $[\partial_t,\partial_u]=0$.

\begin{remark}[Translation BRST convention]
In this appendix, $Q_t,Q_u$ encode the BRST differentials for the \emph{abelian}
translation symmetries in $t$ and $u$. Accordingly, we take the ghosts $c_t,c_u$
to be constant odd parameters (so $\partial_t c_t=\partial_u c_u=0$).
\end{remark}

\begin{lemma}[Nilpotency and anticommutation]
\label{lem:bi-complex}
The operators $Q_t$ and $Q_u$ are Grassmann--odd derivations satisfying
\[
Q_t^2 = 0,\qquad Q_u^2 = 0,\qquad \{Q_t,Q_u\}:=Q_tQ_u+Q_uQ_t=0.
\]
\end{lemma}

\begin{proof}
By the translation convention above, $c_t,c_u$ are constant odd parameters. Hence
the odd vector fields $c_t\partial_t$ and $c_u\partial_u$ satisfy
$(c_t\partial_t)^2=(c_u\partial_u)^2=0$ and
$[c_t\partial_t,c_u\partial_u]=0$ because $[\partial_t,\partial_u]=0$.
Taking Lie derivatives along these commuting square-zero odd vector fields gives
$Q_t^2=Q_u^2=\{Q_t,Q_u\}=0$.
\end{proof}

\begin{definition}[Total two--time differential]
\label{def:total-diff}
Define the total differential
\[
\mathbb Q := Q_t+Q_u.
\]
Then $\mathbb Q^2=0$ by Lemma~\ref{lem:bi-complex}.
\end{definition}

\subsection{Ward identities and SK--admissible states}
\label{subsec:ward}

\begin{definition}[SK--admissible states]
\label{def:SK-admissible}
A state is \emph{SK--admissible} if the initial density matrix is vacuum or KMS
(or a finite-energy perturbation thereof) and the sources on the two SK branches
are equal (physical SK).
\end{definition}

\begin{proposition}[Ward identities]
\label{prop:ward-identities}
For SK--admissible states, insertions of $Q_t$--exact operators have vanishing
expectation value on the closed time path:
\[
\langle Q_t X\rangle = 0.
\]
The $Q_u$ Ward identity holds up to radial boundary/anomaly terms determined by
the chosen renormalized variational principle (GHY + counterterms):
\[
\langle Q_u X\rangle = \textnormal{(radial boundary/anomaly term)}.
\]
At the strict UV fixed point on flat boundary data with trivial Weyl factor,
this radial term vanishes; for controlled Weyl pulses it reduces to the boundary
Weyl-anomaly density, which is the input for the anomaly-floor monotone in
Subprogram~3b.
\end{proposition}

\begin{proof}
Write the (renormalized) SK expectation value in the physical SK sector as
\[
\langle X\rangle
\;:=\;
\frac{1}{Z}\int \mathcal D\Psi_+\mathcal D\Psi_-\,e^{iS_{\mathrm{SK,ren}}[\Psi_+,\Psi_-]}\,X,
\]
where $S_{\mathrm{SK,ren}}$ includes the renormalized bulk action, the initial
state insertion (vacuum/KMS), and the renormalized boundary terms determined by
the variational principle (GHY + counterterms).

\paragraph{$Q_t$ Ward identity.}
Let $\varepsilon$ be a constant Grassmann--odd parameter and perform the change
of variables $\Psi_\pm\mapsto \Psi_\pm+\varepsilon\,Q_t\Psi_\pm$ in the SK path
integral.  For SK--admissible states, the contour is closed and the sources are
equal on the two branches, so the renormalized SK measure and action are
invariant under $t$--reparametrisations generated by $Q_t$ (no boundary term is
picked up on the closed time path).  Therefore the variation of the integral
vanishes:
\[
0
=
\varepsilon\int \mathcal D\Psi_+\mathcal D\Psi_-\,e^{iS_{\mathrm{SK,ren}}}\,(Q_t X),
\]
which implies $\langle Q_t X\rangle=0$.

\paragraph{$Q_u$ Ward identity.}
Similarly, perform $\Psi_\pm\mapsto \Psi_\pm+\varepsilon\,Q_u\Psi_\pm$.
Unlike $t$, the radial/RG direction $u$ has a boundary at the UV cutoff (and
possibly at an IR regulator), so the variation of the renormalized action
produces a boundary/anomaly term determined by the renormalized variational
principle.  Denoting the resulting insertion by $\mathcal A_u(X)$, one obtains
the Ward identity
\[
\langle Q_u X\rangle=\langle \mathcal A_u(X)\rangle.
\]
At a strict UV fixed point on flat boundary data with trivial Weyl factor this
insertion vanishes, whereas for controlled Weyl pulses it reduces to the
boundary Weyl-anomaly density (the standard radial/Weyl Ward insertion in the
Hamilton--Jacobi renormalization scheme).
\end{proof}

\begin{remark}
The $Q_t$ Ward identity is the usual SK statement that BRST-exact insertions
decouple for equal sources; the $Q_u$ Ward identity is the radial/Weyl analogue,
whose nonzero anomaly term is precisely the source of monotonicity away from
fixed points.
\end{remark}

\begin{proposition}[Identification of $\mathcal A^{(1)}$ with the renormalized radial/Weyl anomaly insertion]
\label{prop:A1-from-SK}
In the physical SK sector ($J_+=J_-$) the renormalized SK functional obeys a radial Ward identity
\[
\langle Q_u X\rangle \;=\; \langle \mathcal A_u(X)\rangle,
\]
where $\mathcal A_u(X)$ is a local boundary/anomaly insertion determined by the
renormalized variational principle.  On the monotone representative $\mathcal O^{(0)}(u)$,
the induced insertion defines a local cochain $\mathcal A^{(1)}(u)\in(\Fsloc)_{\bar 1}$ such that
\[
\big[\,Q_u^{\mathrm{phys}}\mathcal O^{(0)}(u)\,\big] \;=\; [\mathcal A^{(1)}(u)]\in H^1(\Fsloc,Q_t),
\]
and the descent equation holds in the locality quotient:
\[
\big[\,Q_u^{\mathrm{phys}}\mathcal O^{(0)}(u)+Q_t\Xi^{(1)}(u)-\mathcal A^{(1)}(u)\,\big]_{\mathrm{loc}}=0.
\]
\end{proposition}

\begin{proof}
By Proposition~\ref{prop:ward-identities}, the $u$--reparametrisation BRST
operator $Q_u$ obeys a Ward identity in the physical SK sector:
\[
\langle Q_u X\rangle \;=\; \langle \mathcal A_u(X)\rangle,
\]
where $\mathcal A_u(X)$ is the boundary/anomaly insertion produced by the
renormalized variational principle.

Apply this to the monotone representative $X=\mathcal O^{(0)}(u)$.
Define $\mathcal A^{(1)}(u)$ to be the induced sufficiently local cochain on
the UV boundary, i.e.\ the restriction of $\mathcal A_u(\mathcal O^{(0)}(u))$ to
the locality class $\Fsloc$ guaranteed by $\mathbf{2T}_{\mathrm{Peierls}}$-\textbf{(P5)}.
Since $Q_t\mathcal O^{(0)}(u)=0$ and $\{Q_t,Q_u^{\mathrm{phys}}\}$ vanishes on
the monotone representative (Definition~\ref{def:A1-class}), the cochain
$Q_u^{\mathrm{phys}}\mathcal O^{(0)}(u)$ is $Q_t$--closed and defines a class in
$H^1(\Fsloc,Q_t)$.
The Ward insertion $\mathcal A^{(1)}(u)$ represents the same class, i.e.
\[
\big[\,Q_u^{\mathrm{phys}}\mathcal O^{(0)}(u)\,\big]=[\mathcal A^{(1)}(u)]\in H^1(\Fsloc,Q_t),
\]
with scheme/reparametrisation effects absorbed in the locality quotient
(Definition~\ref{def:locality-quotient}).

Since $\mathcal A^{(1)}(u)$ represents the same $Q_t$--cohomology class as
$Q_u^{\mathrm{phys}}\mathcal O^{(0)}(u)$, there exists a (sufficiently local)
raw first-descendant potential \(\Xi^{(1)}(u)\in\Fsloc\) such that
\[
\mathcal A^{(1)}(u)-Q_u^{\mathrm{phys}}\mathcal O^{(0)}(u)=Q_t\Xi^{(1)}(u).
\]
Interpreting the equality modulo the standard locality ambiguities encoded in
Definition~\ref{def:locality-quotient} yields the descent equation in the
locality quotient:
\[
\big[\,Q_u^{\mathrm{phys}}\mathcal O^{(0)}(u)+Q_t\Xi^{(1)}(u)-\mathcal A^{(1)}(u)\,\big]_{\mathrm{loc}}=0.
\]
\end{proof}

\begin{lemma}[Locality of the anomaly insertion and existence of a sufficiently local potential]
\label{lem:A1-and-J-local}
Assume $\mathbf{2T}_{\mathrm{Peierls}}$-\textbf{(P5)}.  Then the anomaly insertion
$\mathcal A^{(1)}(u)$ extracted from the renormalized radial Ward identity lies in
$(\Fsloc)_{\bar 1}$, and there exists a first-descendant potential
\(\Xi^{(1)}(u)\in(\Fsloc)_{\bar 0}\) such that the anomalous descent equation holds in the locality quotient:
\[
\big[\,Q_u^{\mathrm{phys}}\mathcal O^{(0)}(u)+Q_t\Xi^{(1)}(u)-\mathcal A^{(1)}(u)\,\big]_{\mathrm{loc}}=0.
\]
\end{lemma}

\begin{proof}
By $\mathbf{2T}_{\mathrm{Peierls}}$-\textbf{(P5)}, the sufficiently local class
$\Fsloc$ is closed under the locality operations used to define renormalized Ward
insertions on the UV boundary.  In particular, the radial Ward insertion
$\mathcal A_u(\mathcal O^{(0)}(u))$ extracted from the renormalized variational
principle is a gauge-invariant local boundary functional, hence determines a
sufficiently local cochain $\mathcal A^{(1)}(u)\in\Fsloc$.  Since it is the
insertion associated to the odd operator $Q_u$, it has cohomological parity
$|\mathcal A^{(1)}(u)|_{\mathrm{ch}}=1$, i.e.\ $\mathcal A^{(1)}(u)\in(\Fsloc)_{\bar 1}$.

By Proposition~\ref{prop:A1-from-SK} and Definition~\ref{def:A1-class},
$\mathcal A^{(1)}(u)$ represents the same $Q_t$--cohomology class as
$Q_u^{\mathrm{phys}}\mathcal O^{(0)}(u)$ in $H^1(\Fsloc,Q_t)$.  Therefore there
exists a cochain \(\Xi^{(1)}(u)\in\Fsloc\) such that
\[
\mathcal A^{(1)}(u)-Q_u^{\mathrm{phys}}\mathcal O^{(0)}(u)=Q_t\Xi^{(1)}(u)
\qquad\text{in }\Fsloc,
\]
and \(\Xi^{(1)}(u)\) necessarily has even cohomological parity
\(|\Xi^{(1)}(u)|_{\mathrm{ch}}=0\), i.e.\ \(\Xi^{(1)}(u)\in(\Fsloc)_{\bar 0}\).
Rewriting yields
\[
Q_u^{\mathrm{phys}}\mathcal O^{(0)}(u)+Q_t\Xi^{(1)}(u)-\mathcal A^{(1)}(u)=0.
\]
Interpreting this identity modulo the standard scheme/contact/counterterm
ambiguities encoded in the locality quotient (Definition~\ref{def:locality-quotient})
gives the stated descent equation in $[\,\Fsloc\,]_{\mathrm{loc}}$.
\end{proof}

\subsection{Restriction to the UV boundary and the BRST interface}
\label{subsec:restriction-boundary}

\begin{theorem}[Boundary BRST interface]
\label{thm:subprogram1a-BRST}
Assume $\mathbf{2T}_{\mathrm{Peierls}}$-\textbf{(P5)} so that the sufficiently
local observable class $\Fsloc(O)$ is available on the UV boundary.  Then on each
bounded region $O$ the differentials $Q_t,Q_u$ restrict to $\Fsloc(O)$ and satisfy
\[
Q_t^2=Q_u^2=\{Q_t,Q_u\}=0,
\]
and on SK--admissible states one has $\langle Q_t X\rangle=0$.  In other words,
the interface $\mathbf{2T}_{\mathrm{BRST}}$ (Definition~\ref{def:2T-BRST}) holds
on the locality class used in the later Hamiltonization step.
\end{theorem}

\begin{proof}
By $\mathbf{2T}_{\mathrm{Peierls}}$-\textbf{(P5)}, for each bounded region $O$
on the UV boundary we have a sufficiently local cochain algebra $\Fsloc(O)$.
The bulk operators $Q_t,Q_u$ are defined on local SK functionals by Lie
derivatives along $c_t\partial_t$ and $c_u\partial_u$ (Subprogram~1a).  Since
restriction/pullback to the UV boundary commutes with Lie derivatives along
$\partial_t$ and $\partial_u$, the operators $Q_t,Q_u$ restrict to odd
derivations on each $\Fsloc(O)$.
The relations $Q_t^2=Q_u^2=\{Q_t,Q_u\}=0$ are algebraic consequences of
Lemma~\ref{lem:bi-complex} and therefore continue to hold after restriction to
the boundary.
Finally, the $Q_t$ Ward identity $\langle Q_t X\rangle=0$ on SK--admissible
states is Proposition~\ref{prop:ward-identities}.
Thus the conditions of Definition~\ref{def:2T-BRST} are satisfied on the
locality class used later in the construction.
\end{proof}

\begin{remark}
This is the point where the BRST/bicomplex operators
\(Q_t,Q_u\) are introduced; later sections treat them as fixed structure on the
chosen locality class.
\end{remark}

% ============================================================
\section{Subprogram 1b: Local net, GNS reconstruction, and implementation template}
\label{sec:subprogram1b}

This subprogram fixes the Minkowski UV boundary setup in which the odd symmetry
constructed in Subprogram~3b will be represented as an operator and classified by
HLS.  We keep the classification step itself in Subprogram~3c (after the odd
derivation is constructed), but we set up the operator-algebraic/Wightman
infrastructure here.

\subsection{Local net and vacuum state}
\label{subsec:net-and-state}

\begin{definition}[Dense Borchers--Uhlmann core]
\label{def:BU-core}
Let $O\subset\mathbb R^{1,3}$ be a double cone. Define $\mathfrak A_{\mathrm{BU}}(O)$
to be the *-algebra generated by symbols $\Phi(F)$ with $F\in\Fsloc(O)$, modulo
locality relations for spacelike separated supports and the *-structure
$\Phi(F)^*=\Phi(\overline F)$.  Write $\mathfrak A_{\mathrm{loc,BU}}:=\bigcup_O
\mathfrak A_{\mathrm{BU}}(O)$.
\end{definition}

\begin{remark}
The Borchers--Uhlmann net $O\mapsto\mathfrak A_{\mathrm{BU}}(O)$ is the dense
algebraic core used for GNS reconstruction.  A physically relevant completion
$\mathfrak A(O)$ will be defined in the vacuum representation obtained below.
\end{remark}

\begin{definition}[Vacuum state from the SK functional]
\label{def:omega-SK}
Assume the SK functional determines a consistent family of vacuum
correlators on the Minkowski UV boundary (equal sources on the two SK branches),
and that on this equal-source locus the renormalized SK functional is invariant
under swapping the two SK branches.
Define a linear functional $\varpi$ on $\mathfrak A_{\mathrm{loc,BU}}$ by
declaring its values on products of generators to be the corresponding SK
vacuum correlators, and extending by linearity.
\end{definition}

\begin{theorem}[Wightman/GNS net on the Minkowski UV boundary]
\label{thm:subprogram1b-Wightman}
Assume $\mathbf{2T}_{\mathrm{Peierls}}$-\textbf{(P5)} so that the locality
presheaf $O\mapsto\Fsloc(O)$ is available on the Minkowski UV chart, and assume
the energy bounds and covariance interface $\mathbf{2T}_{\mathrm{EBC}}$
(Definition~\ref{def:2T-EBC}).  Let $(\mathcal H,\pi,\Omega)$ denote the vacuum
GNS triple supplied by $\mathbf{2T}_{\mathrm{EBC}}$-\textnormal{(0i)}.
Defining local algebras by completion in this representation,
\[
\mathfrak A(O)\ :=\ C^*\!\big(\pi(\mathfrak A_{\mathrm{BU}}(O))\big)
\quad(\text{or equivalently the von Neumann closure}),
\]
the assignment $O\mapsto\mathfrak A(O)$ is a Haag--Kastler net on $\mathbb R^{1,3}$
with vacuum vector $\Omega$ and positive-energy Poincar\'e implementation.
Equivalently, the AQFT interface $\mathbf{2T}_{\mathrm{AQFT}}$
(Definition~\ref{def:2T-AQFT}) holds.
\end{theorem}

\begin{proof}
We split the argument into the state/representation, the net, locality, and
covariance.

\paragraph{State and GNS representation.}
By Definition~\ref{def:2T-EBC}-(0i), $\varpi$ is a vacuum state on the algebraic
inductive limit $*$-algebra
$
\mathfrak A_{\mathrm{loc,BU}}:=\bigcup_O \mathfrak A_{\mathrm{BU}}(O)
$
and $(\mathcal H,\pi,\Omega)$ is the corresponding vacuum GNS triple. In
particular,
\[
\varpi(A)=\langle\Omega,\pi(A)\Omega\rangle
\qquad(A\in\mathfrak A_{\mathrm{loc,BU}}).
\]

\paragraph{Local net and isotony.}
For each double cone $O$, define the local algebra by completion in the vacuum
representation:
\[
\mathfrak A(O):=C^*\!\big(\pi(\mathfrak A_{\mathrm{BU}}(O))\big).
\]
If $O_1\subset O_2$, then by construction
$\mathfrak A_{\mathrm{BU}}(O_1)\subset\mathfrak A_{\mathrm{BU}}(O_2)$, hence
$\pi(\mathfrak A_{\mathrm{BU}}(O_1))\subset\pi(\mathfrak A_{\mathrm{BU}}(O_2))$,
and taking C$^*$-closures yields isotony:
$\mathfrak A(O_1)\subset\mathfrak A(O_2)$.

\paragraph{Locality.}
Let $O_1,O_2$ be spacelike separated. Locality holds on the BU core by the
defining Borchers--Uhlmann relations, i.e.
$[A_1,A_2]=0$ for all
$A_i\in \pi(\mathfrak A_{\mathrm{BU}}(O_i))$.
To pass to the C$^*$-closures, take arbitrary $A_i\in\mathfrak A(O_i)$ and
choose sequences $A_{i,n}\in\pi(\mathfrak A_{\mathrm{BU}}(O_i))$ with
$\|A_{i,n}-A_i\|\to 0$.
Then $[A_{1,n},A_{2,m}]=0$ for all $n,m$, and using the operator norm estimate
\[
\|[B,C]\|\le 2\|B\|\,\|C-C'\|+2\|C'\|\,\|B-B'\|
\]
with $B=A_1$, $B'=A_{1,n}$, $C=A_2$, $C'=A_{2,m}$, we obtain
$\|[A_1,A_2]\|\to 0$ as $n,m\to\infty$. Hence $[A_1,A_2]=0$.

\paragraph{Poincar\'e covariance and spectrum.}
By Definition~\ref{def:2T-EBC}-(0ii), there exists a strongly continuous
positive-energy unitary representation $U_g$ of the Poincar\'e group implementing
the induced action $\alpha_g$ on the locality presheaf $O\mapsto\Fsloc(O)$ and
fixing the vacuum vector, $U_g\Omega=\Omega$.
The induced action extends to the BU core by
$\alpha_g(\Phi(F)):=\Phi(\alpha_gF)$ and multiplicativity, and the
implementation identity
$
U_g\,\pi(A)\,U_g^{-1}=\pi(\alpha_g(A))
$
holds on $\pi(\mathfrak A_{\mathrm{BU}}(O))$ by construction. Since conjugation
by $U_g$ is isometric, this identity extends by continuity to the C$^*$-closures
and yields
$
U_g\,\mathfrak A(O)\,U_g^{-1}=\mathfrak A(gO)
$
for all double cones $O$.
Finally, the spectrum condition and strong continuity are part of
Definition~\ref{def:2T-EBC}-(0ii), completing the verification that
$O\mapsto\mathfrak A(O)$ satisfies the Haag--Kastler properties with vacuum and
positive-energy Poincar\'e implementation, i.e.\ the AQFT interface
$\mathbf{2T}_{\mathrm{AQFT}}$ (Definition~\ref{def:2T-AQFT}) holds.
\end{proof}

\begin{remark}
One convenient reference mechanism for positive-energy implementation of
translations in a vacuum representation is the Borchers--Arveson theorem; in
this appendix spine we treat the spectrum condition and strong continuity as
part of the energy bounds and covariance interface (Definition~\ref{def:2T-EBC}).
\end{remark}

\subsection{Dense core and inner implementation template}
\label{subsec:inner-implementation-template}

\begin{definition}[Finite polynomial algebra and GNS core]
\label{def:A-fin}
Let $\mathfrak A_{\mathrm{fin}}\subset\mathfrak A_{\mathrm{loc,BU}}$ denote the
*-subalgebra generated by finite polynomials in the local generators
$\Phi(F)$ (with $F$ compactly supported in a double cone).  Define the common
GNS core
\[
\mathcal D_0 := \pi(\mathfrak A_{\mathrm{fin}})\Omega.
\]
\end{definition}

\begin{proposition}[Inner implementation on the GNS core]
\label{prop:inner}
Assume Theorem~\ref{thm:subprogram1b-Wightman} and let
$\delta:\mathfrak A_{\mathrm{fin}}\to \mathfrak A_{\mathrm{fin}}$ be a graded
*-derivation which is local, Poincar\'e covariant, and vacuum-invariant in the
sense that $\varpi\circ\delta=0$.  Assume moreover that $\delta$ is implementable
by closable operators on a common invariant core (as will be verified for the
constructed $\delta$ in Subprogram~3c). Then there exists a densely defined
closable operator $Q$ on $\mathcal H$ with $\mathcal D_0\subset\mathcal D(Q)$ such
that, for all $A\in\mathfrak A_{\mathrm{fin}}$,
\[
\delta(A)=i[Q,\pi(A)]_{\mathrm s}\quad\text{on }\mathcal D_0,
\]
where $[X,Y]_{\mathrm s}:=XY-(-1)^{|X||Y|}YX$ is the graded commutator.
\end{proposition}

\begin{proof}
Define an operator $Q_0$ on the GNS core $\mathcal D_0=\pi(\mathfrak A_{\mathrm{fin}})\Omega$
by
\[
Q_0\,\pi(B)\Omega \;:=\; i\,\pi(\delta(B))\Omega
\qquad(B\in\mathfrak A_{\mathrm{fin}}).
\]

\paragraph{Well-definedness.}
Suppose $\pi(B)\Omega=0$.  For any $A\in\mathfrak A_{\mathrm{fin}}$ one has
\[
\langle \pi(A)\Omega,\pi(\delta(B))\Omega\rangle
\;=\;
\varpi\!\big(A^*\delta(B)\big).
\]
Using the graded Leibniz rule and vacuum invariance $\varpi\circ\delta=0$,
\[
0
=\varpi\!\big(\delta(A^*B)\big)
=\varpi\!\big(\delta(A^*)B\big)+(-1)^{|A|}\varpi\!\big(A^*\delta(B)\big),
\]
so
$
\varpi(A^*\delta(B))=-(-1)^{|A|}\varpi(\delta(A^*)B)
=-(-1)^{|A|}\varpi((\delta(A))^*B).
$
Since $\pi(B)\Omega=0$, we have $\varpi(C^*B)=0$ for all $C\in\mathfrak A_{\mathrm{fin}}$,
and in particular for $C=\delta(A)$; hence $\varpi(A^*\delta(B))=0$ for all $A$.
Therefore $\pi(\delta(B))\Omega$ is orthogonal to the dense set
$\pi(\mathfrak A_{\mathrm{fin}})\Omega$, so $\pi(\delta(B))\Omega=0$ and $Q_0$ is well-defined.

\paragraph{Graded commutator identity on the core.}
For $A,B\in\mathfrak A_{\mathrm{fin}}$, using the definition of $Q_0$ and the graded
Leibniz rule for $\delta$,
\[
\begin{aligned}
[Q_0,\pi(A)]_{\mathrm s}\,\pi(B)\Omega
&=Q_0\,\pi(AB)\Omega-(-1)^{|A|}\pi(A)\,Q_0\,\pi(B)\Omega \\
&=i\,\pi(\delta(AB))\Omega-(-1)^{|A|}i\,\pi(A\delta(B))\Omega \\
&=i\,\pi(\delta(A)B)\Omega
\;=\; i\,\pi(\delta(A))\,\pi(B)\Omega,
\end{aligned}
\]
so $\delta(A)=i[Q_0,\pi(A)]_{\mathrm s}$ on $\mathcal D_0$.

\paragraph{Closability.}
By the implementability hypothesis, $Q_0$ admits a closable extension on a
common invariant core.  Let $Q$ denote such a densely defined closable
extension with $\mathcal D_0\subset\mathcal D(Q)$.  Then the graded commutator
identity continues to hold on $\mathcal D_0$, completing the proof.
\end{proof}

% ============================================================
\section{Subprogram 2a: Results from the 2T Limited Scope Theorem}
\label{sec:limited-scope-import}

\begin{theorem}[Limited Scope import]
\label{thm:LS-import}
Assume $\mathbf{2T}_{\mathrm{LS}}$. Then the Limited Scope Theorem provides:
\begin{enumerate}[label=\textnormal{(LS\arabic*)},leftmargin=*,itemsep=0.4em]
\item A presented record-process groupoid $\ProcG$ and a parity cocycle
  $\omegaLS:\ProcG\to B\mathbb Z_2$ with
  $\omegaLS(\sqLoop(t,u,r))=1$ for the canonical square-loop.
\item A non-splitting $\mathbb Z_2$-cover
  $\widetilde{\ProcG}\to\ProcG$ (no global section / no state-label trivialization).
\item A nontrivial action on $\Diff=\mathrm{Cofib}(\diag)$ witnessed by
  $y\in\Diff$ with $\flip(y)\neq y$ (detector nontriviality).
\item Bicomplex identities on free abelian chains generated by the
  bisimplicial set $\mathcal C^{\mathrm{SK}}_{m,n}$, yielding
  $(Q_t^{\mathrm{ch}})^2=(Q_u^{\mathrm{ch}})^2=
  Q_t^{\mathrm{ch}}Q_u^{\mathrm{ch}}+Q_u^{\mathrm{ch}}Q_t^{\mathrm{ch}}=0$.
\end{enumerate}
\end{theorem}

\begin{remark}[Two bicomplexes are not identified]
The LS bicomplex $(Q_t^{\mathrm{ch}},Q_u^{\mathrm{ch}})$ acts on record-process chains
generated by the bisimplicial set $\mathcal C^{\mathrm{SK}}_{m,n}$ and is used only to
produce the parity cocycle $\omega_{\mathrm{LS}}$.
The SK/BRST bicomplex $(Q_t,Q_u)$ acts on sufficiently local boundary cochains $\Fsloc$
and is used to formulate Ward identities and descent.
These bicomplexes interact only through the tag-action/parity-transport step
(Definition~\ref{def:actFill-obs} and Lemma~\ref{lem:parity-transport}); they are not
identified.
\end{remark}

% ============================================================
\section{Subprogram 2b: Cosmology theory Peierls bracket and swap invariance}
\label{sec:cosmo-peierls-import}

\begin{theorem}[Cosmology theory Peierls bracket]
\label{thm:cosmo-derived-Peierls}
Assume $\mathbf{2T}_{\mathrm{Peierls}}$. Then the Peierls bracket
$\{-,-\}_{\mathrm P}$ is a well-defined graded Poisson bracket on the algebra
$\Fsloc$ of gauge-invariant sufficiently local
boundary functionals, satisfying antisymmetry, Leibniz, and Jacobi.  This is
the cosmology theory Peierls package (P1--P6), recorded here in abstract
interface form.
\end{theorem}

\begin{remark}
For general background on the Peierls bracket and its covariant phase space
formulation, see \cite{Peierls1952,Khavkine2014Peierls}.
\end{remark}

\begin{lemma}[Transpose/swap equivariance implies equivariance of $E^\pm$ and $\Delta$]
\label{lem:GammaStar-commutes-Greens}
Assume $\mathbf{2T}_{\mathrm{Peierls}}$-\textbf{(P2)} and
\textbf{(P6)}, and let $\Gamma_*$ be the involution supplied by
\textbf{(P6)}.
Then
\[
E^\pm\circ\Gamma_*=\Gamma_*\circ E^\pm,
\qquad
\Delta\circ\Gamma_*=\Gamma_*\circ\Delta,
\quad\text{where}\quad \Delta:=E^- - E^+.
\]
\end{lemma}

\begin{proof}
Fix a test source $f$ and consider
\[
\psi^\pm := E^\pm(\Gamma_* f),
\qquad
\phi^\pm := \Gamma_*(E^\pm f).
\]
By definition, $\psi^\pm$ solves $P_{\mathrm{gf,ren}}\psi^\pm=\Gamma_* f$ with
$B_\Lambda(\psi^\pm)=0$ and advanced/retarded support. Using the commutation
assumption from $\mathbf{2T}_{\mathrm{Peierls}}$-\textbf{(P6)},
\[
P_{\mathrm{gf,ren}}\phi^\pm
= P_{\mathrm{gf,ren}}\Gamma_*E^\pm f
= \Gamma_* P_{\mathrm{gf,ren}}E^\pm f
= \Gamma_* f,
\qquad
B_\Lambda(\phi^\pm)
= B_\Lambda\Gamma_*E^\pm f
= \Gamma_* B_\Lambda(E^\pm f)
= 0.
\]
Moreover, $\phi^\pm$ has the same advanced/retarded support property as $E^\pm
f$ because $\Gamma_*$ is pointwise and does not enlarge support. By uniqueness
of $E^\pm$, one has $\psi^\pm=\phi^\pm$, i.e.\ $E^\pm\Gamma_*=\Gamma_*E^\pm$.
Subtracting yields $\Delta\Gamma_*=\Gamma_*\Delta$.
\end{proof}

\begin{lemma}[Swap induces $\Gamma_*$ on first functional variations]
\label{lem:Gamma-star-dictionary}
Let $\Gamma$ act on functionals by pullback along the induced swap on boundary
data, and let $\Gamma_*$ denote the induced involution on test sources/linearized
boundary data (the tangent map of the swap). Then for every $F\in\Fsloc$ one has
\[
(\Gamma F)^{(1)}=\Gamma_*\big(F^{(1)}\big),
\]
and similarly for second variations. In the pointer-basis instantiation,
$\Gamma_*$ is the transpose/real structure on linearized boundary data.
\end{lemma}

\begin{proof}
This is the chain rule for pullback: writing $(\Gamma F)(J):=F(\Gamma_*J)$ for
boundary data $J$, differentiation gives $(\Gamma F)^{(1)}=\Gamma_* F^{(1)}$.
\end{proof}

\begin{lemma}[Swap-equivariance and Poisson automorphism]
\label{lem:Gamma-Poisson}
Assume $\mathbf{2T}_{\mathrm{Peierls}}$, let $\Gamma_*$ be the involution of
\textbf{(P6)}, and let $\Gamma$ act on functionals by pullback along the
induced swap on boundary data. Then
\[
\Gamma\{F,G\}_{\mathrm P} = \{\Gamma F,\Gamma G\}_{\mathrm P}
\quad\text{for all }F,G\in\Fsloc.
\]
\end{lemma}

\begin{proof}
By definition, $\{F,G\}_{\mathrm P}=\langle F^{(1)},\Delta\,G^{(1)}\rangle$.
By Lemma~\ref{lem:Gamma-star-dictionary}, pullback by $\Gamma$ sends first
functional derivatives by $\Gamma_*$, and by
Lemma~\ref{lem:GammaStar-commutes-Greens} one has $\Delta\Gamma_*=\Gamma_*\Delta$.
Using the $\Gamma_*$--invariance of the pairing from
$\mathbf{2T}_{\mathrm{Peierls}}$-\textbf{(P6)} gives the claim.
\end{proof}

% ============================================================
\section{Subprogram 2c: Parity transport and physical $\mathbb Z_2$ grading}
\label{sec:parity-bridge}

\begin{definition}[SK parity candidate]
\label{def:GammaSK}
Let $\Gamma_{\mathrm{SK}}$ be the involutive *-automorphism induced by swapping
SK branches (equivalently: sending difference/$a$-type generators to minus
themselves and fixing $r$-type ones).
\end{definition}

\begin{lemma}[Swap invariance of the SK vacuum state]
\label{lem:varpi-Gamma-invariant}
Assume the vacuum state $\varpi$ is induced from the physical SK functional with
equal sources in the sense of Definition~\ref{def:omega-SK}. Then $\varpi$ is
$\Gamma_{\mathrm{SK}}$-invariant:
\[
\varpi\circ\Gamma_{\mathrm{SK}}=\varpi.
\]
\end{lemma}

\begin{proof}
On products of generators $\Phi(F)$ the state $\varpi$ is, by construction,
given by SK vacuum correlators at equal sources. Swap-invariance of the
equal-source SK functional implies these correlators are unchanged when every
insertion is acted on by $\Gamma_{\mathrm{SK}}$. By linearity (and continuity on
the completion), $\varpi\circ\Gamma_{\mathrm{SK}}=\varpi$ on the net.
\end{proof}

\begin{lemma}[$\Gamma$ is unitarily implemented in the vacuum GNS representation]
\label{lem:Gamma-implemented}
Assume Lemma~\ref{lem:varpi-Gamma-invariant}. Let $(\mathcal H,\pi,\Omega)$ be
the GNS triple of the vacuum state $\varpi$.
Then there exists a unitary $U_\Gamma$ on $\mathcal H$ with
$U_\Gamma\Omega=\Omega$ and
\[
U_\Gamma\,\pi(A)\,U_\Gamma^{-1}=\pi(\Gamma A)
\qquad(A\in\mathfrak A_{\mathrm{loc,BU}}),
\]
and $U_\Gamma^2=\mathbf 1$.
\end{lemma}

\begin{proof}
Define $U_\Gamma$ on the cyclic subspace $\pi(\mathfrak A_{\mathrm{loc,BU}})\Omega$
by
\[
U_\Gamma\,\pi(A)\Omega := \pi(\Gamma A)\Omega.
\]
This is well-defined and isometric because $\varpi\circ\Gamma=\varpi$ implies
\[
\|U_\Gamma\pi(A)\Omega\|^2
=\langle\Omega,\pi((\Gamma A)^*(\Gamma A))\Omega\rangle
=\varpi(\Gamma(A^*A))=\varpi(A^*A)=\|\pi(A)\Omega\|^2.
\]
Hence $U_\Gamma$ extends uniquely to a unitary on $\mathcal H$. The
implementation identity follows by construction, and $U_\Gamma^2=\mathbf 1$
because $\Gamma^2=\mathrm{id}$.
\end{proof}

\begin{lemma}[$\Gamma_{\mathrm{SK}}$ commutes with $Q_t$ and respects the RG decomposition]
\label{lem:GammaSK-comm}
On the locality class used in Subprograms~3a and~3b, $\Gamma_{\mathrm{SK}}$ commutes with
the two odd differentials:
\[
\Gamma_{\mathrm{SK}}Q_t = Q_t\Gamma_{\mathrm{SK}},\qquad
\Gamma_{\mathrm{SK}}Q_u = Q_u\Gamma_{\mathrm{SK}}.
\]
Moreover, for the decomposition $Q_u=Q_u^{\mathrm{phys}}+Q_u^{\mathrm{scheme}}$
used in Definition~\ref{def:2T-RG}, we assume that $\Gamma_{\mathrm{SK}}$
commutes with $Q_u^{\mathrm{phys}}$ and $Q_u^{\mathrm{scheme}}$ on the monotone sector
${\Fsloc}_{\mathrm{mon}}$ (Definition~\ref{def:Fsloc-mon}):
\[
\Gamma_{\mathrm{SK}}Q_u^{\mathrm{phys}} = Q_u^{\mathrm{phys}}\Gamma_{\mathrm{SK}},
\qquad
\Gamma_{\mathrm{SK}}Q_u^{\mathrm{scheme}} = Q_u^{\mathrm{scheme}}\Gamma_{\mathrm{SK}}
\quad\text{on }{\Fsloc}_{\mathrm{mon}}.
\]
\end{lemma}

\begin{proof}
On SK fields and sources, $\Gamma_{\mathrm{SK}}$ exchanges the two SK branches
and fixes the base ghosts. The differentials $Q_t,Q_u$ are defined diagonally
on the two branches (Subprogram~1a) as Lie derivatives along
$c_t\partial_t$ and $c_u\partial_u$. Since these definitions are symmetric under
interchanging the branches, $\Gamma_{\mathrm{SK}}$ commutes with $Q_t$ and $Q_u$
on generators; the graded Leibniz rule then extends the commutation to all
local cochains in the chosen locality class.
The additional equivariance for the decomposition $Q_u=Q_u^{\mathrm{phys}}+Q_u^{\mathrm{scheme}}$
is assumed on the monotone sector ${\Fsloc}_{\mathrm{mon}}$.
\end{proof}

\begin{lemma}[Parity involution on observables]
\label{lem:Gamma-def}
We write $\Gamma:=\Gamma_{\mathrm{SK}}$. Let $F$ be any observable functional on
$\SKstates$. Define
\[
(\Gamma F)(x):=F(\swap x).
\]
Then $\Gamma^2=\mathrm{id}$ and for all $r\in\Sys$,
$(\Gamma F)(\diag r)=F(\diag r)$.
\end{lemma}

\begin{remark}
In the pointer basis, $\swap$ acts as transposition on states, so for
$F(\rho)=\mathrm{Tr}(\rho\,\mathcal O)$ one has
$(\Gamma F)(\rho)=F(\rho^{\mathsf T})=\mathrm{Tr}(\rho\,\mathcal O^{\mathsf T})$.
\end{remark}

\begin{lemma}[Detector witness implies $\Gamma\neq \mathrm{id}$ on observables]
\label{lem:detector-Gamma-nontrivial}
Assume $\mathbf{2T}_{\mathrm{det}}$ in the pointer-basis instantiation, so there
exists a swap-odd observable functional $F$ on $\SKstates$ that is constant on
diagonal record states and is not identically zero. Then
\[
\Gamma(F) = -F,
\]
and in particular $\Gamma\neq \mathrm{id}$ on observables.
\end{lemma}

\begin{proof}
Swap-oddness gives $F(\swap x)=-F(x)$ for some $x\in\SKstates$, hence by
definition $(\Gamma F)(x)=F(\swap x)=-F(x)\neq F(x)$.
\end{proof}

\begin{remark}
By the non-splitting obstruction from the Limited Scope import, this nontrivial
parity cannot be removed by a relabeling of record states; it is a physical
grading.
\end{remark}

\begin{lemma}[$\omega$-action on additive targets]
\label{lem:omega-action-on-algebra}
Let $G$ be any abelian group. Define
\[
\mathrm{Act}_\omega^G(p)(g):=(-1)^{\omegaLS(p)}g.
\]
Then $\mathrm{Act}_\omega^G$ is a functorial action of $\ProcG$ on $G$. If
$\Gamma$ acts on an algebra $\mathcal A$ of observables, the action
$\mathrm{Act}(p):=\Gamma^{\omegaLS(p)}$ is a functorial representation of
$\ProcG$ on $\mathcal A$.
\end{lemma}

\begin{definition}[Tag action on observables]
\label{def:actFill-obs}
Define an action of the two square tags on observables by
\[
\mathrm{act}_{\mathrm{fill}}^{\mathrm{obs}}(\mathrm{diag})=\mathrm{id},\qquad
\mathrm{act}_{\mathrm{fill}}^{\mathrm{obs}}(\mathrm{bridge})=\Gamma,
\]
where $\Gamma$ is the swap-induced involution of Lemma~\ref{lem:Gamma-def}.
\end{definition}

\begin{lemma}[$\Gamma$ preserves the locality quotient]
\label{lem:Gamma-preserves-quotient}
For each double cone $O$, if $F\simeq F'$ in $\Fsloc(O)$
(Definition~\ref{def:locality-quotient}), then $\Gamma(F)\simeq \Gamma(F')$.
Equivalently, $\Gamma$ induces an involution on $[\,\Fsloc(O)\,]_{\mathrm{loc}}$.
\end{lemma}

\begin{proof}
Since $\simeq$ is generated by the elementary moves
Definition~\ref{def:locality-quotient}\textnormal{(a)}--\textnormal{(d)}, it
suffices to check stability of each move under $\Gamma$.

\paragraph{Move \textnormal{(a)} ($Q_t$--exact shifts).}
For $X\in\Fsloc(O)$,
\[
\Gamma(F+Q_tX)=\Gamma(F)+\Gamma(Q_tX)=\Gamma(F)+Q_t(\Gamma X),
\]
using $\Gamma Q_t=Q_t\Gamma$ (Lemma~\ref{lem:GammaSK-comm}).

\paragraph{Move \textnormal{(b)} (scheme move on the monotone sector).}
Let $Y\in{\Fsloc}_{\mathrm{mon}}(O)$.  By Definition~\ref{def:Fsloc-mon},
$\Gamma Y\in{\Fsloc}_{\mathrm{mon}}(O)$, and by Lemma~\ref{lem:GammaSK-comm} (on
${\Fsloc}_{\mathrm{mon}}$),
\[
\Gamma(Q_u^{\mathrm{scheme}}Y)=Q_u^{\mathrm{scheme}}(\Gamma Y).
\]
Moreover, $Q_u^{\mathrm{scheme}}$ preserves each $\Fsloc(O)$ by
Definition~\ref{def:2T-RG}-\textbf{(M2)}, so the shift is well-typed on $\Fsloc(O)$.

\paragraph{Move \textnormal{(c)} (total divergences).}
Since $\Gamma$ acts on fields/sources and not on the spacetime coordinates,
\[
\Gamma(F+\partial_\mu K^\mu)=\Gamma(F)+\partial_\mu(\Gamma K^\mu).
\]

\paragraph{Move \textnormal{(d)} (counterterm coboundaries).}
Counterterm coboundaries transform covariantly under the SK swap, so
$\Gamma(\Delta_{\mathrm{ct}})$ is again a counterterm coboundary.

Thus each generating move is preserved by $\Gamma$, hence $\Gamma$ preserves
$\simeq$ and induces an involution on $[\,\Fsloc(O)\,]_{\mathrm{loc}}$.
\end{proof}

\begin{lemma}[$\Gamma$--evenness of the anomaly insertion (modulo contact/counterterm ambiguities)]
\label{lem:Gamma-even-A1}
Assume the equal-source SK functional is invariant under SK branch swap, and let
$\mathcal A^{(1)}(u)$ be the anomaly insertion extracted from the renormalized radial Ward identity
in the physical SK sector (Proposition~\ref{prop:A1-from-SK}).
Assume moreover that any swap-odd remainder in the renormalized Ward insertion is a
contact/counterterm ambiguity (i.e.\ representable by a total divergence or a counterterm
coboundary). Then
\[
\big[\,\Gamma(\mathcal A^{(1)}(u))-\mathcal A^{(1)}(u)\,\big]_{\mathrm{loc}}=0,
\]
where the locality quotient is generated by $Q_t$--exact shifts
(Definition~\ref{def:locality-quotient}\textnormal{(a)}), total divergences
\textnormal{(c)}, and counterterm coboundaries \textnormal{(d)} (no use of the scheme move
\textnormal{(b)} is required here).
\end{lemma}

\begin{proof}
On the physical SK locus (equal sources), swapping the SK branches leaves the
renormalized functional invariant by hypothesis.  The corresponding radial/Weyl
Ward insertion extracted from the renormalized variational principle is
therefore $\Gamma$--invariant up to the usual local contact and counterterm
ambiguities of renormalized Ward identities.

By the additional hypothesis of the lemma, the swap-odd remainder
$\Gamma(\mathcal A^{(1)}(u))-\mathcal A^{(1)}(u)$ can be represented as a total
divergence $\partial_\mu K^\mu$ plus a counterterm coboundary $\Delta_{\mathrm{ct}}$
(and, if present, a $Q_t$--exact shift).  Each of these representatives is
trivial in the locality quotient by
Definition~\ref{def:locality-quotient}\textnormal{(a)}, \textnormal{(c)}, and
\textnormal{(d)}.  Hence
\[
\big[\,\Gamma(\mathcal A^{(1)}(u))-\mathcal A^{(1)}(u)\,\big]_{\mathrm{loc}}=0.
\]
No use of the scheme move \textnormal{(b)} is required.
\end{proof}

\begin{remark}[Parity of the anomaly insertion]
\label{rem:anomaly-parity}
Applying $\Gamma$ to the descent equation sends the anomaly insertion
$\mathcal A^{(1)}$ to $\Gamma(\mathcal A^{(1)})$. By
Lemma~\ref{lem:Gamma-even-A1}, the anomaly insertion is $\Gamma$--even in the
locality quotient, so the $\mathrm{diag}$ and $\mathrm{bridge}$ fillers correspond
to the same anomaly insertion modulo the standard locality ambiguities.
\end{remark}

\begin{remark}
In the parity-transport step we work with the physical radial differential
$Q_u^{\mathrm{phys}}$; scheme terms $Q_u^{\mathrm{scheme}}$ are treated as
locality-quotient ambiguities (counterterm reparametrisations) and do not affect
the transported $\Gamma$--odd sector.
\end{remark}

\begin{lemma}[Bridge filler acts by $\Gamma$ on first-descendant potentials]
\label{lem:bridgeActsByGamma}
Assume Lemma~\ref{lem:GammaSK-comm}. Let $\mathcal O^{(0)}$ be $Q_t$-closed and
$\Gamma$--even, i.e.\ $\Gamma(\mathcal O^{(0)})=\mathcal O^{(0)}$, and let
$\mathcal J$ and $\mathcal A^{(1)}$ satisfy the (possibly anomalous) first
descent equation
\[
\big[\,Q_u^{\mathrm{phys}}\mathcal O^{(0)} + Q_t\mathcal J - \mathcal A^{(1)}\,\big]_{\mathrm{loc}}=0,
\qquad
\big[\,\Gamma(\mathcal A^{(1)})-\mathcal A^{(1)}\,\big]_{\mathrm{loc}}=0.
\]
Define two tag-associated representatives by
\[
\mathcal J_{\mathrm{diag}}
:=\mathrm{act}_{\mathrm{fill}}^{\mathrm{obs}}(\mathrm{diag})(\mathcal J)
=\mathcal J,
\qquad
\mathcal J_{\mathrm{bridge}}
:=\mathrm{act}_{\mathrm{fill}}^{\mathrm{obs}}(\mathrm{bridge})(\mathcal J)
=\Gamma(\mathcal J).
\]
Then $\mathcal J_{\mathrm{diag}}$ and $\mathcal J_{\mathrm{bridge}}$
are both valid first-descendant potentials of $\mathcal O^{(0)}$ (for the same
anomaly insertion), and switching the
record-level square filler tag from $\mathrm{diag}$ to $\mathrm{bridge}$ induces
\[
\mathcal J_{\mathrm{bridge}}=\Gamma(\mathcal J_{\mathrm{diag}}).
\]
\end{lemma}

\begin{proof}
Apply $\Gamma$ to the descent identity in the locality quotient. By
Lemma~\ref{lem:GammaSK-comm} and Lemma~\ref{lem:Gamma-preserves-quotient},
\[
\big[\,Q_u^{\mathrm{phys}}(\Gamma\mathcal O^{(0)}) + Q_t(\Gamma\mathcal J) - \Gamma(\mathcal A^{(1)})\,\big]_{\mathrm{loc}}=0.
\]
Since $\Gamma\mathcal O^{(0)}=\mathcal O^{(0)}$ by assumption and
$[\,\Gamma(\mathcal A^{(1)})-\mathcal A^{(1)}\,]_{\mathrm{loc}}=0$, we obtain
\[
\big[\,Q_u^{\mathrm{phys}}\mathcal O^{(0)} + Q_t(\Gamma\mathcal J) - \mathcal A^{(1)}\,\big]_{\mathrm{loc}}=0,
\]
so $\Gamma\mathcal J$ is also a valid first-descendant potential (for the same
anomaly insertion) in the locality quotient. By Definition~\ref{def:actFill-obs}
it is exactly the $\mathrm{bridge}$-tag representative.
\end{proof}

\begin{definition}[First-descendant ambiguity group]
\label{def:desc-ambiguity}
Let $\mathcal V^1$ be the additive group of first-descendant potentials in the
chosen locality class, modulo the standard $Q_t$-- and $Q_u$--exact shifts:
\[
\mathcal V^1 :=
\frac{\Fsloc}{Q_t(\Fsloc) + Q_u(\Fsloc)}.
\]
We write $[\mathcal J]\in\mathcal V^1$ for the class of a chosen first-descendant
potential $\mathcal J\in\Fsloc$.
\end{definition}

\begin{lemma}[$\Gamma$--odd representative choice]
\label{lem:Gamma-odd-rep}
Assume Lemma~\ref{lem:GammaSK-comm} and that the locality class is linear over a
ring in which $2$ is invertible (e.g.\ over $\mathbb R$). Let
$\mathcal J$ be a representative of a class $[\mathcal J]\in\mathcal V^1$
satisfying $[\Gamma(\mathcal J)]=-[\mathcal J]$ in $\mathcal V^1$ and
define
\[
\mathcal J_-:=\tfrac12\big(\mathcal J-\Gamma\mathcal J\big).
\]
Then $\Gamma(\mathcal J_-)=-\mathcal J_-$ and
$[\mathcal J_-]=[\mathcal J]$ in $\mathcal V^1$.
\end{lemma}

\begin{proof}
The parity identity is immediate from $\Gamma^2=\mathrm{id}$. The class
condition implies $\Gamma(\mathcal J)+\mathcal J$ lies in the ambiguity subgroup
$Q_t(\Fsloc) + Q_u(\Fsloc)$. Therefore
\[
\mathcal J-\mathcal J_-=\tfrac12\big(\Gamma(\mathcal J)+\mathcal J\big)
\]
 lies in the ambiguity subgroup, so $[\mathcal J_-]=[\mathcal J]$ in
 $\mathcal V^1$.
\end{proof}

\begin{lemma}[Odd representative choice preserves the descent identity]
\label{lem:J-odd-choice}
Assume the setup of Lemma~\ref{lem:bridgeActsByGamma} and that the locality class
is linear over a ring in which $2$ is invertible (e.g.\ over $\mathbb R$).
Define
\[
\mathcal J_-:=\tfrac12\big(\mathcal J-\Gamma\mathcal J\big).
\]
Then $\Gamma(\mathcal J_-)=-\mathcal J_-$ and $\mathcal J_-$ is a valid
first-descendant potential for the same anomaly insertion, i.e.
\[
\big[\,Q_u^{\mathrm{phys}}\mathcal O^{(0)} + Q_t\mathcal J_- - \mathcal A^{(1)}\,\big]_{\mathrm{loc}}=0.
\]
Moreover, $\mathcal J_-\in\Fsloc$ and $\mathcal J_-$ is EB3-admissible whenever
$\mathcal J$ is.
\end{lemma}

\begin{proof}
By Lemma~\ref{lem:bridgeActsByGamma}, both $\mathcal J$ and $\Gamma\mathcal J$
are valid first-descendant potentials for the same anomaly insertion, so
\[
\big[\,Q_u^{\mathrm{phys}}\mathcal O^{(0)}+Q_t\mathcal J-\mathcal A^{(1)}\,\big]_{\mathrm{loc}}=0,
\qquad
\big[\,Q_u^{\mathrm{phys}}\mathcal O^{(0)}+Q_t(\Gamma\mathcal J)-\mathcal A^{(1)}\,\big]_{\mathrm{loc}}=0.
\]
Subtracting gives
\[
\big[\,Q_t(\mathcal J-\Gamma\mathcal J)\,\big]_{\mathrm{loc}}=0.
\]
Define $\mathcal J_-:=\tfrac12(\mathcal J-\Gamma\mathcal J)$, so $\Gamma(\mathcal J_-)=-\mathcal J_-$.
Now compute the descent residual for $\mathcal J_-$ in the locality quotient:
\[
\begin{aligned}
\big[\,Q_u^{\mathrm{phys}}\mathcal O^{(0)}+Q_t\mathcal J_- - \mathcal A^{(1)}\,\big]_{\mathrm{loc}}
&=\Big[\,Q_u^{\mathrm{phys}}\mathcal O^{(0)}-\mathcal A^{(1)}+\tfrac12\big(Q_t\mathcal J-Q_t(\Gamma\mathcal J)\big)\,\Big]_{\mathrm{loc}} \\
&=\Big[\,-Q_t\mathcal J+\tfrac12\big(Q_t\mathcal J-Q_t(\Gamma\mathcal J)\big)\,\Big]_{\mathrm{loc}}
\quad\text{(using }\big[\,Q_u^{\mathrm{phys}}\mathcal O^{(0)}+Q_t\mathcal J-\mathcal A^{(1)}\,\big]_{\mathrm{loc}}=0\text{)}\\
&=\Big[\,-Q_t\!\big(\tfrac12(\mathcal J+\Gamma\mathcal J)\big)\,\Big]_{\mathrm{loc}}
\;=\;0,
\end{aligned}
\]
because $Q_t$--exact shifts are trivial by Definition~\ref{def:locality-quotient}\textnormal{(a)}.
Finally, $\Gamma$ preserves $\Fsloc$ and the defining seminorm controls on each $O$,
so $\mathcal J_-\in\Fsloc$ and inherits EB3-admissibility from $\mathcal J$.
\end{proof}

\begin{remark}
Lemma~\ref{lem:J-odd-choice} shows that the $\Gamma$--odd projection
$\mathcal J_-=\tfrac12(\mathcal J-\Gamma\mathcal J)$ preserves the descent data
in the locality quotient.
The involution $\Gamma$ preserves the locality class $\Fsloc(O)$ and its defining
seminorm controls; hence $\mathcal J_-=\tfrac12(\mathcal J-\Gamma\mathcal J)$ remains
EB3-admissible whenever $\mathcal J$ is.
\par
Because the descent equation is imposed only in the locality quotient, it fixes the
cohomology class $[\,Q_u^{\mathrm{phys}}\mathcal O^{(0)}\,]\in H^1(\Fsloc,Q_t)$ but does not
uniquely fix the potential. The LS parity transport identifies the physically relevant
\(\Gamma\)--odd descendant class, and Theorem~\ref{thm:finite-cutoff-Jcan} then selects its
canonical representative by calibration minimization.
\end{remark}

\begin{lemma}[Filler--descent dictionary]
\label{lem:filler-descent-dictionary}
Assume the setup of Lemma~\ref{lem:bridgeActsByGamma} so that
$\mathcal J_{\mathrm{diag}}$ and $\mathcal J_{\mathrm{bridge}}$ are defined from
a first-descendant potential $\mathcal J$.
Let $[\mathcal J]\in\mathcal V^1$ denote the class of a first-descendant
modulo the standard $Q_t$/$Q_u$ ambiguities. By the Limited Scope import,
the canonical square-loop $\sqLoop(t,u,r)$ satisfies $\omegaLS(\sqLoop(t,u,r))=1$.
In the sign transport induced by $\omegaLS$ on the additive group $\mathcal V^1$
(Lemma~\ref{lem:omega-action-on-algebra}), changing the filler tag from
$\mathrm{diag}$ to $\mathrm{bridge}$ multiplies the descendant class by $-1$:
\[
[\,\mathcal J_{\mathrm{bridge}}\,]
=-\, [\,\mathcal J_{\mathrm{diag}}\,]
\quad\text{in }\mathcal V^1.
\]
Equivalently, the descendant sector is a torsor for the $\omega$-twisted local system.
\end{lemma}

\begin{proof}
Apply Lemma~\ref{lem:omega-action-on-algebra} to the abelian group
$G=\mathcal V^1$. This defines a functorial action of the record-process groupoid
$\ProcG$ on $\mathcal V^1$ by
\[
\mathrm{Act}_\omega^{\mathcal V^1}(p)([\mathcal J])=(-1)^{\omegaLS(p)}[\mathcal J].
\]
By the Limited Scope import (Theorem~\ref{thm:LS-import}), the canonical square-loop
$p:=\sqLoop(t,u,r)$ satisfies $\omegaLS(p)=1$, hence
\[
\mathrm{Act}_\omega^{\mathcal V^1}(p)([\mathcal J])=-[\mathcal J].
\]
Changing the record-level square filler tag from $\mathrm{diag}$ to $\mathrm{bridge}$
corresponds to transporting the first-descendant class along this canonical square-loop,
so $[\mathcal J_{\mathrm{bridge}}]=\mathrm{Act}_\omega^{\mathcal V^1}(p)([\mathcal J_{\mathrm{diag}}])$.
Therefore $[\mathcal J_{\mathrm{bridge}}]=- [\mathcal J_{\mathrm{diag}}]$ in $\mathcal V^1$.
\end{proof}

\begin{lemma}[Parity transport: LS square-loop acts as $\Gamma$]
\label{lem:parity-transport}
Assume $\mathbf{2T}_{\mathrm{LS}}$ and the tag action
Definition~\ref{def:actFill-obs}.  On the first-descendant ambiguity group
$\mathcal V^1$ (Definition~\ref{def:desc-ambiguity}), changing the square filler
from $\mathrm{diag}$ to $\mathrm{bridge}$ acts by $\Gamma$ on representatives and
by multiplication by $-1$ on classes:
\[
[\Gamma(\mathcal J)]=-[\mathcal J]\quad\text{in }\mathcal V^1.
\]
In particular, if $[\mathcal J]\neq 0$ then it lies in the odd sector of
the $\Gamma$--grading.
\end{lemma}

\begin{proof}
By Lemma~\ref{lem:bridgeActsByGamma},
$\mathcal J_{\mathrm{bridge}}=\Gamma(\mathcal J_{\mathrm{diag}})$,
while Lemma~\ref{lem:filler-descent-dictionary} gives
$[\mathcal J_{\mathrm{bridge}}]=-[\mathcal J_{\mathrm{diag}}]$ in
$\mathcal V^1$. Thus $[\Gamma(\mathcal J)]=-[\mathcal J]$.
\end{proof}

% ============================================================
\section{Subprogram 3: Selected-line nonvanishing, Hamiltonized odd derivation, AQFT and HLS}
\label{sec:subprogram3}

\subsection{Subprogram 3a: Selected-line nonvanishing, canonicalization, and characteristic activity}
\label{sec:subprogram3a}

In this subsection we work directly with the two-time bicomplex
\((\Fsloc,Q_t,Q_u)\) supplied by Subprogram~1a and allow an explicit anomaly
insertion as in \eqref{eq:anom-descent}. This keeps the mixed-anomaly
bookkeeping in the full two-ghost complex.

\begin{theorem}[Finite-cutoff canonical odd representative]
\label{thm:finite-cutoff-Jcan}
Assume the calibration characterization package
\(\mathbf{2T}_{\mathrm{char}}^{\mathrm{cal}}\), namely
Definition~\ref{def:2T-char-weak} together with
Definition~\ref{def:2T-char-cal}. For each ledger step \(k\), let
\(
[\Xi^{(1)}_{\Lambda_k}]=J_{0,\Lambda_k}+\mathcal A^-_{\Lambda_k}\subset E^-_{\Lambda_k}
\)
be the admissible odd descendant class. Then there exists a unique
representative
\[
\mathcal J^{\mathrm{can}}_{\Lambda_k}\in [\Xi^{(1)}_{\Lambda_k}]
\]
such that
\begin{equation}
\mathcal J^{\mathrm{can}}_{\Lambda_k}
=
\arg\min_{J\in [\Xi^{(1)}_{\Lambda_k}]}
\mathsf B_{\Lambda_k}(J,J).
\label{eq:Jcan-min}
\end{equation}
\end{theorem}

\begin{proof}
By Definition~\ref{def:2T-char-cal}, the form \(\mathsf B_{\Lambda_k}\) Hilbertizes
\(E^-_{\Lambda_k}\); write \(H_{\Lambda_k}\) for the resulting Hilbert space with
inner product still denoted \(\mathsf B_{\Lambda_k}\).
Since \(\mathcal A^-_{\Lambda_k}\subset E^-_{\Lambda_k}\) is closed, the affine set
\(
[\Xi^{(1)}_{\Lambda_k}]=J_{0,\Lambda_k}+\mathcal A^-_{\Lambda_k}
\)
is a closed affine subspace of \(H_{\Lambda_k}\). The Hilbert projection theorem
therefore gives a unique element of this affine subspace of minimal
\(\mathsf B_{\Lambda_k}\)-norm, which is exactly
\(\mathcal J^{\mathrm{can}}_{\Lambda_k}\).
\end{proof}

\begin{proposition}[Orthogonality characterization of \(\mathcal J^{\mathrm{can}}_{\Lambda_k}\)]
\label{prop:Jcan-orth}
Under the hypotheses of Theorem~\ref{thm:finite-cutoff-Jcan}, the unique minimizer
\(\mathcal J^{\mathrm{can}}_{\Lambda_k}\) is characterized equivalently by
\begin{equation}
\mathsf B_{\Lambda_k}(\mathcal J^{\mathrm{can}}_{\Lambda_k},A)=0
\qquad
\forall\,A\in\mathcal A^-_{\Lambda_k}.
\label{eq:Jcan-orth}
\end{equation}
\end{proposition}

\begin{proof}
Write \(\mathcal J^{\mathrm{can}}_{\Lambda_k}=J_{0,\Lambda_k}+A_{\ast}\) with
\(A_{\ast}\in\mathcal A^-_{\Lambda_k}\).
For any \(A\in\mathcal A^-_{\Lambda_k}\) and \(s\in\mathbb R\), the point
\(\mathcal J^{\mathrm{can}}_{\Lambda_k}+sA\) lies in the same affine class.
Minimality of \(\mathcal J^{\mathrm{can}}_{\Lambda_k}\) implies
\[
0=\frac{d}{ds}\Big\rvert_{s=0}
\mathsf B_{\Lambda_k}(\mathcal J^{\mathrm{can}}_{\Lambda_k}+sA,\,
\mathcal J^{\mathrm{can}}_{\Lambda_k}+sA)
=2\,\mathsf B_{\Lambda_k}(\mathcal J^{\mathrm{can}}_{\Lambda_k},A),
\]
which is \eqref{eq:Jcan-orth}. Conversely, if \(J\in[\Xi^{(1)}_{\Lambda_k}]\)
 satisfies \eqref{eq:Jcan-orth}, then for every \(A\in\mathcal A^-_{\Lambda_k}\),
\[
\mathsf B_{\Lambda_k}(J+A,J+A)
=
\mathsf B_{\Lambda_k}(J,J)+\mathsf B_{\Lambda_k}(A,A)
\ge
\mathsf B_{\Lambda_k}(J,J),
\]
so \(J\) is the unique minimizer on the affine class.
\end{proof}

\begin{remark}[Finite-dimensional projection formula]
\label{rem:Jcan-finite-dimensional-projector}
In any finite-dimensional realization of the affine odd descendant class
\[
[\Xi^{(1)}_{\Lambda_k}]
=
J_{0,\Lambda_k}+\operatorname{im}(U_{\Lambda_k}),
\]
where \(U_{\Lambda_k}\) is any full-column-rank basis matrix for
\(\mathcal A^-_{\Lambda_k}\), let \(B_{\Lambda_k}\) be the matrix of the
calibration form \(\mathsf B_{\Lambda_k}\).
Then the \(\mathsf B_{\Lambda_k}\)-orthogonal projector onto the ambiguity
subspace is
\[
P_{\mathcal A^-_{\Lambda_k}}^{\mathsf B_{\Lambda_k}}
:=
U_{\Lambda_k}
\bigl(U_{\Lambda_k}^{\top}B_{\Lambda_k}U_{\Lambda_k}\bigr)^{-1}
U_{\Lambda_k}^{\top}B_{\Lambda_k},
\]
and the canonical representative is
\[
\mathcal J^{\mathrm{can}}_{\Lambda_k}
=
\bigl(I-P_{\mathcal A^-_{\Lambda_k}}^{\mathsf B_{\Lambda_k}}\bigr)
J_{0,\Lambda_k}.
\]
Equivalently, if
\[
a_{\ast}
:=
-\bigl(U_{\Lambda_k}^{\top}B_{\Lambda_k}U_{\Lambda_k}\bigr)^{-1}
U_{\Lambda_k}^{\top}B_{\Lambda_k}J_{0,\Lambda_k},
\]
then
\[
\mathcal J^{\mathrm{can}}_{\Lambda_k}
=
J_{0,\Lambda_k}+U_{\Lambda_k}a_{\ast}.
\]
This formula is independent of the chosen basis \(U_{\Lambda_k}\) of
\(\mathcal A^-_{\Lambda_k}\): replacing \(U_{\Lambda_k}\) by
\(U_{\Lambda_k}T\) with \(T\in \mathrm{GL}(r,\mathbb R)\) leaves the projector
\(P_{\mathcal A^-_{\Lambda_k}}^{\mathsf B_{\Lambda_k}}\) unchanged.
It is also independent of the chosen representative
\(J_{0,\Lambda_k}\in[\Xi^{(1)}_{\Lambda_k}]\): replacing \(J_{0,\Lambda_k}\)
by \(J_{0,\Lambda_k}+U_{\Lambda_k}h\) leaves
\(
\bigl(I-P_{\mathcal A^-_{\Lambda_k}}^{\mathsf B_{\Lambda_k}}\bigr)
J_{0,\Lambda_k}
\)
unchanged.
\end{remark}

\begin{remark}[Proof strategy for same-tick obstruction nonvanishing]
\label{rem:Main-nonvanishing-strategy}
The nonvanishing step is purely quotient-theoretic for the same-tick obstruction lane once one has a
faithful LS-to-descendant intertwiner.
One first isolates a selected odd descendant line inside the raw first-descendant
ambiguity group \(\mathcal V^1\) and uses Lemma~\ref{lem:parity-transport} to
see that the LS square-loop acts there by sign.
One then maps that line into the characterization quotient
\(E^-_{\Lambda_k}/\mathcal A^-_{\Lambda_k}\).
If this realization map is faithful on the selected line, then the image of a
nonzero generator is itself nonzero; equivariance therefore upgrades the sign
transport on the raw line to a nontrivial sign transport on the selected quotient
point.
At that stage the conclusion \(0\notin[\Xi^{(1)}_{\Lambda_k}]\) is immediate,
because \(0\) is the only fixed point of the sign involution on a real vector
space. This is the abstract same-tick obstruction route; the physically correct
delayed route is the lift recorded below.
\end{remark}

\begin{remark}[Proof skeleton for the same-tick LS-to-descendant faithfulness step]
\label{rem:Main-faithfulness-skeleton}
Fix a ledger step \(k\).
The structured proof has four steps:
\begin{enumerate}[label=\textnormal{(\roman*)},leftmargin=*,itemsep=0.3em]
\item choose a selected odd descendant line
  \(L^{\mathrm{sel}}_{\Lambda_k}=\mathbb R\,\ell_{\Lambda_k}\subset\mathcal V^1\)
  with \(\ell_{\Lambda_k}\neq 0\);
\item use Lemma~\ref{lem:parity-transport} to identify the LS square-loop action on
  that line with multiplication by \(-1\);
\item realize the selected line on the characterization quotient by a linear map
  \(\mathfrak R_{\Lambda_k}:L^{\mathrm{sel}}_{\Lambda_k}\to
  E^-_{\Lambda_k}/\mathcal A^-_{\Lambda_k}\) whose value on
  \(\ell_{\Lambda_k}\) is the selected quotient point
  \(\overline{\Xi}^{(1)}_{\Lambda_k}\);
\item prove faithfulness on the selected line, i.e.\ \(\ker\mathfrak R_{\Lambda_k}=0\).
\end{enumerate}
Then \(\overline{\Xi}^{(1)}_{\Lambda_k}=\mathfrak R_{\Lambda_k}(\ell_{\Lambda_k})\neq 0\),
and equivariance gives
\[
\mathrm{Act}^{\mathrm{desc}}_{\sqLoop}(\overline{\Xi}^{(1)}_{\Lambda_k})
=
-\,\overline{\Xi}^{(1)}_{\Lambda_k}
\neq
\overline{\Xi}^{(1)}_{\Lambda_k}.
\]
The theorem below expands this same-tick obstruction skeleton into a full proof.
\end{remark}

\begin{proposition}[Selected-line compatibility suffices for the same-tick faithfulness hypotheses]
\label{prop:Main-selected-line-compatibility}
Assume \(\mathbf{2T}_{\mathrm{LS}}+\mathbf{2T}_{\mathrm{char}}^{\mathrm{weak}}\).
Fix a ledger step \(k\), and let
\[
q_{\Lambda_k}:E^-_{\Lambda_k}\longrightarrow E^-_{\Lambda_k}/\mathcal A^-_{\Lambda_k}
\]
be the quotient map.
Assume the selected-line supplement package
\(\mathbf{2T}_{\mathrm{sel}}\) of Definition~\ref{def:2T-sel}, and write
\[
\ell_{\Lambda_k}
:=
[J^{\mathrm{sel}}_{\Lambda_k}]_{\mathcal V^1}
\neq 0.
\]
Define
\[
L^{\mathrm{sel}}_{\Lambda_k}
:=
\mathbb R\,\ell_{\Lambda_k},
\qquad
\mathfrak R_{\Lambda_k}(a\,\ell_{\Lambda_k})
:=
a\,q_{\Lambda_k}(J^{\mathrm{sel}}_{\Lambda_k}).
\]
Then hypotheses \textnormal{(a)}--\textnormal{(e)} of
Theorem~\ref{thm:Main-faithful-LS-desc} hold.
\end{proposition}

\begin{proof}
Condition \textnormal{(a)} of Theorem~\ref{thm:Main-faithful-LS-desc} holds by construction:
\(
L^{\mathrm{sel}}_{\Lambda_k}=\mathbb R\,\ell_{\Lambda_k}
\)
is one-dimensional and nonzero because
\(
\ell_{\Lambda_k}\neq 0
\).

Condition \textnormal{(b)} holds because \(\mathfrak R_{\Lambda_k}\) is defined on a real line by
scalar multiplication of the fixed quotient element \(q_{\Lambda_k}(J^{\mathrm{sel}}_{\Lambda_k})\).

Condition \textnormal{(c)} holds by direct substitution:
\[
\mathfrak R_{\Lambda_k}(\ell_{\Lambda_k})
=
q_{\Lambda_k}(J^{\mathrm{sel}}_{\Lambda_k}).
\]

For condition \textnormal{(d)}, linearity on the one-dimensional real line gives
\[
\mathfrak R_{\Lambda_k}(-\ell_{\Lambda_k})
=
-\,\mathfrak R_{\Lambda_k}(\ell_{\Lambda_k}).
\]

It remains to prove faithfulness \textnormal{(e)}.
Let \(a\,\ell_{\Lambda_k}\in\ker\mathfrak R_{\Lambda_k}\).
Then
\[
a\,q_{\Lambda_k}(J^{\mathrm{sel}}_{\Lambda_k})=0.
\]
If \(a\neq 0\), then
\(
q_{\Lambda_k}(J^{\mathrm{sel}}_{\Lambda_k})=0
\),
so the selected-line compatibility implication
\eqref{eq:selected-line-compatibility} yields
\[
[J^{\mathrm{sel}}_{\Lambda_k}]_{\mathcal V^1}=0,
\]
contradicting the definition of \(\ell_{\Lambda_k}\neq 0\).
Hence \(a=0\), and therefore
\[
\ker\mathfrak R_{\Lambda_k}=\{0\}.
\]
So condition \textnormal{(e)} holds as well.
\end{proof}

\begin{definition}[Scheme-refined raw descendant quotient]
\label{def:Main-Vhat}
Fix a ledger step \(k\).
Define the \emph{scheme-refined raw descendant quotient} by
\[
\widehat{\mathcal V}^1_{\Lambda_k}
:=
\frac{\Fsloc}{Q_t(\Fsloc)+Q_u^{\mathrm{scheme}}(\Fsloc)}.
\]
The coarse ambiguity group \(\mathcal V^1\) of
Definition~\ref{def:desc-ambiguity} is the further quotient obtained by killing the
remaining \(Q_u^{\mathrm{phys}}\)-exact shifts.  We write
\[
\pi_{\Lambda_k}:\widehat{\mathcal V}^1_{\Lambda_k}\twoheadrightarrow \mathcal V^1
\]
for that quotient map.
\end{definition}

\begin{remark}[What the refined quotient isolates]
\label{rem:Main-Vhat-scope}
The quotient \(\widehat{\mathcal V}^1_{\Lambda_k}\) is introduced only to separate the genuinely
physical \(Q_u^{\mathrm{phys}}\)-collapse from the already isolated scheme sector.
It does not attempt to incorporate boundary-admissible divergences or local counterterm
coboundaries into the raw descendant group itself.
Those local improvement directions are already handled separately:
they appear in the locality quotient relation
\textnormal{(Definition~\ref{def:locality-quotient})} and, on the characterization side, are
absorbed into \(\mathcal A^-_{\Lambda_k}\) by the characterization supplement.
\end{remark}

\begin{theorem}[Scheme-refined same-tick no-extra-collapse yields the selected-line supplement]
\label{thm:Main-selected-line-contract-refined}
Assume \(\mathbf{2T}_{\mathrm{LS}}+\mathbf{2T}_{\mathrm{char}}^{\mathrm{weak}}\).
Fix a ledger step \(k\), and let
\[
q_{\Lambda_k}:E^-_{\Lambda_k}\longrightarrow E^-_{\Lambda_k}/\mathcal A^-_{\Lambda_k}
\]
be the quotient map.
Assume there exist:
\begin{enumerate}[label=\textnormal{(\alph*)},leftmargin=*,itemsep=0.35em]
\item a one-dimensional selected odd line
  \[
  \widehat L^{\mathrm{sel}}_{\Lambda_k}
  =
  \mathbb R\,\widehat\ell_{\Lambda_k}
  \subset \widehat{\mathcal V}^1_{\Lambda_k};
  \]
\item noncollapse under the coarse quotient,
  \[
  \pi_{\Lambda_k}(\widehat\ell_{\Lambda_k})\neq 0;
  \]
\item a selected representative
  \[
  J^{\mathrm{sel}}_{\Lambda_k}\in[\Xi^{(1)}_{\Lambda_k}]\cap E^-_{\Lambda_k};
  \]
\item a linear realization map
  \[
  \widehat{\mathfrak R}_{\Lambda_k}:\widehat L^{\mathrm{sel}}_{\Lambda_k}
  \longrightarrow
  E^-_{\Lambda_k}/\mathcal A^-_{\Lambda_k}
  \]
  satisfying
  \[
  \widehat{\mathfrak R}_{\Lambda_k}(\widehat\ell_{\Lambda_k})
  =
  q_{\Lambda_k}(J^{\mathrm{sel}}_{\Lambda_k});
  \]
\item exact no-extra-collapse on the selected line,
  \[
  \ker\!\bigl(\widehat{\mathfrak R}_{\Lambda_k}\!\mid_{\widehat L^{\mathrm{sel}}_{\Lambda_k}}\bigr)
  =
  \ker\!\bigl(\pi_{\Lambda_k}\!\mid_{\widehat L^{\mathrm{sel}}_{\Lambda_k}}\bigr).
  \]
\end{enumerate}
Then the selected-line Main supplement package
\(\mathbf{2T}_{\mathrm{sel}}\) of Definition~\ref{def:2T-sel} holds with
\[
\ell_{\Lambda_k}
:=
\pi_{\Lambda_k}(\widehat\ell_{\Lambda_k})
\in\mathcal V^1.
\]
Consequently, Proposition~\ref{prop:Main-selected-line-compatibility} applies.
\end{theorem}

\begin{proof}
By assumption \textnormal{(b)}, the class
\(
\ell_{\Lambda_k}:=\pi_{\Lambda_k}(\widehat\ell_{\Lambda_k})
\)
is nonzero in \(\mathcal V^1\), so the first part of
Definition~\ref{def:2T-sel} holds with
\(
L^{\mathrm{sel}}_{\Lambda_k}=\mathbb R\,\ell_{\Lambda_k}
\).
The second part holds by the chosen representative \(J^{\mathrm{sel}}_{\Lambda_k}\).

It remains to verify the selected-line compatibility implication
\eqref{eq:selected-line-compatibility}.
Assume
\[
q_{\Lambda_k}(J^{\mathrm{sel}}_{\Lambda_k})=0.
\]
Using assumption \textnormal{(d)}, this means
\[
\widehat{\mathfrak R}_{\Lambda_k}(\widehat\ell_{\Lambda_k})=0,
\]
so
\(
\widehat\ell_{\Lambda_k}\in
\ker(\widehat{\mathfrak R}_{\Lambda_k}\!\mid_{\widehat L^{\mathrm{sel}}_{\Lambda_k}})
\).
By assumption \textnormal{(e)},
\[
\widehat\ell_{\Lambda_k}\in
\ker(\pi_{\Lambda_k}\!\mid_{\widehat L^{\mathrm{sel}}_{\Lambda_k}}),
\]
hence
\[
\pi_{\Lambda_k}(\widehat\ell_{\Lambda_k})=0.
\]
By definition of \(\ell_{\Lambda_k}\), this is exactly
\[
[J^{\mathrm{sel}}_{\Lambda_k}]_{\mathcal V^1}=0.
\]
So the selected-line compatibility implication holds, completing the selected-line supplement
contract.

The final sentence follows by Proposition~\ref{prop:Main-selected-line-compatibility}.
\end{proof}

\begin{theorem}[Same-tick faithful LS-to-descendant intertwiner yields a nontrivial square-loop witness]
\label{thm:Main-faithful-LS-desc}
Assume \(\mathbf{2T}_{\mathrm{LS}}+\mathbf{2T}_{\mathrm{char}}^{\mathrm{weak}}\).
Fix a ledger step \(k\), and let
\[
q_{\Lambda_k}:E^-_{\Lambda_k}\longrightarrow E^-_{\Lambda_k}/\mathcal A^-_{\Lambda_k}
\]
be the quotient map.
Assume there exist:
\begin{enumerate}[label=\textnormal{(\alph*)},leftmargin=*,itemsep=0.35em]
\item a one-dimensional selected odd descendant line
  \[
  L^{\mathrm{sel}}_{\Lambda_k}
  =
  \mathbb R\,\ell_{\Lambda_k}
  \subset\mathcal V^1
  \qquad\text{with}\qquad
  \ell_{\Lambda_k}\neq 0;
  \]
\item a linear realization map
  \[
  \mathfrak R_{\Lambda_k}:L^{\mathrm{sel}}_{\Lambda_k}
  \longrightarrow
  E^-_{\Lambda_k}/\mathcal A^-_{\Lambda_k};
  \]
\item a representative \(J_{\Lambda_k}\in[\Xi^{(1)}_{\Lambda_k}]\) such that
  \[
  \mathfrak R_{\Lambda_k}(\ell_{\Lambda_k})
  =
  q_{\Lambda_k}(J_{\Lambda_k});
  \]
\item equivariance of the realization map with respect to the LS square-loop:
  \[
  \mathfrak R_{\Lambda_k}(-\ell_{\Lambda_k})
  =
  -\,\mathfrak R_{\Lambda_k}(\ell_{\Lambda_k});
  \]
\item faithfulness on the selected line:
  \[
  \ker \mathfrak R_{\Lambda_k}=\{0\}.
  \]
\end{enumerate}
Define
\[
\overline{\Xi}^{(1)}_{\Lambda_k}
:=
q_{\Lambda_k}(J)
\qquad
\bigl(J\in[\Xi^{(1)}_{\Lambda_k}]\bigr),
\]
which is well defined because \([\Xi^{(1)}_{\Lambda_k}]\) is an affine class modulo
\(\mathcal A^-_{\Lambda_k}\).
Then
\begin{equation}
\mathrm{Act}^{\mathrm{desc}}_{\sqLoop}\!\big(\overline{\Xi}^{(1)}_{\Lambda_k}\big)
=
-\,\overline{\Xi}^{(1)}_{\Lambda_k}
\neq
\overline{\Xi}^{(1)}_{\Lambda_k}.
\label{eq:Main-faithful-LS-desc}
\end{equation}
Consequently the hypotheses of Theorem~\ref{thm:Main-nonvanishing-from-square-loop}
hold at ledger step \(k\).
\end{theorem}

\begin{proof}
By Lemma~\ref{lem:parity-transport}, the LS square-loop acts by sign on raw
first-descendant classes in \(\mathcal V^1\):
\[
[\Gamma(\mathcal J)] = -[\mathcal J]
\qquad\text{in }\mathcal V^1.
\]
Since \(L^{\mathrm{sel}}_{\Lambda_k}=\mathbb R\,\ell_{\Lambda_k}\subset\mathcal V^1\) is
assumed to be a selected odd descendant line, this sign transport restricts to
\[
\mathrm{Act}^{\mathcal V^1}_{\sqLoop}(\ell_{\Lambda_k})
=
-\,\ell_{\Lambda_k}.
\]

Next, identify the selected quotient point with the image of the chosen generator.
By assumption \textnormal{(c)},
\[
\mathfrak R_{\Lambda_k}(\ell_{\Lambda_k})
=
q_{\Lambda_k}(J_{\Lambda_k})
\]
for some \(J_{\Lambda_k}\in[\Xi^{(1)}_{\Lambda_k}]\).
But by definition of \(\overline{\Xi}^{(1)}_{\Lambda_k}\), the quotient value
\(q_{\Lambda_k}(J)\) is independent of the chosen representative
\(J\in[\Xi^{(1)}_{\Lambda_k}]\). Therefore
\[
\mathfrak R_{\Lambda_k}(\ell_{\Lambda_k})
=
\overline{\Xi}^{(1)}_{\Lambda_k}.
\]

We now prove that this quotient point is nonzero.
Because \(\ell_{\Lambda_k}\neq 0\) by assumption \textnormal{(a)} and
\(\ker\mathfrak R_{\Lambda_k}=\{0\}\) by assumption \textnormal{(e)}, one has
\[
\mathfrak R_{\Lambda_k}(\ell_{\Lambda_k})\neq 0.
\]
Using the identification just proved, this gives
\[
\overline{\Xi}^{(1)}_{\Lambda_k}\neq 0
\qquad\text{in}\qquad
E^-_{\Lambda_k}/\mathcal A^-_{\Lambda_k}.
\]

Next, transport the square-loop action to the quotient.
By assumption \textnormal{(d)},
\[
\mathfrak R_{\Lambda_k}(-\ell_{\Lambda_k})
=
-\,\mathfrak R_{\Lambda_k}(\ell_{\Lambda_k}).
\]
Substituting \(\mathfrak R_{\Lambda_k}(\ell_{\Lambda_k})=\overline{\Xi}^{(1)}_{\Lambda_k}\),
we obtain
\[
\mathrm{Act}^{\mathrm{desc}}_{\sqLoop}\!\big(\overline{\Xi}^{(1)}_{\Lambda_k}\big)
=
-\,\overline{\Xi}^{(1)}_{\Lambda_k}.
\]
Since \(\overline{\Xi}^{(1)}_{\Lambda_k}\neq 0\) and the quotient
\(E^-_{\Lambda_k}/\mathcal A^-_{\Lambda_k}\) is a real vector space, the only fixed
point of \(x\mapsto -x\) is \(x=0\). Hence
\[
-\,\overline{\Xi}^{(1)}_{\Lambda_k}
\neq
\overline{\Xi}^{(1)}_{\Lambda_k}.
\]
This proves \eqref{eq:Main-faithful-LS-desc}.

Finally, compare with Theorem~\ref{thm:Main-nonvanishing-from-square-loop}.
That theorem requires exactly two inputs on the selected quotient point:
sign transport and nontriviality of the square-loop action.
These are precisely what was established above, namely
\[
\mathrm{Act}^{\mathrm{desc}}_{\sqLoop}\!\big(\overline{\Xi}^{(1)}_{\Lambda_k}\big)
=
-\,\overline{\Xi}^{(1)}_{\Lambda_k}
\qquad\text{and}\qquad
\mathrm{Act}^{\mathrm{desc}}_{\sqLoop}\!\big(\overline{\Xi}^{(1)}_{\Lambda_k}\big)
\neq
\overline{\Xi}^{(1)}_{\Lambda_k}.
\]
Therefore the hypotheses of Theorem~\ref{thm:Main-nonvanishing-from-square-loop}
hold at ledger step \(k\).
\end{proof}

\begin{theorem}[Nonvanishing from a same-tick square-loop witness on the transported descendant quotient]
\label{thm:Main-nonvanishing-from-square-loop}
Assume the weak polarized characterization interface
\(\mathbf{2T}_{\mathrm{char}}^{\mathrm{weak}}\)
\textnormal{(Definition~\ref{def:2T-char-weak})}. Fix a ledger step \(k\), and let
\[
q_{\Lambda_k}:E^-_{\Lambda_k}\longrightarrow E^-_{\Lambda_k}/\mathcal A^-_{\Lambda_k}
\]
be the quotient map.
Define the transported descendant quotient point by
\[
\overline{\Xi}^{(1)}_{\Lambda_k}
:=
q_{\Lambda_k}(J)
\qquad
\bigl(J\in[\Xi^{(1)}_{\Lambda_k}]\bigr),
\]
which is well defined because \([\Xi^{(1)}_{\Lambda_k}]\) is an affine class modulo
\(\mathcal A^-_{\Lambda_k}\).
Assume moreover that the abstract LS square-loop transports this selected quotient point by sign:
\begin{equation}
\mathrm{Act}^{\mathrm{desc}}_{\sqLoop}\!\big(\overline{\Xi}^{(1)}_{\Lambda_k}\big)
=
-\,\overline{\Xi}^{(1)}_{\Lambda_k},
\label{eq:Main-square-loop-sign}
\end{equation}
and that this transport is nontrivial on the selected descendant point:
\begin{equation}
\mathrm{Act}^{\mathrm{desc}}_{\sqLoop}\!\big(\overline{\Xi}^{(1)}_{\Lambda_k}\big)
\neq
\overline{\Xi}^{(1)}_{\Lambda_k}.
\label{eq:Main-square-loop-nontrivial}
\end{equation}
Then
\[
0\notin[\Xi^{(1)}_{\Lambda_k}].
\]
\end{theorem}

\begin{proof}
First note that \(\overline{\Xi}^{(1)}_{\Lambda_k}\) is well defined.
Indeed, if \(J_1,J_2\in[\Xi^{(1)}_{\Lambda_k}]\), then
\(
J_2-J_1\in\mathcal A^-_{\Lambda_k}
\),
so
\(
q_{\Lambda_k}(J_1)=q_{\Lambda_k}(J_2)
\)
in the quotient \(E^-_{\Lambda_k}/\mathcal A^-_{\Lambda_k}\).

Assume for contradiction that
\[
0\in[\Xi^{(1)}_{\Lambda_k}].
\]
Then, since \([\Xi^{(1)}_{\Lambda_k}]\) is an affine class modulo \(\mathcal A^-_{\Lambda_k}\), its
quotient point is the zero class:
\[
\overline{\Xi}^{(1)}_{\Lambda_k}
=
q_{\Lambda_k}(0)
=
0
\qquad\text{in}\qquad
E^-_{\Lambda_k}/\mathcal A^-_{\Lambda_k}.
\]
Applying the square-loop transport law \eqref{eq:Main-square-loop-sign} then gives
\[
\mathrm{Act}^{\mathrm{desc}}_{\sqLoop}\!\big(\overline{\Xi}^{(1)}_{\Lambda_k}\big)
=
-\,\overline{\Xi}^{(1)}_{\Lambda_k}
=
-0
=
0
=
\overline{\Xi}^{(1)}_{\Lambda_k},
\]
which contradicts the nontriviality assumption \eqref{eq:Main-square-loop-nontrivial}.
Therefore the assumption \(0\in[\Xi^{(1)}_{\Lambda_k}]\) is false, proving
\[
0\notin[\Xi^{(1)}_{\Lambda_k}].
\]

\end{proof}

\begin{theorem}[Delayed selected line yields a nonzero quotient point]
\label{thm:Main-delayed-selected-line-contract}
Assume \(\mathbf{2T}_{\mathrm{char}}^{\mathrm{weak}}\).
Fix consecutive ledger steps \(k\) and \(k+1\).
Assume there are delayed-lift data
\(
\mathfrak L^{\mathrm{split}}_{k\to k+1},
J^{\mathrm{split}}_{\Lambda_{k+1}},
\overline\rho^{\mathrm{dom}}_{\Lambda_{k+1}}
\),
and a scalar \(c_{\Lambda_{k+1}}\) as displayed below.
In particular, let
\[
\overline\ell^{\mathrm{Main}}_{\Lambda_{k+1}}
:=
\mathfrak L^{\mathrm{split}}_{k\to k+1}(\widehat\ell^{\mathrm{delay}}_k)
=
q_{\Lambda_{k+1}}(J^{\mathrm{split}}_{\Lambda_{k+1}})
\in
\mathfrak D^-_{\Lambda_{k+1}},
\]
with
\[
J^{\mathrm{split}}_{\Lambda_{k+1}}
\in
[\Xi^{(1)}_{\Lambda_{k+1}}]\cap E^-_{\Lambda_{k+1}}.
\]
Assume moreover that
\[
\overline\ell^{\mathrm{Main}}_{\Lambda_{k+1}}\neq 0
\qquad\text{and}\qquad
\overline\rho^{\mathrm{dom}}_{\Lambda_{k+1}}
(\overline\ell^{\mathrm{Main}}_{\Lambda_{k+1}})
=
c_{\Lambda_{k+1}}\neq 0.
\]
Then:
\begin{enumerate}[label=\textnormal{(\roman*)},leftmargin=*,itemsep=0.3em]
\item
  \(
  \overline\ell^{\mathrm{Main}}_{\Lambda_{k+1}}\neq 0
  \);
\item
  \[
  \overline\rho^{\mathrm{dom}}_{\Lambda_{k+1}}
  (\overline\ell^{\mathrm{Main}}_{\Lambda_{k+1}})
  =
  c_{\Lambda_{k+1}}\neq 0;
  \]
\item the map
  \(
  \mathfrak L^{\mathrm{split}}_{k\to k+1}
  \)
  is faithful on
  \(
  \widehat L^{\mathrm{sel,delay}}_k
  \);
\item the delayed selected-line supplement package
  \(\mathbf{2T}_{\mathrm{sel}}^{\mathrm{delay}}\)
  of Definition~\ref{def:2T-sel-delay} holds.
\end{enumerate}
\end{theorem}

\begin{proof}
Items \textnormal{(i)} and \textnormal{(ii)} are assumed.
Since \(\widehat L^{\mathrm{sel,delay}}_k\) is one-dimensional and
\(
\mathfrak L^{\mathrm{split}}_{k\to k+1}
(\widehat\ell^{\mathrm{delay}}_k)
=
\overline\ell^{\mathrm{Main}}_{\Lambda_{k+1}}
\neq 0,
\)
the map \(\mathfrak L^{\mathrm{split}}_{k\to k+1}\) is faithful on
\(\widehat L^{\mathrm{sel,delay}}_k\), proving
\textnormal{(iii)}.
The representative
\(
J^{\mathrm{split}}_{\Lambda_{k+1}}
\)
there satisfies
\[
q_{\Lambda_{k+1}}(J^{\mathrm{split}}_{\Lambda_{k+1}})
=
\overline\ell^{\mathrm{Main}}_{\Lambda_{k+1}}.
\]
So every clause of Definition~\ref{def:2T-sel-delay} is satisfied.
\end{proof}

\begin{theorem}[Nonvanishing from the delayed lift]
\label{thm:Main-nonvanishing-from-delayed-lift}
Assume \(\mathbf{2T}_{\mathrm{char}}^{\mathrm{weak}}\).
Fix consecutive ledger steps \(k\) and \(k+1\).
Assume the delayed selected-line supplement package
\(\mathbf{2T}_{\mathrm{sel}}^{\mathrm{delay}}\)
from Definition~\ref{def:2T-sel-delay}.
Then
\[
0\notin[\Xi^{(1)}_{\Lambda_{k+1}}].
\]
\end{theorem}

\begin{proof}
By Definition~\ref{def:2T-sel-delay}, there exists
\(
J^{\mathrm{split}}_{\Lambda_{k+1}}
\in
[\Xi^{(1)}_{\Lambda_{k+1}}]\cap E^-_{\Lambda_{k+1}}
\)
with
\[
q_{\Lambda_{k+1}}(J^{\mathrm{split}}_{\Lambda_{k+1}})
=
\overline\ell^{\mathrm{Main}}_{\Lambda_{k+1}}.
\]
The same definition gives
\(
\overline\ell^{\mathrm{Main}}_{\Lambda_{k+1}}\neq 0
\).
Assume for contradiction that
\(
0\in[\Xi^{(1)}_{\Lambda_{k+1}}]
\).
Then every representative of the affine class has quotient image zero, hence in particular
\[
q_{\Lambda_{k+1}}(J^{\mathrm{split}}_{\Lambda_{k+1}})=0,
\]
contradicting
\(
\overline\ell^{\mathrm{Main}}_{\Lambda_{k+1}}\neq 0
\).
Therefore
\(
0\notin[\Xi^{(1)}_{\Lambda_{k+1}}]
\).
\end{proof}

\begin{corollary}[Canonical representative is nonzero on the delayed route]
\label{cor:Jcan-nonzero-from-delayed-lift}
Assume \(\mathbf{2T}_{\mathrm{char}}^{\mathrm{cal}}\) and the hypotheses of
Theorem~\ref{thm:Main-nonvanishing-from-delayed-lift}.
Then
\[
\mathcal J^{\mathrm{can}}_{\Lambda_{k+1}}\neq 0.
\]
\end{corollary}

\begin{proof}
By Theorem~\ref{thm:Main-nonvanishing-from-delayed-lift},
\(
0\notin[\Xi^{(1)}_{\Lambda_{k+1}}]
\).
If \(\mathcal J^{\mathrm{can}}_{\Lambda_{k+1}}=0\), then since
\(
\mathcal J^{\mathrm{can}}_{\Lambda_{k+1}}\in[\Xi^{(1)}_{\Lambda_{k+1}}]
\)
by Theorem~\ref{thm:finite-cutoff-Jcan}, one would have
\(
0\in[\Xi^{(1)}_{\Lambda_{k+1}}]
\), contradiction.
\end{proof}

\begin{corollary}[Canonical dominant-kernel reentry on the delayed selected quotient line]
\label{cor:Main-delayed-dominant-kernel-reentry}
Fix consecutive ledger steps \(k\) and \(k+1\).
Assume the hypotheses of
Theorem~\ref{thm:Main-delayed-selected-line-contract}.
Define
\[
N^{\mathrm{dom-null}}_{\Lambda_{k+1}}
:=
\ker \overline\rho^{\mathrm{dom}}_{\Lambda_{k+1}}
\subset
\mathfrak D^-_{\Lambda_{k+1}}.
\]
Assume furthermore that \(N^{\mathrm{dom-null}}_{\Lambda_{k+1}}\) is admissible
as the mod-null datum and satisfies the required transversality condition.
In particular, the mod-null selected line and its faithful realization map are canonically defined at
\(\Lambda_{k+1}\).
\end{corollary}

\begin{proof}
By the additional hypothesis, the dominant-kernel datum
\(
N^{\mathrm{dom-null}}_{\Lambda_{k+1}}
\)
is admissible as the mod-null datum, the required transversality holds,
and the conclusions of the canonical mod-null reentry theorem hold at \(\Lambda_{k+1}\).
\end{proof}

\begin{remark}[Exact boundary of the delayed selected-line lane]
\label{rem:Main-delayed-boundary}
The delayed selected-line lane establishes a quotient-first delayed route:
there is a delayed selected carrier at step \(k\), a faithful delayed lift to a nonzero dominant-active
quotient point in
\(\mathfrak D^-_{\Lambda_{k+1}}\),
and a canonical dominant-kernel reentry lane.
What is not yet derived from the delayed quotient data alone is the existence
of one subordinate unit descendant realization refining that quotient point.
Accordingly, this lane should not be presented as a
fully closed microscopic reconstruction theorem.
\end{remark}

\begin{corollary}[Canonical representative is nonzero once the class is nonzero]
\label{cor:Jcan-nonzero-from-class}
Assume \(\mathbf{2T}_{\mathrm{char}}^{\mathrm{cal}}\) and the hypotheses of
Theorem~\ref{thm:Main-nonvanishing-from-square-loop}. Then
\[
\mathcal J^{\mathrm{can}}_{\Lambda_k}\neq 0.
\]
\end{corollary}

\begin{proof}
By Theorem~\ref{thm:Main-nonvanishing-from-square-loop},
\(
0\notin[\Xi^{(1)}_{\Lambda_k}]
\).
If \(\mathcal J^{\mathrm{can}}_{\Lambda_k}=0\), then since
\(
\mathcal J^{\mathrm{can}}_{\Lambda_k}\in[\Xi^{(1)}_{\Lambda_k}]
\)
by Theorem~\ref{thm:finite-cutoff-Jcan}, one would have
\(
0\in[\Xi^{(1)}_{\Lambda_k}]
\), contradiction.
\end{proof}

\begin{theorem}[The characterization algebra is derived from the canonical representative]
\label{thm:char-cal-to-alg}
Assume \(\mathbf{2T}_{\mathrm{char}}^{\mathrm{cal}}\). For each ledger step \(k\), let
\[
\mathcal O^{\mathrm{char}}_{\Lambda_k}
:=
\mathrm{alg}\Bigl\{\mathcal J^{\mathrm{can}}_{\Lambda_k},
\bigl(\mathcal J^{\mathrm{can}}_{\Lambda_k}\bigr)^2\Bigr\}
\subset \Fsloc.
\]
Then the characterization-algebra supplement
\(\mathbf{2T}_{\mathrm{alg}}\) \textnormal{(Definition~\ref{def:2T-alg})} holds.
\end{theorem}

\begin{proof}
By Theorem~\ref{thm:finite-cutoff-Jcan}, the canonical representative
\(\mathcal J^{\mathrm{can}}_{\Lambda_k}\in E^-_{\Lambda_k}\subset\Fsloc\) is defined for each
ledger step \(k\). Since \(\Fsloc\) is an algebra, the algebra generated by
\(\mathcal J^{\mathrm{can}}_{\Lambda_k}\) and its even square
\(\bigl(\mathcal J^{\mathrm{can}}_{\Lambda_k}\bigr)^2\) is a distinguished subalgebra of
\(\Fsloc\). This is exactly the datum required in Definition~\ref{def:2T-alg}.

\end{proof}

The nonzero-class input used below may now be supplied either by the same-tick
obstruction/square-loop route of Theorem~\ref{thm:Main-nonvanishing-from-square-loop}
or by the delayed lift route of
Theorem~\ref{thm:Main-nonvanishing-from-delayed-lift}; the latter passes through
the delayed \(k\to k+1\) step.

\begin{proposition}[Canonical activity functional from a nonzero odd descendant class]
\label{prop:char-activity-from-nonzero-class}
Assume \(\mathbf{2T}_{\mathrm{char}}^{\mathrm{cal}}\). Fix a ledger step \(k\).
If \(0\notin[\Xi^{(1)}_{\Lambda_k}]\), equivalently
\(\mathcal J^{\mathrm{can}}_{\Lambda_k}\neq 0\), define
\[
\chi^{\mathrm{can}}_{\Lambda_k}(F)
:=
\frac{\mathsf B_{\Lambda_k}(F,\mathcal J^{\mathrm{can}}_{\Lambda_k})}
{\mathsf B_{\Lambda_k}(\mathcal J^{\mathrm{can}}_{\Lambda_k},\mathcal J^{\mathrm{can}}_{\Lambda_k})}
\qquad
(F\in E^-_{\Lambda_k}).
\]
Then \(\chi^{\mathrm{can}}_{\Lambda_k}\) is a continuous nonzero linear functional on
\(E^-_{\Lambda_k}\) satisfying
\[
\chi^{\mathrm{can}}_{\Lambda_k}(A)=0
\qquad
\forall\,A\in\mathcal A^-_{\Lambda_k},
\]
and
\[
\chi^{\mathrm{can}}_{\Lambda_k}(J)=1
\qquad
\forall\,J\in[\Xi^{(1)}_{\Lambda_k}].
\]
In particular, once the class is known to be nonzero, activity/orientation are
derived canonically and require no extra weak-interface input.
\end{proposition}

\begin{proof}
Continuity and linearity are immediate from the bilinearity of \(\mathsf B_{\Lambda_k}\),
and the denominator is strictly positive because
\(\mathcal J^{\mathrm{can}}_{\Lambda_k}\neq 0\) and \(\mathsf B_{\Lambda_k}\) is positive
definite.
For \(A\in\mathcal A^-_{\Lambda_k}\), Proposition~\ref{prop:Jcan-orth} gives
\[
\mathsf B_{\Lambda_k}(\mathcal J^{\mathrm{can}}_{\Lambda_k},A)=0,
\]
hence \(\chi^{\mathrm{can}}_{\Lambda_k}(A)=0\).
Now let \(J\in[\Xi^{(1)}_{\Lambda_k}]\). Since
\(
J-\mathcal J^{\mathrm{can}}_{\Lambda_k}\in\mathcal A^-_{\Lambda_k},
\)
write \(J=\mathcal J^{\mathrm{can}}_{\Lambda_k}+A\) with
\(A\in\mathcal A^-_{\Lambda_k}\). Then
\[
\chi^{\mathrm{can}}_{\Lambda_k}(J)
=
\frac{\mathsf B_{\Lambda_k}(\mathcal J^{\mathrm{can}}_{\Lambda_k},\mathcal J^{\mathrm{can}}_{\Lambda_k}+A)}
{\mathsf B_{\Lambda_k}(\mathcal J^{\mathrm{can}}_{\Lambda_k},\mathcal J^{\mathrm{can}}_{\Lambda_k})}
=
1
\]
by the orthogonality just proved.  In particular
\(\chi^{\mathrm{can}}_{\Lambda_k}(\mathcal J^{\mathrm{can}}_{\Lambda_k})=1\), so the functional is
nonzero.
\end{proof}

\begin{remark}[Canonicalization lives on the descendant class]
\label{rem:class-not-whole-space}
Theorem~\ref{thm:finite-cutoff-Jcan} is intentionally a theorem about the affine
odd descendant class \([\Xi^{(1)}_{\Lambda_k}]\), not about all of
\(E^-_{\Lambda_k}\). Once the class is known to be nonzero,
Proposition~\ref{prop:char-activity-from-nonzero-class} supplies a canonical orientation
functional, but that functional still lives on the already identified affine class and does not by
itself pick a unique element of the full odd generator space.
\end{remark}

\begin{remark}[Selected-line lanes]
\label{rem:Main-activity-gap}
There are two selected-line sublanes.
The same-tick package \(\mathbf{2T}_{\mathrm{sel}}\) is the abstract obstruction lane used
to transport the square-loop sign and rule out extra collapse on a same-tick selected carrier.
The delayed package \(\mathbf{2T}_{\mathrm{sel}}^{\mathrm{delay}}\) is the quotient-first
delayed lane: it yields a nonzero dominant-active quotient point at
\(\Lambda_{k+1}\), feeds the canonical dominant-kernel reentry theorem, and
supplies the delayed odd-to-even route.
Its nondegeneracy clause \(c_{\Lambda_{k+1}}\neq 0\) is built into the delayed package as an
imported hypothesis; deriving that clause from stronger upstream transport/overlap input remains
open, and the bridge-closing existence of one subordinate realization remains open.
\end{remark}

\subsection{Subprogram 3b: RG monotone and Hamiltonized odd derivation}
\label{sec:subprogram3b}

In this subsection we return to the RG monotone package and combine it with the
canonicalized odd data obtained above to produce the Hamiltonized local odd
derivation and its physical consequences on the sufficiently local boundary
algebra.

\begin{definition}[RG monotone]
\label{def:RG-monotone}
Assume $\mathbf{2T}_{\mathrm{RG}}$. We call $\mathcal O^{(0)}(u_0)\in\Fsloc$ a
\emph{good observable} if it is $Q_t$--closed, represents a nontrivial class in
$H^0(\Fsloc,Q_t)$, and defines the monotone via
$M(u)=\langle \mathcal O^{(0)}(u)\rangle_{\rho_u}$.
\par
\end{definition}

\begin{lemma}[RG slope equals the anomaly insertion]
\label{lem:Qu-M}
Assume $\mathbf{2T}_{\mathrm{RG}}$ and the anomalous descent
\eqref{eq:anom-descent}. Then at the UV ceiling,
\[
\partial_u M(u)\big\rvert_{u_0}
\;=\;
-\big\langle \mathcal A^{(1)}(u_0)\big\rangle_{\rho_{u_0}}.
\]
In particular, if $\partial_u M(u)\rvert_{u_0}<0$, then
$\big\langle \mathcal A^{(1)}(u_0)\big\rangle_{\rho_{u_0}}>0$.
\end{lemma}

\begin{proof}
By $\mathbf{2T}_{\mathrm{RG}}$-\textbf{(M2)},
$\partial_u M(u_0)=-\langle (Q_u^{\mathrm{phys}}\mathcal O^{(0)})(u_0)\rangle_{\rho_{u_0}}$.
Taking the expectation value of \eqref{eq:anom-descent} and using
$\langle Q_t X\rangle_{\rho_{u_0}}=0$ gives
$\langle (Q_u^{\mathrm{phys}}\mathcal O^{(0)})(u_0)\rangle_{\rho_{u_0}}
=\langle \mathcal A^{(1)}(u_0)\rangle_{\rho_{u_0}}$.
\end{proof}

\begin{proposition}[Nontrivial descendant sector detected by the RG slope]
\label{prop:Qu-O0-nontrivial}
Assume $\mathbf{2T}_{\mathrm{BRST}}+\mathbf{2T}_{\mathrm{RG}}$. If the monotone
has nontrivial UV slope $\partial_u M(u)\rvert_{u_0}<0$, then the degree-one
cohomology class
\[
\big[\, (Q_u^{\mathrm{phys}}\mathcal O^{(0)})(u_0)\,\big]\ \neq\ 0
\qquad\text{in}\qquad H^1(\Fsloc,Q_t).
\]
\end{proposition}

\begin{proof}
By $\mathbf{2T}_{\mathrm{RG}}$-\textbf{(M2)}, $(Q_u^{\mathrm{phys}}\mathcal O^{(0)})(u_0)$
is $Q_t$--closed (so it defines a class in $H^1(\Fsloc,Q_t)$).
Suppose for contradiction that it is $Q_t$--exact:
$(Q_u^{\mathrm{phys}}\mathcal O^{(0)})(u_0)=Q_tY$ for some $Y\in\Fsloc$.
Then $\mathbf{2T}_{\mathrm{BRST}}$-\textbf{(B2)} gives
$\langle (Q_u^{\mathrm{phys}}\mathcal O^{(0)})(u_0)\rangle_{\rho_{u_0}}
=\langle Q_tY\rangle_{\rho_{u_0}}=0$.
This contradicts Lemma~\ref{lem:Qu-M} under $\partial_u M(u)\rvert_{u_0}<0$.
\end{proof}

\begin{theorem}[Good observable from RG monotone]
\label{thm:good-O0-from-RG}
Under $\mathbf{2T}_{\mathrm{BRST}}+\mathbf{2T}_{\mathrm{RG}}$, there exists a
nontrivial $Q_t$--closed representative $\mathcal O^{(0)}$ of the RG monotone and,
by \eqref{eq:anom-descent}, a raw first-descendant potential \(\Xi^{(1)}\) together
with an anomaly insertion $\mathcal A^{(1)}$ implementing the anomalous first
descent.
Moreover, every sufficiently local representative \(J\in[\Xi^{(1)}]\) is
EB3-admissible in the sense of Definition~\ref{def:EB3-adm-J}
(Lemma~\ref{lem:J-EB3-adm}).
Moreover, if the UV slope is nontrivial ($\partial_u M(u)\rvert_{u_0}<0$), then
$(Q_u^{\mathrm{phys}}\mathcal O^{(0)})(u_0)$ defines a nonzero class in
$H^1(\Fsloc,Q_t)$ (Proposition~\ref{prop:Qu-O0-nontrivial}).
\end{theorem}

\begin{remark}[Cohomological vs.\ physical parity]
\label{rem:two-gradings}
Recall Definition~\ref{def:two-gradings}.  In particular, the first-descendant
data \(\Xi^{(1)}\) have cohomological degree
\(0\) at representative level, whereas the canonical odd representative
\(\mathcal J^{\mathrm{can}}_{\Lambda_k}\) is chosen in the physically odd sector
\(\Gamma(\mathcal J^{\mathrm{can}}_{\Lambda_k})=-\mathcal J^{\mathrm{can}}_{\Lambda_k}\).
\end{remark}

\begin{definition}[Hamiltonized odd derivation]
\label{def:Main-delta}
Assume \(\mathbf{2T}_{\mathrm{Peierls}}+\mathbf{2T}_{\mathrm{char}}^{\mathrm{cal}}\) so the
Peierls bracket is defined on \(\Fsloc\) and the canonical representative
\(\mathcal J^{\mathrm{can}}_{\Lambda_k}\) is available on the ledger-step
skeleton. For a fixed ledger step \(k\), define
\[
\delta_{\Lambda_k}(A) := \{\mathcal J^{\mathrm{can}}_{\Lambda_k},A\}_{\mathrm P,\Lambda_k}
\]
for $A\in\Fsloc$.
\end{definition}

\begin{lemma}[Hamiltonian parity $\Rightarrow$ derivation parity]
\label{lem:Gamma-parity-of-delta}
Assume $\Gamma$ is a Poisson automorphism on $\Fsloc$ (Lemma~\ref{lem:Gamma-Poisson}).
If \(\Gamma(\mathcal J^{\mathrm{can}}_{\Lambda_k})=-\mathcal J^{\mathrm{can}}_{\Lambda_k}\),
then the Hamiltonian derivation
\(\delta_{\Lambda_k}:=\{\mathcal J^{\mathrm{can}}_{\Lambda_k},-\}_{\mathrm P,\Lambda_k}\)
is \(\Gamma\)--odd:
\[
\Gamma\delta_{\Lambda_k}\Gamma^{-1}=-\delta_{\Lambda_k}.
\]
\end{lemma}

\begin{proof}
For $F\in\Fsloc$,
$
\Gamma\delta_{\Lambda_k}(F)
=\Gamma\{\mathcal J^{\mathrm{can}}_{\Lambda_k},F\}_{\mathrm P,\Lambda_k}
=\{\Gamma\mathcal J^{\mathrm{can}}_{\Lambda_k},\Gamma F\}_{\mathrm P,\Lambda_k}
=-\{\mathcal J^{\mathrm{can}}_{\Lambda_k},\Gamma F\}_{\mathrm P,\Lambda_k}
=-\delta_{\Lambda_k}(\Gamma F).
$
\end{proof}

\begin{remark}[Optional Peierls closure identity]
\label{lem:Main-closure-target}
For the chosen first-descendant potential
\(\mathcal J^{\mathrm{can}}_{\Lambda_k}\), one may prove the existence of a translation-direction Hamiltonian
$\mathcal H_\xi\in\Fsloc$ (for some constant Minkowski vector $\xi^\mu$) such
that, modulo Peierls-central terms and the standard $Q_t$-exact/contact
ambiguities in the locality quotient,
\begin{equation}
\label{eq:Main-closure-selfbracket}
\{\mathcal J^{\mathrm{can}}_{\Lambda_k},\mathcal J^{\mathrm{can}}_{\Lambda_k}\}_{\mathrm P,\Lambda_k}\ =\ 2\,\mathcal H_\xi.
\end{equation}
If \eqref{eq:Main-closure-selfbracket} holds, then for all $F\in\Fsloc$,
\[
\delta_{\Lambda_k}^2(F)
\;=\;\tfrac12\,\big\{\{\mathcal J^{\mathrm{can}}_{\Lambda_k},\mathcal J^{\mathrm{can}}_{\Lambda_k}\}_{\mathrm P,\Lambda_k},\,F\big\}_{\mathrm P,\Lambda_k}
\;=\;\{\mathcal H_\xi,F\}_{\mathrm P},
\]
so \(\delta_{\Lambda_k}^2\) acts as the (inner) translation derivation generated by
$\mathcal H_\xi\).
\end{remark}

\begin{proposition}[Characterization activity implies nontriviality of $\delta$ on the BU/AQFT core]
\label{prop:delta-nontrivial}
Assume
\(\mathbf{2T}_{\mathrm{Peierls}}
+\mathbf{2T}_{\mathrm{char}}^{\mathrm{cal}}
+\mathbf{2T}_{\mathrm{char}}^{\mathrm{act}}\).
Fix a ledger step \(k\), and let
\(\mathcal J^{\mathrm{can}}_{\Lambda_k}\in[\Xi^{(1)}_{\Lambda_k}]\)
be the canonical representative of Theorem~\ref{thm:finite-cutoff-Jcan}.
Then there exists an even sufficiently local probe
\[
A\in\mathcal P^+_{\Lambda_k}\subset\Fsloc
\]
such that
\[
\{\mathcal J^{\mathrm{can}}_{\Lambda_k},A\}_{\mathrm P,\Lambda_k}\neq 0.
\]
Equivalently, the Hamiltonized derivation
\(\delta_{\Lambda_k}=\{\mathcal J^{\mathrm{can}}_{\Lambda_k},-\}_{\mathrm P,\Lambda_k}\)
is nontrivial on \(\Fsloc\).  After passage to the Borchers--Uhlmann core
\textnormal{(Definition~\ref{def:BU-core} and Lemma~\ref{lem:delta-net})}, this gives a
nontrivial graded *-derivation on \(\mathfrak A_{\mathrm{fin}}\).
\end{proposition}

\begin{proof}
Because
\(\mathcal J^{\mathrm{can}}_{\Lambda_k}\in[\Xi^{(1)}_{\Lambda_k}]\),
Definition~\ref{def:2T-char-weak}\textnormal{\textbf{(C5)}} gives
\[
\chi_{\Lambda_k}(\mathcal J^{\mathrm{can}}_{\Lambda_k})
=
\varepsilon_{\Lambda_k}\neq 0.
\]
If instead
\[
\{\mathcal J^{\mathrm{can}}_{\Lambda_k},A\}_{\mathrm P,\Lambda_k}=0
\qquad
\forall\,A\in\mathcal P^+_{\Lambda_k},
\]
then Definition~\ref{def:2T-char-act} would force
\[
\chi_{\Lambda_k}(\mathcal J^{\mathrm{can}}_{\Lambda_k})=0,
\]
contradiction. Therefore some
\(
A\in\mathcal P^+_{\Lambda_k}\subset\Fsloc
\)
satisfies
\[
\{\mathcal J^{\mathrm{can}}_{\Lambda_k},A\}_{\mathrm P,\Lambda_k}\neq 0.
\]
For that same \(A\), the BU-core extension from Lemma~\ref{lem:delta-net} satisfies
\[
\delta_{\Lambda_k}(\Phi(A))
=
\Phi\!\bigl(\{\mathcal J^{\mathrm{can}}_{\Lambda_k},A\}_{\mathrm P,\Lambda_k}\bigr).
\]
Since the Borchers--Uhlmann core imposes only the linearity, *-structure, and
locality relations on generator symbols, a nonzero local functional gives a
nonzero generator symbol. Hence \(\delta_{\Lambda_k}\) is nontrivial on
\(\mathfrak A_{\mathrm{fin}}\) as well.
\end{proof}

\begin{lemma}[Derivation and locality of the canonical Hamiltonian action]
\label{lem:Main-delta-derivation}
Assume
\(\mathbf{2T}_{\mathrm{Peierls}}+\mathbf{2T}_{\mathrm{EBC}}+\mathbf{2T}_{\mathrm{char}}^{\mathrm{cal}}\),
and fix the EB3-admissible canonical representative
\(\mathcal J^{\mathrm{can}}_{\Lambda_k}\in E^-_{\Lambda_k}\)
\textnormal{(Theorem~\ref{thm:finite-cutoff-Jcan} and Lemma~\ref{lem:J-EB3-adm})}.
Then
\[
\delta_{\Lambda_k}(F):=\{\mathcal J^{\mathrm{can}}_{\Lambda_k},F\}_{\mathrm P,\Lambda_k}
\]
is a graded derivation on \(\Fsloc\). Moreover, for every double cone \(O\),
\[
\delta_{\Lambda_k}(\Fsloc(O))\subset\Fsloc(O),
\]
so the BU-core assignment
\(\delta_{\Lambda_k}(\Phi(F)):=\Phi(\delta_{\Lambda_k}F)\) from
Definition~\ref{def:2T-EBC}\textnormal{(EB3)} is local on generators.
\end{lemma}

\begin{proof}
By Theorem~\ref{thm:finite-cutoff-Jcan}, the canonical representative
\(\mathcal J^{\mathrm{can}}_{\Lambda_k}\) belongs to the odd sufficiently local class
\(E^-_{\Lambda_k}\subset\Fsloc\), so it is an admissible input for the Peierls
bracket.  Theorem~\ref{thm:cosmo-derived-Peierls} identifies the Peierls bracket as
a graded Poisson bracket on \(\Fsloc\); for fixed
\(\mathcal J^{\mathrm{can}}_{\Lambda_k}\), the Leibniz and Jacobi identities therefore
imply that the Hamiltonian adjoint action
\(
F\mapsto\{\mathcal J^{\mathrm{can}}_{\Lambda_k},F\}_{\mathrm P,\Lambda_k}
\)
is a graded derivation on \(\Fsloc\).

For locality, let \(F\in\Fsloc(O)\). The sufficiently local class is closed under
the Peierls construction in the same locality regime used in
\(\mathbf{2T}_{\mathrm{Peierls}}\)-\textbf{(P5)} and the EB3 estimate, so the
bracket with the fixed local generator \(\mathcal J^{\mathrm{can}}_{\Lambda_k}\) remains
supported in the same double cone \(O\). Hence
\(\delta_{\Lambda_k}(\Fsloc(O))\subset\Fsloc(O)\). The BU-core extension from
Definition~\ref{def:2T-EBC}\textnormal{(EB3)} is therefore local on generators.
\end{proof}

\begin{lemma}[$\Gamma$-oddness of the Hamiltonized derivation]
\label{lem:Main-delta-omega-odd}
Assume \(\mathbf{2T}_{\mathrm{LS}}+\mathbf{2T}_{\mathrm{char}}^{\mathrm{cal}}\) and
Lemma~\ref{lem:Gamma-Poisson}. For the canonical representative
\(\mathcal J^{\mathrm{can}}_{\Lambda_k}\in E^-_{\Lambda_k}\), one has
\[
\Gamma\delta_{\Lambda_k}\Gamma^{-1}=-\delta_{\Lambda_k}.
\]
\end{lemma}

\begin{proof}
With \(\Gamma(\mathcal J^{\mathrm{can}}_{\Lambda_k})=-\mathcal J^{\mathrm{can}}_{\Lambda_k}\), Lemma~\ref{lem:Gamma-Poisson} gives for
all $A\in\Fsloc$:
\[
\Gamma\delta_{\Lambda_k}(A)
=\Gamma\{\mathcal J^{\mathrm{can}}_{\Lambda_k},A\}_{\mathrm P,\Lambda_k}
=\{\Gamma\mathcal J^{\mathrm{can}}_{\Lambda_k},\Gamma A\}_{\mathrm P,\Lambda_k}
=-\{\mathcal J^{\mathrm{can}}_{\Lambda_k},\Gamma A\}_{\mathrm P,\Lambda_k}
=-\delta_{\Lambda_k}(\Gamma A),
\]
so \(\Gamma\delta_{\Lambda_k}\Gamma^{-1}=-\delta_{\Lambda_k}\).
\end{proof}

\begin{theorem}[Peierls Hamiltonization in $\mathbf{2T}_{\mathrm{Main}}^{\mathrm{spine}}$]
\label{thm:Main-Peierls-Hamiltonization}
Assume \(\mathbf{2T}_{\mathrm{Main}}^{\mathrm{spine}}\)
\textnormal{(Definition~\ref{def:2T-Main-spine})}.
Then for each ledger step \(k\), the canonical odd representative
\(\mathcal J^{\mathrm{can}}_{\Lambda_k}\) Hamiltonizes to a well-defined local odd
\[
\delta_{\Lambda_k}
:=
\{\mathcal J^{\mathrm{can}}_{\Lambda_k},-\}_{\mathrm P,\Lambda_k}
\]
on \(\Fsloc\), and \(\delta_{\Lambda_k}\) is physically \(\Gamma\)-odd.
\end{theorem}

\begin{proof}
By Theorem~\ref{thm:finite-cutoff-Jcan}, each ledger step \(k\) carries a canonical
odd representative \(\mathcal J^{\mathrm{can}}_{\Lambda_k}\in E^-_{\Lambda_k}\), and
Lemma~\ref{lem:J-EB3-adm} shows that this representative is EB3-admissible.  The
Hamiltonian action
\(
\delta_{\Lambda_k}=\{\mathcal J^{\mathrm{can}}_{\Lambda_k},-\}_{\mathrm P,\Lambda_k}
\)
is therefore well-defined on \(\Fsloc\), and
Lemma~\ref{lem:Main-delta-derivation} gives both the graded derivation property and
locality.  Lemma~\ref{lem:Main-delta-omega-odd} supplies the parity statement
\(
\Gamma\delta_{\Lambda_k}\Gamma^{-1}=-\delta_{\Lambda_k}
\).
Hence the canonical representative Hamiltonizes to a well-defined local physically
odd derivation, as claimed.
\end{proof}

\begin{proposition}[Physical invariance under ambiguity-triviality]
\label{prop:Main-physical-invariance-ambiguity}
Fix a ledger step \(k\), and let \(J_1,J_2\in[\Xi^{(1)}_{\Lambda_k}]\), so
\(J_2-J_1\in\mathcal A^-_{\Lambda_k}\). Define
\[
\delta^{(i)}_{\Lambda_k}:=\{J_i,-\}_{\mathrm P,\Lambda_k}
\qquad (i=1,2).
\]
If every ambiguity direction \(A\in\mathcal A^-_{\Lambda_k}\) acts trivially on the
physical observable sector, then \(\delta^{(1)}_{\Lambda_k}\) and
\(\delta^{(2)}_{\Lambda_k}\) agree there. More generally, if for every
\(A\in\mathcal A^-_{\Lambda_k}\) and every physical local observable \(F\) the
Peierls bracket \(\{A,F\}_{\mathrm P,\Lambda_k}\) vanishes in the physical quotient
\textnormal{(equivalently: belongs to the physical null ideal)}, then the induced
physical Peierls derivation is independent of the chosen representative in
\([\Xi^{(1)}_{\Lambda_k}]\).
\end{proposition}

\begin{proof}
Since \(J_2-J_1\in\mathcal A^-_{\Lambda_k}\), bilinearity of the Peierls bracket gives
\[
\delta^{(2)}_{\Lambda_k}(F)-\delta^{(1)}_{\Lambda_k}(F)
=\{J_2-J_1,F\}_{\mathrm P,\Lambda_k}.
\]
If every ambiguity direction acts trivially on the physical observable sector, the
right-hand side vanishes there, so the two derivations agree on physical observables.
Under the weaker null-ideal hypothesis, the same identity shows that the difference
vanishes after passage to the physical quotient, hence the induced physical derivation
depends only on the class \([\Xi^{(1)}_{\Lambda_k}]\).
\end{proof}

\begin{remark}[Need for the ambiguity-triviality hypothesis]
The weak package \(\mathbf{2T}_{\mathrm{char}}^{\mathrm{weak}}\) does not by
itself imply ambiguity-triviality on physical observables. A separate ambiguity
supplement may impose ambiguity-triviality on the characterization algebra and
representative-independence of the induced action of
\([\Xi^{(1)}_{\Lambda_k}]\) on the relevant reduced sector.
Proposition~\ref{prop:Main-physical-invariance-ambiguity} is the corresponding
statement after passage to the AQFT/HLS physical sector.
\end{remark}

% ============================================================
\subsection{Subprogram 3c: AQFT net and HLS on the Minkowski boundary}
\label{sec:subprogram3c}

\paragraph{Minkowski boundary chart.}
Cosmology theory provides the boundary local algebraic data and the Peierls
bracket.  For the AQFT/HLS analysis below we work on a chosen Minkowski
boundary chart on which \(\mathbf{2T}_{\mathrm{EBC}}\) is assumed.  No further
details of the cosmology construction enter in this part of the argument.

\begin{definition}[Borchers--Uhlmann local net]
\label{def:BU-net}
Recall Definition~\ref{def:BU-core}. The assignment $O\mapsto\mathfrak A_{\mathrm{BU}}(O)$
defines a net of local *-algebras on $\mathbb R^{1,3}$.  We write $\mathfrak A(O)$ for
the completion in the vacuum representation obtained in Subprogram~1b
(Theorem~\ref{thm:subprogram1b-Wightman}).
\end{definition}

\begin{proposition}[Haag--Kastler properties]
\label{prop:HK}
Assume Theorem~\ref{thm:subprogram1b-Wightman}. Then the net
$O\mapsto\mathfrak A(O)$ satisfies isotony and locality and carries a
Poincar\'e covariant action $\alpha_{(a,\Lambda)}$ with positive-energy
representation in the vacuum GNS space.
\end{proposition}

\begin{lemma}[Parity on the net]
\label{lem:Gamma-net}
The involution $\Gamma$ extends to a *-automorphism of the net,
$\Gamma:\mathfrak A(O)\to\mathfrak A(O)$, and defines a $\mathbb Z_2$ grading
$\mathfrak A=\mathfrak A^{+}\oplus\mathfrak A^{-}$.
\end{lemma}

\begin{lemma}[Local graded derivation]
\label{lem:delta-net}
Assume Theorem~\ref{thm:subprogram1b-Wightman} and
Theorem~\ref{thm:Main-Peierls-Hamiltonization}. Define \(\delta_{\Lambda_k}\) on BU generators
by
\[
\delta_{\Lambda_k}\big(\Phi(F)\big)\ :=\ \Phi\big(\delta_{\Lambda_k} F\big),
\qquad
\delta_{\Lambda_k} F := \{\mathcal J^{\mathrm{can}}_{\Lambda_k},F\}_{\mathrm P,\Lambda_k}.
\]
Then \(\delta_{\Lambda_k}\) extends uniquely to a graded *-derivation on the polynomial core
$\mathfrak A_{\mathrm{fin}}$ (Definition~\ref{def:A-fin}) and is local in the
sense that \(\delta_{\Lambda_k}(\mathfrak A_{\mathrm{fin}}\cap\mathfrak A_{\mathrm{BU}}(O))
\subset \mathfrak A_{\mathrm{BU}}(O)$.
\end{lemma}

\begin{proof}
The Borchers--Uhlmann algebra $\mathfrak A_{\mathrm{BU}}(O)$ is generated by the
symbols $\Phi(F)$ with $F\in\Fsloc(O)$, modulo the defining *-structure
$\Phi(F)^*=\Phi(\overline F)$ and the locality relations for spacelike separated
supports.  We must check that the assignment on generators is compatible with
these relations.

\paragraph{Compatibility with the *-structure.}
Using $\Phi(F)^*=\Phi(\overline F)$ and the definition of $\delta$ on
generators,
\[
\delta_{\Lambda_k}(\Phi(F)^*)
=\delta_{\Lambda_k}(\Phi(\overline F))
=\Phi(\delta_{\Lambda_k}(\overline F)).
\]
Since \(\delta_{\Lambda_k} F=\{\mathcal J^{\mathrm{can}}_{\Lambda_k},F\}_{\mathrm P,\Lambda_k}\) is defined by the (real)
Peierls bracket, it respects complex conjugation on sufficiently local
functionals: \(\delta_{\Lambda_k}(\overline F)=\overline{\delta_{\Lambda_k}F}\).  Therefore
$
\delta_{\Lambda_k}(\Phi(F)^*)=\Phi(\overline{\delta_{\Lambda_k}F})=\delta_{\Lambda_k}(\Phi(F))^*,
$
so \(\delta_{\Lambda_k}\) is a *-derivation on generators.

\paragraph{Compatibility with locality relations.}
Let $F_i\in\Fsloc(O_i)$ with $O_1$ spacelike separated from $O_2$, so the defining
relations impose $[\Phi(F_1),\Phi(F_2)]=0$ in the BU algebra.
By Lemma~\ref{lem:Main-delta-derivation}, \(\delta_{\Lambda_k}\) preserves locality on \(\Fsloc\),
so \(\delta_{\Lambda_k} F_i\in\Fsloc(O_i)\).
Hence \( [\Phi(\delta_{\Lambda_k} F_1),\Phi(F_2)]=[\Phi(F_1),\Phi(\delta_{\Lambda_k} F_2)]=0 \) by the same
BU locality relations, and applying the graded Leibniz rule to the commutator
shows \(\delta_{\Lambda_k}([\Phi(F_1),\Phi(F_2)])=0\).  Thus \(\delta_{\Lambda_k}\) descends to the quotient
defining $\mathfrak A_{\mathrm{BU}}(O)$.

\paragraph{Extension and uniqueness.}
The polynomial core $\mathfrak A_{\mathrm{fin}}$ is generated by finite
polynomials in the generators $\Phi(F)$.  Since \(\delta_{\Lambda_k}\) is fixed on generators
and required to satisfy the graded Leibniz rule, it extends uniquely to a graded
*-derivation on $\mathfrak A_{\mathrm{fin}}$.  Locality of the extension follows
from \(\delta_{\Lambda_k}(\Fsloc(O))\subset\Fsloc(O)\).
\end{proof}

\begin{lemma}[Vacuum invariance on the even subalgebra]
\label{lem:delta-vacuum-invariant}
Assume Theorem~\ref{thm:subprogram1b-Wightman} and
Lemma~\ref{lem:varpi-Gamma-invariant} (so $\varpi\circ\Gamma=\varpi$). Assume
moreover that \(\delta_{\Lambda_k}\) is \(\Gamma\)-odd in the sense that
\(\Gamma\delta_{\Lambda_k}\Gamma^{-1}=-\delta_{\Lambda_k}\) (Lemma~\ref{lem:Main-delta-omega-odd}). Then
\[
\varpi(\delta_{\Lambda_k}(A))=0\qquad\text{for all }A\in\mathfrak A^{+}.
\]
\end{lemma}

\begin{proof}
If $A\in\mathfrak A^{+}$, then $\Gamma(A)=A$, hence
\[
\varpi(\delta_{\Lambda_k}(A))
=\varpi(\Gamma(\delta_{\Lambda_k}(A)))
=\varpi\big((\Gamma\delta_{\Lambda_k}\Gamma^{-1})(\Gamma A)\big)
=\varpi\big(-\delta_{\Lambda_k}(A)\big)
=-\varpi(\delta_{\Lambda_k}(A)).
\]
\end{proof}

\begin{theorem}[Uniform graph bound for the canonical implementer candidate]
\label{thm:Q0-uniform-graph-bound}
Assume Theorem~\ref{thm:subprogram1b-Wightman}, Lemma~\ref{lem:delta-net},
Definition~\ref{def:2T-EBC} \textnormal{(EB3)}, and the fixed-generator
EB3-admissibility package of Definition~\ref{def:EB3-adm-J} for the canonical
representative \(\mathcal J^{\mathrm{can}}_{\Lambda_k}\). Define \(Q_0\) on the
vacuum core \(\mathcal D_0\) by
\[
Q_0\,\pi(B)\Omega := i\,\pi(\delta_{\Lambda_k}(B))\Omega
\qquad(B\in\mathfrak A_{\mathrm{fin}}).
\]
Then there exist an integer \(m\ge 0\) and a constant \(C_*>0\), depending only
on the fixed canonical representative and the vacuum representation, such that
\[
\|Q_0\psi\|
\le
C_*\,\|(1+H)^m\psi\|
\qquad(\psi\in\mathcal D_0).
\]
Equivalently, \(Q_0\) is graph-bounded on \(\mathcal D_0\) by one fixed power of
\((1+H)\).
\end{theorem}

\begin{proof}
Definition~\ref{def:EB3-adm-J} fixes one Sobolev order
\(\sigma>\frac52\) for the chosen canonical generator, and
Lemma~\ref{lem:J-EB3-adm} shows that the corresponding finitely many locality/Sobolev
seminorms are cutoff-uniform for that fixed
\(\mathcal J^{\mathrm{can}}_{\Lambda_k}\). In the cosmology proof of EB3, these
seminorms are the only source of the graph-bound constants, while the power of
\((1+H)\) depends only on the fixed Sobolev order. Hence the regionwise EB3 bounds
for \(\delta_{\Lambda_k}\) may be restated, for this fixed generator, in the form
\[
\|\pi(\delta_{\Lambda_k}(B))\Omega\|
\le
C_*\,\|(1+H)^m\pi(B)\Omega\|
\qquad(B\in\mathfrak A_{\mathrm{fin}})
\]
with one common exponent \(m\) and one common constant \(C_*\). Writing
\(\psi=\pi(B)\Omega\) gives the displayed graph bound for \(Q_0\).
\end{proof}

\begin{proposition}[Implementability from an energy bound]
\label{prop:delta-implementable}
Assume the hypotheses of Theorem~\ref{thm:Q0-uniform-graph-bound}. Then there
exists a densely defined closable operator \(Q\) on \(\mathcal H\) with
\(\Omega\in\mathcal D(Q)\), \(Q\Omega=0\), such that on the core \(\mathcal D_0\) one has
\[
\delta_{\Lambda_k}(A)=i[Q,\pi(A)]_{\mathrm s}\qquad(A\in\mathfrak A_{\mathrm{fin}}).
\]
\end{proposition}

\begin{proof}
Define an operator $Q_0$ on the GNS core $\mathcal D_0=\pi(\mathfrak A_{\mathrm{fin}})\Omega$
by
\[
Q_0\,\pi(B)\Omega \;:=\; i\,\pi(\delta_{\Lambda_k}(B))\Omega
\qquad(B\in\mathfrak A_{\mathrm{fin}}).
\]

\paragraph{Well-definedness.}
If $\pi(B)\Omega=0$, then Theorem~\ref{thm:Q0-uniform-graph-bound} gives
\[
\|\pi(\delta_{\Lambda_k}(B))\Omega\|
\le
C_*\,\|(1+H)^m\,\pi(B)\Omega\|
=0,
\]
so \(\pi(\delta_{\Lambda_k}(B))\Omega=0\) and $Q_0$ is well-defined.

\paragraph{Implementation identity on $\mathcal D_0$.}
Let $A,B\in\mathfrak A_{\mathrm{fin}}$.  Using the definition of $Q_0$ and the
graded Leibniz rule for $\delta$,
\[
\begin{aligned}
[Q_0,\pi(A)]_{\mathrm s}\,\pi(B)\Omega
&=Q_0\,\pi(AB)\Omega-(-1)^{|A|}\pi(A)\,Q_0\,\pi(B)\Omega \\
&=i\,\pi(\delta_{\Lambda_k}(AB))\Omega-(-1)^{|A|}i\,\pi(A\delta_{\Lambda_k}(B))\Omega \\
&=i\,\pi(\delta_{\Lambda_k}(A)B)\Omega
\;=\;i\,\pi(\delta_{\Lambda_k}(A))\,\pi(B)\Omega,
\end{aligned}
\]
so \(\delta_{\Lambda_k}(A)=i[Q_0,\pi(A)]_{\mathrm s}\) on \(\mathcal D_0\).

\paragraph{Closability and extension.}
By Theorem~\ref{thm:Q0-uniform-graph-bound},
\[
\|Q_0\psi\|
\le
C_*\,\|(1+H)^m\psi\|
\qquad(\psi\in\mathcal D_0).
\]
Thus \(Q_0\) is bounded relative to the closed positive selfadjoint operator
\((1+H)^m\) on the invariant core \(\mathcal D_0\). Relative graph-boundedness
therefore implies that \(Q_0\) is closable; let \(Q\) be its closure. Then \(Q\)
is densely defined and closable, contains \(\Omega\) in its domain, and the
implementation identity remains valid on \(\mathcal D_0\).

Finally, since \(1\in\mathfrak A_{\mathrm{fin}}\) and every derivation annihilates
the unit,
\[
Q_0\Omega
=
Q_0\,\pi(1)\Omega
=
i\,\pi(\delta_{\Lambda_k}(1))\Omega
=
0.
\]
Because \(Q\) is the closure of \(Q_0\), it follows that \(\Omega\in\mathcal D(Q)\)
and \(Q\Omega=0\).
\end{proof}

\begin{lemma}[Noether current datum for a Peierls-inner odd derivation]
\label{lem:peierls-noether-current-datum}
Assume Theorem~\ref{thm:subprogram1b-Wightman},
Definition~\ref{def:2T-EBC},
and \(\mathbf{2T}_{\mathrm{Peierls}}\)\textnormal{-(P2)--(P5)}.
Let \(\mathcal J\in\Fsloc(O_{\mathcal J})\) be physically odd and EB3-admissible, and
define the local graded *-derivation
\[
\delta_{\mathcal J}(F):=\{\mathcal J,F\}_{\mathrm P}
\qquad(F\in\Fsloc),
\]
extended to \(\mathfrak A_{\mathrm{fin}}\) by
\(\delta_{\mathcal J}(\Phi(F)):=\Phi(\delta_{\mathcal J}F)\) and the graded Leibniz rule.
Then there exists a local vector-valued operator-valued distribution
\[
j^\mu_{\delta_{\mathcal J}}(x)
\]
on the vacuum core \(\mathcal D_0\) with the following properties.
\begin{enumerate}[label=\textnormal{(\roman*)},leftmargin=*,itemsep=0.35em]
\item \textbf{Locality, *-compatibility, and temperedness.}
  For every compactly supported smooth test vector \(h_\mu\), the smeared operator
  \[
  j_{\delta_{\mathcal J}}(h):=\int d^4x\,h_\mu(x)\,j^\mu_{\delta_{\mathcal J}}(x)
  \]
  is localized in \(\operatorname{supp} h\), acts on \(\mathcal D_0\), and satisfies
  \[
  j_{\delta_{\mathcal J}}(h)^*=j_{\delta_{\mathcal J}}(\overline h)
  \]
  on \(\mathcal D_0\).
  For every \(\psi,\phi\in\mathcal D_0\), the map
  \(h\mapsto \langle \psi,j_{\delta_{\mathcal J}}(h)\phi\rangle\)
  is a tempered distribution of finite order.
\item \textbf{Conservation.}
  \[
  \partial_\mu j^\mu_{\delta_{\mathcal J}}=0
  \]
  in the distributional/Wightman sense.
\item \textbf{Poincar\'e covariance.}
  Writing \(g=(a,\Lambda)\) and
  \((g\cdot h)_\mu(x):=\Lambda_\mu{}^\nu h_\nu(g^{-1}x)\),
  one has on \(\mathcal D_0\)
  \[
  U(g)\,j_{\delta_{\mathcal J}}(h)\,U(g)^{-1}
  =
  j_{\delta_{\mathcal J}}(g\cdot h).
  \]
\item \textbf{Charge-implementation on local vectors.}
  Fix a standard time mollifier \(\chi_\varepsilon\in C_c^\infty(\mathbb R)\) with
  \(\int \chi_\varepsilon(s)\,ds=1\), and let \(f_R\in C_c^\infty(\mathbb R^3)\) be an
  equal-time spatial cutoff family with \(0\le f_R\le 1\) and \(f_R\to 1\) pointwise.
  Define
  \[
  h^\mu_{t,R,\varepsilon}(x):=\delta^\mu_0\,\chi_\varepsilon(x^0-t)\,f_R(\vec x),
  \qquad
  Q_{R,\varepsilon}(t):=j_{\delta_{\mathcal J}}(h_{t,R,\varepsilon}).
  \]
  After the vacuum normalization \(Q_{R,\varepsilon}(t)\Omega=0\), the following holds:
  for every double cone \(O\) and every \(\psi=\pi(B)\Omega\in\mathcal D_0(O)\),
  there exists \(R_0=R_0(O,t,\varepsilon)\) such that for all \(R\ge R_0\),
  \[
  Q_{R,\varepsilon}(t)\psi
  =
  i\,\pi(\delta_{\mathcal J}(B))\Omega.
  \]
  Equivalently, for every local \(A\in\mathfrak A_{\mathrm{fin}}(O)\) and every
  \(\psi\in\mathcal D_0(O)\),
  \[
  i\,[Q_{R,\varepsilon}(t),\pi(A)]_{\mathrm s}\psi
  =
  \pi(\delta_{\mathcal J}(A))\psi
  \qquad(R\gg 1).
  \]
\item \textbf{Uniqueness of the datum.}
  Any other current datum with the same properties differs from
  \(j^\mu_{\delta_{\mathcal J}}\) by a local improvement
  \(
  \partial_\nu B^{[\nu\mu]}
  \)
  and a conserved current whose cutoff charges act trivially on the induced
  derivation.
\end{enumerate}
\end{lemma}

\begin{proof}
The proof has three steps.

\paragraph{Step 1: classical Noether current from the symplectic current.}
By \(\mathbf{2T}_{\mathrm{Peierls}}\)\textnormal{-(P2)--(P4)}, the causal propagator
\(\Delta\) is defined on first functional derivatives and is compatible with the
renormalized symplectic current \(\omega^\mu_{\mathrm{ren}}\) of
\(\mathbf{2T}_{\mathrm{Peierls}}\)\textnormal{-(P3)}.
Set
\[
X_{\mathcal J}:=\Delta\,\mathcal J^{(1)}.
\]
This is the Hamiltonian linearized solution associated with \(\mathcal J\).
For every compactly supported linearized solution \(\eta\), define the classical current
density
\[
c^\mu_{\mathcal J}[\eta](x)
:=
\omega^\mu_{\mathrm{ren}}(X_{\mathcal J},\eta)(x),
\]
with the argument order chosen so that the associated Cauchy-surface charge equals
\(\{\mathcal J,-\}_{\mathrm P}\) rather than \(\{-,\mathcal J\}_{\mathrm P}\).

Because both entries satisfy the linearized equations, conservation of the symplectic
current gives
\[
\partial_\mu c^\mu_{\mathcal J}[\eta]=0.
\]
For every Cauchy surface \(\Sigma\),
\[
\int_\Sigma d\Sigma_\mu\,c^\mu_{\mathcal J}[\eta]
=
\Omega_{\mathrm{ren},\Sigma}(X_{\mathcal J},\eta),
\]
and by the compatibility of \(\Omega_{\mathrm{ren}}\) with \(\Delta\),
\[
\Omega_{\mathrm{ren},\Sigma}(X_{\mathcal J},\eta_F)
=
\{\mathcal J,F\}_{\mathrm P}
\qquad
(\eta_F:=\Delta F^{(1)}).
\]
Thus the classical Cauchy-surface charge of \(c^\mu_{\mathcal J}\) implements exactly
the Peierls derivation \(\delta_{\mathcal J}\).

The locality of \(\omega^\mu_{\mathrm{ren}}\), the retarded/advanced support of \(\Delta\),
and \(\mathcal J\in\Fsloc(O_{\mathcal J})\) imply that \(c^\mu_{\mathcal J}\) is
sufficiently local. By \(\mathbf{2T}_{\mathrm{Peierls}}\)\textnormal{-(P5)}, compactly
smearing \(c^\mu_{\mathcal J}\) against a test vector \(h_\mu\) again produces an element
of \(\Fsloc\).

\paragraph{Step 2: Wightman/GNS realization on the vacuum core.}
For every test vector \(h_\mu\), let \(C_{\mathcal J}(h)\in\Fsloc\) denote the smeared
current functional constructed in Step~1. Theorem~\ref{thm:subprogram1b-Wightman}
realizes every element of \(\Fsloc\) as a local operator on the vacuum GNS core, so set
\[
j_{\delta_{\mathcal J}}(h):=\pi(\Phi(C_{\mathcal J}(h))).
\]
By Definition~\ref{def:2T-EBC}\textnormal{-(EB1)}, every such smeared local field obeys a
polynomial \(H\)-bound on \(\mathcal D_0\); hence it is closable and its matrix elements
against vectors in \(\mathcal D_0\) are tempered distributions of finite order.
This gives item~\textnormal{(i)}.

The distributional conservation in item~\textnormal{(ii)} is the quantum realization of
the classical identity \(\partial_\mu c^\mu_{\mathcal J}=0\): for every test scalar
\(g\in C_c^\infty(\mathbb R^{1,3})\),
\[
j_{\delta_{\mathcal J}}(\partial^\mu g)=0
\]
on \(\mathcal D_0\).
Likewise, the covariance clause of Definition~\ref{def:2T-EBC} and the covariance of
\(\omega^\mu_{\mathrm{ren}}\) imply
\[
U(g)\,j_{\delta_{\mathcal J}}(h)\,U(g)^{-1}
=
j_{\delta_{\mathcal J}}(g\cdot h),
\]
which is item~\textnormal{(iii)}.

\paragraph{Step 3: charge identity on local vectors.}
Let \(O\) be a double cone and \(\psi=\pi(B)\Omega\in \mathcal D_0(O)\).
Fix \(t\) and \(\varepsilon\). Choose \(R\) so large that the shell on which
\(\nabla f_R\) is supported is spacelike to \(O\) after thickening by the time slab
\(\operatorname{supp}\chi_\varepsilon(\cdot-t)\).
Conservation of the current then allows one to deform the cutoff charge
\(Q_{R,\varepsilon}(t)\) to any other Cauchy-surface representative with the same
interior region but no shell contribution against vectors localized in \(O\).
Replacing the cutoff surface by a full Cauchy surface and using the Step~1
surface-charge identity shows that, in vacuum matrix elements against local vectors,
\[
\langle \pi(C)\Omega,Q_{R,\varepsilon}(t)\pi(B)\Omega\rangle
=
\langle \pi(C)\Omega,\, i\,\pi(\delta_{\mathcal J}(B))\Omega\rangle
\]
for every \(\pi(C)\Omega\in\mathcal D_0\).
Because \(\mathcal D_0\) is dense, this proves
\[
Q_{R,\varepsilon}(t)\pi(B)\Omega
=
i\,\pi(\delta_{\mathcal J}(B))\Omega
\qquad(R\gg 1),
\]
after the vacuum normalization \(Q_{R,\varepsilon}(t)\Omega=0\).
Applying the same argument to graded commutators with a local observable \(A\)
gives
\[
i\,[Q_{R,\varepsilon}(t),\pi(A)]_{\mathrm s}\psi
=
\pi(\delta_{\mathcal J}(A))\psi
\qquad(R\gg 1),
\]
which is item~\textnormal{(iv)}.

Finally, if \(\widetilde j^\mu\) is another current datum with the same charge action,
their difference \(k^\mu:=\widetilde j^\mu-j^\mu_{\delta_{\mathcal J}}\) is conserved and
has vanishing cutoff charges on all local vectors. On each double cone, the local
Poincar\'e lemma writes a conserved compactly supported current with vanishing Cauchy
charge as a divergence \(\partial_\nu B^{[\nu\mu]}\), up to a conserved remainder whose
charges act trivially on the induced derivation. This proves item~\textnormal{(v)}.
\end{proof}

\begin{theorem}[Conserved-current implementation theorem on the vacuum core]
\label{thm:odd-derivation-current-reconstruction}
Assume Theorem~\ref{thm:subprogram1b-Wightman}. Let \(\delta\) be a local graded
*-derivation on \(\mathfrak A_{\mathrm{fin}}\). Assume there exists a current datum
\(j^\mu_\delta(x)\) on \(\mathcal D_0\) with the properties listed in
Lemma~\ref{lem:peierls-noether-current-datum}\textnormal{(i)--(v)} for \(\delta\).
Then:
\begin{enumerate}[label=\textnormal{(\roman*)},leftmargin=*,itemsep=0.35em]
\item for every \(\psi\in\mathcal D_0\) and every time \(t\), the equal-time cutoff
  charges \(Q_{R,\varepsilon}(t)\) stabilize for \(R\gg 1\) to a value
  \(Q_\delta(t)\psi\) independent of \(R\) and \(\varepsilon\);
\item the operator \(Q_\delta(t)\) on \(\mathcal D_0\) is densely defined, odd,
  and closable;
\item \(Q_\delta(t)\) is independent of \(t\); we write the common operator simply
  as \(Q_\delta\);
\item on \(\mathcal D_0\),
  \[
  \delta(A)=i[Q_\delta,\pi(A)]_{\mathrm s}
  \qquad(A\in\mathfrak A_{\mathrm{fin}});
  \]
\item \(j^\mu_\delta\) is unique up to local improvements and currents acting
  trivially on the induced derivation.
\end{enumerate}
\end{theorem}

\begin{proof}
\smallskip\noindent\textit{Step 1: stabilization on local vectors.}
Let \(\psi=\pi(B)\Omega\in\mathcal D_0(O)\). By the charge-implementation property of
the current datum, there exists \(R_0\) such that for all \(R\ge R_0\),
\[
Q_{R,\varepsilon}(t)\psi=i\,\pi(\delta(B))\Omega.
\]
Thus the family stabilizes on every vector of the local core \(\mathcal D_0(O)\), hence
on all of \(\mathcal D_0=\bigcup_O\mathcal D_0(O)\).
This proves item~\textnormal{(i)} and defines \(Q_\delta(t)\) on \(\mathcal D_0\).

\smallskip\noindent\textit{Step 2: implementation identity.}
Let \(A,B\in\mathfrak A_{\mathrm{fin}}\). Choose \(R\) so large that the charge datum
applies simultaneously to \(B\) and to \(AB\). Using the graded Leibniz rule and the
vacuum normalization \(Q_{R,\varepsilon}(t)\Omega=0\),
\[
\begin{aligned}
[Q_\delta(t),\pi(A)]_{\mathrm s}\,\pi(B)\Omega
&=
Q_\delta(t)\pi(AB)\Omega-(-1)^{|A|}\pi(A)Q_\delta(t)\pi(B)\Omega \\
&=
i\,\pi(\delta(AB))\Omega-(-1)^{|A|}i\,\pi(A\,\delta(B))\Omega \\
&=
i\,\pi(\delta(A)B)\Omega \\
&=
i\,\pi(\delta(A))\,\pi(B)\Omega.
\end{aligned}
\]
Since such vectors span \(\mathcal D_0\), this gives
\[
\delta(A)=i[Q_\delta(t),\pi(A)]_{\mathrm s}
\quad\text{on }\mathcal D_0.
\]
This is item~\textnormal{(iv)}.

\smallskip\noindent\textit{Step 3: independence of the time slice.}
Let \(t_1,t_2\in\mathbb R\) and \(\psi\in\mathcal D_0(O)\). Choose \(R\) so large that the
shell region of the two cutoffs between the slices \(t_1\) and \(t_2\) is spacelike to
\(O\). Conservation of the current datum implies that the two cutoff charges have the
same action on \(\psi\). Hence
\[
Q_\delta(t_1)\psi=Q_\delta(t_2)\psi
\qquad(\psi\in\mathcal D_0),
\]
which proves item~\textnormal{(iii)}.

\smallskip\noindent\textit{Step 4: closability.}
Fix \(t\) and write \(Q_\delta:=Q_\delta(t)\).
Suppose \(\psi_n\in\mathcal D_0\), \(\psi_n\to 0\), and \(Q_\delta\psi_n\to\eta\) in
\(\mathcal H\).
Take any local vector \(\phi=\pi(C)\Omega\in\mathcal D_0(O)\). For \(R\gg 1\), the same
stabilization argument applied to the adjoint current datum gives a vector
\(Q_\delta^\sharp\phi\in\mathcal H\) such that
\[
\langle \phi,Q_\delta\psi_n\rangle
=
\langle Q_\delta^\sharp\phi,\psi_n\rangle
\]
for all sufficiently large \(n\).
Because \(\psi_n\to 0\), the right-hand side tends to \(0\), so
\[
\langle \phi,\eta\rangle = 0
\qquad\forall\,\phi\in\mathcal D_0.
\]
Since \(\mathcal D_0\) is dense, \(\eta=0\). Therefore \(Q_\delta\) is closable.
This proves item~\textnormal{(ii)}.

\smallskip\noindent\textit{Step 5: uniqueness.}
If \(\widetilde j^\mu_\delta\) is another current datum for the same derivation, then
\(k^\mu:=\widetilde j^\mu_\delta-j^\mu_\delta\) is conserved and its cutoff charges vanish
on \(\mathcal D_0\). By the local Poincar\'e lemma on Minkowski space, \(k^\mu\) is a local
improvement \(\partial_\nu B^{[\nu\mu]}\) plus a conserved current whose charges act
trivially on the induced derivation. This is item~\textnormal{(v)}.
\end{proof}

\begin{corollary}[Canonical current reconstruction for the Main Hamiltonized derivation]
\label{cor:canonical-current-reconstruction}
Assume \(\mathbf{2T}_{\mathrm{Main}}^{\mathrm{spine}}\),
Theorem~\ref{thm:subprogram1b-Wightman},
Theorem~\ref{thm:Main-Peierls-Hamiltonization},
and Lemma~\ref{lem:J-EB3-adm}.
Then the Hamiltonized derivation
\(\delta_{\Lambda_k}=\{\mathcal J^{\mathrm{can}}_{\Lambda_k},-\}_{\mathrm P,\Lambda_k}\)
admits a Poincar\'e-covariant conserved local current
\[
\mathcal J^\mu_{\Lambda_k}(x)
\]
on the vacuum core \(\mathcal D_0\), unique up to local improvements and
derivation-trivial terms, whose equal-time cutoff charges stabilize on \(\mathcal D_0\)
to a densely defined closable odd operator \(Q_{\Lambda_k}\) satisfying
\[
\delta_{\Lambda_k}(A)=i[Q_{\Lambda_k},\pi(A)]_{\mathrm s}
\qquad(A\in\mathfrak A_{\mathrm{fin}})
\]
on \(\mathcal D_0\).
\end{corollary}

\begin{proof}
By Theorem~\ref{thm:Main-Peierls-Hamiltonization}, the canonical representative
\(\mathcal J^{\mathrm{can}}_{\Lambda_k}\) defines a local odd Peierls-inner derivation
\[
\delta_{\Lambda_k}=\{\mathcal J^{\mathrm{can}}_{\Lambda_k},-\}_{\mathrm P,\Lambda_k}.
\]
Lemma~\ref{lem:J-EB3-adm} gives the required EB3-admissibility of the fixed canonical
generator, and \(\mathbf{2T}_{\mathrm{Main}}^{\mathrm{spine}}\) contains both
\(\mathbf{2T}_{\mathrm{Peierls}}\) and \(\mathbf{2T}_{\mathrm{EBC}}\).
Therefore Lemma~\ref{lem:peierls-noether-current-datum} applies and supplies a current
datum \(j^\mu_{\delta_{\Lambda_k}}\) for \(\delta_{\Lambda_k}\).
Applying Theorem~\ref{thm:odd-derivation-current-reconstruction} to that datum yields
the current \(\mathcal J^\mu_{\Lambda_k}\) and the closable operator \(Q_{\Lambda_k}\)
with the stated implementation identity.
\end{proof}

\begin{theorem}[Translation covariance of the Hamiltonized odd derivation]
\label{lem:delta-poincare}
Assume Corollary~\ref{cor:canonical-current-reconstruction} and
Definition~\ref{def:2T-EBC} \textnormal{(Covariance)}. Then
\(\delta_{\Lambda_k}\) commutes with translations on
\(\mathfrak A_{\mathrm{fin}}\):
\[
\alpha_a\circ \delta_{\Lambda_k} \circ \alpha_a^{-1} = \delta_{\Lambda_k}.
\]
\end{theorem}

\begin{proof}
Let \(Q_{\Lambda_k,R,\varepsilon}(t)\) be the cutoff charges built from the current
\(\mathcal J^\mu_{\Lambda_k}\) of
Corollary~\ref{cor:canonical-current-reconstruction}.
By the Poincar\'e covariance of that current,
\[
U(a)\,Q_{\Lambda_k,R,\varepsilon}(t)\,U(a)^{-1}
=
Q_{\Lambda_k,R,\varepsilon}^{(a)}(t+a^0),
\]
where the superscript \({}^{(a)}\) means that the spatial cutoff is translated by
\(\vec a\).

Fix \(\psi\in\mathcal D_0(O)\). Choose \(R\) so large that both the original and
translated cutoff charges have stabilized on \(\psi\). Since the stabilized charge is
independent of the precise cutoff representative and of the time slice,
\[
U(a)\,Q_{\Lambda_k}\,U(a)^{-1}\psi
=
Q_{\Lambda_k}\psi
\qquad(\psi\in\mathcal D_0).
\]
Hence \(U(a)Q_{\Lambda_k}U(a)^{-1}=Q_{\Lambda_k}\) on \(\mathcal D_0\). Therefore, for
every \(A\in\mathfrak A_{\mathrm{fin}}\),
\[
\alpha_a(\delta_{\Lambda_k}(A))
=
i[U(a)Q_{\Lambda_k}U(a)^{-1},\alpha_a(\pi(A))]_{\mathrm s}
=
i[Q_{\Lambda_k},\alpha_a(\pi(A))]_{\mathrm s}
=
\delta_{\Lambda_k}(\alpha_a A).
\]
This is exactly
\(
\alpha_a\circ\delta_{\Lambda_k}\circ\alpha_a^{-1}=\delta_{\Lambda_k}.
\)
\end{proof}

\begin{definition}[Lorentz-covariant family of odd derivations]
\label{def:delta-family}
Let $\delta$ be a local derivation on $\mathfrak A_{\mathrm{fin}}$. For each
proper orthochronous Lorentz transformation $\Lambda\in SO^+(1,3)$ define
\[
\delta_\Lambda:=\alpha_\Lambda\circ\delta\circ\alpha_\Lambda^{-1},
\qquad
\mathscr D:=\mathrm{span}_{\mathbb R}\{\delta_\Lambda:\Lambda\in SO^+(1,3)\}.
\]
We say $\delta$ is \emph{Lorentz-covariant of finite type} if $\mathscr D$ is
finite-dimensional.
\end{definition}

\begin{lemma}[Hyperplane realization of the reconstructed charge]
\label{lem:delta-hyperplane-charge}
Assume Corollary~\ref{cor:canonical-current-reconstruction}. Let
\(\Sigma_0:=\{x^0=0\}\), let \(Q_{\Lambda_k}\) be the equal-time reconstructed global
charge, and let \(\Sigma=g\Sigma_0\) be any oriented affine Cauchy hyperplane
obtained from \(\Sigma_0\) by a proper orthochronous Poincar\'e transformation
\(g=(a,\Lambda)\). For a transported cutoff function on \(\Sigma\), let
\(Q_{R,\Sigma}\) denote the cutoff flux of the conserved current
\(\mathcal J^\mu_{\Lambda_k}\) through \(\Sigma\). Then:
\begin{enumerate}[(i)]
\item \(Q_{R,\Sigma}\) stabilizes on \(\mathcal D_0\) for \(R\gg 1\) to a closable odd
      operator \(Q_{\Lambda_k}(\Sigma)\);
\item \(Q_{\Lambda_k}(\Sigma)\) is independent of \(\Sigma\);
\item \(U(g)Q_{\Lambda_k}(\Sigma_0)U(g)^{-1}=Q_{\Lambda_k}(g\Sigma_0)\) on
      \(\mathcal D_0\).
\end{enumerate}
\end{lemma}

\begin{proof}
Because \(\mathcal J^\mu_{\Lambda_k}\) is Poincar\'e-covariant, the transported
cutoff flux is related to the standard equal-time cutoff charge by
\[
Q_{R,g\Sigma_0}=U(g)\,Q_{R,\Sigma_0}\,U(g)^{-1}
\]
on \(\mathcal D_0\). Since the equal-time charges stabilize on \(\mathcal D_0\) by
Theorem~\ref{thm:odd-derivation-current-reconstruction}, the same is true for the
transported family, proving~(i) and~(iii).

For~(ii), fix a local vector \(\psi=\pi(B)\Omega\in\mathcal D_0(O)\). Choose \(R\)
so large that the cutoff shell between \(\Sigma_0\) and \(\Sigma\) lies spacelike to
\(O\). Exactly the shell argument used in Step~3 of the proof of
Theorem~\ref{thm:odd-derivation-current-reconstruction} then gives
\[
Q_{R,\Sigma}\psi=Q_{R,\Sigma_0}\psi .
\]
Passing to the stabilized values yields
\[
Q_{\Lambda_k}(\Sigma)\psi=Q_{\Lambda_k}(\Sigma_0)\psi .
\]
Since the local vectors span \(\mathcal D_0\), the two operators agree on
\(\mathcal D_0\). As \(\Sigma\) was arbitrary, the reconstructed global charge is
hyperplane-independent.
\end{proof}

\begin{theorem}[Lorentz finite type from the canonical conserved current]
\label{thm:delta-lorentz-finite-type}
Assume Corollary~\ref{cor:canonical-current-reconstruction}. Then the Hamiltonized
odd derivation \(\delta_{\Lambda_k}\) is Lorentz-covariant of finite type in the sense
of Definition~\ref{def:delta-family}. In fact,
\[
\alpha_\Lambda\circ\delta_{\Lambda_k}\circ\alpha_\Lambda^{-1}
=
\delta_{\Lambda_k}
\qquad
(\Lambda\in SO^+(1,3)).
\]
\end{theorem}

\begin{proof}
Let \(\Sigma_0=\{x^0=0\}\). By Lemma~\ref{lem:delta-hyperplane-charge},
\[
U(\Lambda)\,Q_{\Lambda_k}\,U(\Lambda)^{-1}
=
Q_{\Lambda_k}(\Lambda\Sigma_0)
=
Q_{\Lambda_k}(\Sigma_0)
=
Q_{\Lambda_k}
\]
on \(\mathcal D_0\). Therefore for every local observable \(A\),
\[
\begin{aligned}
(\alpha_\Lambda\delta_{\Lambda_k}\alpha_\Lambda^{-1})(A)
&=
i\,[U(\Lambda)Q_{\Lambda_k}U(\Lambda)^{-1},\pi(A)]_{\mathrm s}\\
&=
i\,[Q_{\Lambda_k},\pi(A)]_{\mathrm s}
=
\delta_{\Lambda_k}(A).
\end{aligned}
\]
Hence every Lorentz conjugate derivation equals \(\delta_{\Lambda_k}\) itself, so
\[
\mathscr D
=
\mathrm{span}_{\mathbb R}\{\alpha_\Lambda\delta_{\Lambda_k}\alpha_\Lambda^{-1}\}
=
\mathbb R\,\delta_{\Lambda_k}
\]
is one-dimensional. This is stronger than finite type.
\end{proof}

\begin{lemma}[Covariance of the Lorentz orbit]
\label{lem:delta-lorentz-orbit}
Assume Theorem~\ref{thm:subprogram1b-Wightman} so the net carries a Poincar\'e
action $\alpha_g$. Let $\delta$ be a local derivation on $\mathfrak A_{\mathrm{fin}}$ and
define $\delta_\Lambda$ as in Definition~\ref{def:delta-family}. Then each
$\delta_\Lambda$ is again a local derivation on $\mathfrak A_{\mathrm{fin}}$, and
$\Lambda\mapsto\delta_\Lambda$ defines a representation of $SO^+(1,3)$ on $\mathscr D$
by conjugation.
\end{lemma}

\begin{proof}
Since each $\alpha_\Lambda$ is a *-automorphism of the net, conjugation preserves
the graded Leibniz rule and the *-structure, so $\delta_\Lambda$ is a graded *-derivation.
If $A$ is localized in a double cone $O$, then $\alpha_\Lambda^{-1}(A)$ is localized
in $\Lambda^{-1}O$ and locality of $\delta$ gives
$\delta(\alpha_\Lambda^{-1}(A))$ localized in $\Lambda^{-1}O$; applying $\alpha_\Lambda$
shows $\delta_\Lambda(A)$ is localized in $O$. Thus $\delta_\Lambda$ is local.
Finally,
\[
\delta_{\Lambda_2\Lambda_1}
=\alpha_{\Lambda_2\Lambda_1}\delta\alpha_{\Lambda_2\Lambda_1}^{-1}
=\alpha_{\Lambda_2}\big(\alpha_{\Lambda_1}\delta\alpha_{\Lambda_1}^{-1}\big)\alpha_{\Lambda_2}^{-1}
=\alpha_{\Lambda_2}\delta_{\Lambda_1}\alpha_{\Lambda_2}^{-1},
\]
so $\Lambda\mapsto\delta_\Lambda$ defines a representation by conjugation on the span
$\mathscr D$.
\end{proof}

\begin{theorem}[HLS admissibility of \(\delta\)]
\label{thm:subprogram3b-HLS-hypotheses}
Assume the Main spine package \(\mathbf{2T}_{\mathrm{Main}}^{\mathrm{spine}}\)
\textnormal{(Definition~\ref{def:2T-Main-spine})}. Fix a ledger step \(k\) for which
Subprogram~1b, the Main Hamiltonization chain, the current-reconstruction theorem,
and the UV conformal/common-sector route have been established, namely
Theorem~\ref{thm:subprogram1b-Wightman},
Theorem~\ref{thm:Main-Peierls-Hamiltonization},
Lemma~\ref{lem:varpi-Gamma-invariant},
Corollary~\ref{cor:canonical-current-reconstruction},
Theorem~\ref{lem:delta-poincare},
and Theorem~\ref{thm:delta-lorentz-finite-type}. Then the derivation
\(\delta_{\Lambda_k}\) satisfies the HLS interface
\(\mathbf{2T}_{\mathrm{HLS}}\) \textnormal{(Definition~\ref{def:2T-HLS})}.
\end{theorem}

\begin{proof}
Lemma~\ref{lem:delta-net} gives a densely defined local graded *-derivation on
\(\mathfrak A_{\mathrm{fin}}\).
Proposition~\ref{prop:delta-nontrivial} gives nontriviality of the induced
derivation.
Lemma~\ref{lem:Main-delta-omega-odd} gives the required physical \(\Gamma\)-oddness.
Lemma~\ref{lem:delta-vacuum-invariant} gives vacuum invariance on \(\mathfrak A^{+}\).
Corollary~\ref{cor:canonical-current-reconstruction} gives implementability by a densely
defined closable operator on the common invariant core \(\mathcal D_0\).
Theorem~\ref{lem:delta-poincare} gives translation covariance.
Theorem~\ref{thm:delta-lorentz-finite-type} gives the finite-dimensional Lorentz
covariance datum.
These are exactly the clauses of the HLS interface
\(\mathbf{2T}_{\mathrm{HLS}}\) of Definition~\ref{def:2T-HLS}.
\end{proof}

\begin{theorem}[Haag--\L opusza\'nski--Sohnius (invoked under $\mathbf{2T}_{\mathrm{HLS}}$)
  \cite{HaagLopuszanskiSohnius1975}]
\label{thm:HLS}
Assume the HLS classification theorem applies under $\mathbf{2T}_{\mathrm{HLS}}$.
Let $(\mathcal H,\pi,U,\Omega)$ satisfy the Wightman/AQFT axioms and the spectrum
condition. Suppose there exists a nontrivial $\Gamma$--odd local derivation
$\delta$ obeying $\mathbf{2T}_{\mathrm{HLS}}$. Then there exist:
\begin{itemize}[leftmargin=*,itemsep=0.2em]
\item a finite-dimensional complex odd charge module \(\mathcal V_Q\),
\item a nonzero vector \(v_0\in\mathcal V_Q\),
\item and densely defined operators \(Q_\alpha(v)\), complex-linear in
      \(v\in\mathcal V_Q\), acting on a common invariant core \(\mathcal D_Q\),
\end{itemize}
such that
\begin{enumerate}[label=\textnormal{(\roman*)},leftmargin=*,itemsep=0.3em]
\item the original derivation is implemented by one nonzero element of
  this module:
  \[
  \delta=i[Q(v_0),\cdot]_{\mathrm s}\quad\text{on }\mathcal D_Q,
  \qquad
  Q(v_0):=(Q_1(v_0),Q_2(v_0));
  \]
\item for each \(v\in\mathcal V_Q\), the pair \(Q(v)=(Q_1(v),Q_2(v))\)
  acts on the common invariant core \(\mathcal D_Q\) as a local odd symmetry
  generator and transforms as a Weyl spinor doublet under the Lorentz generators;
\item after choosing any orthonormal basis \(\{e_I\}_{I=1}^N\) of
  \(\mathcal V_Q\) and writing
  \[
  Q_\alpha^I:=Q_\alpha(e_I),\qquad
  \bar Q_{\dot\alpha I}:=\bigl(Q_\alpha^I\bigr)^\dagger,
  \]
  the graded anticommutators close on translations in indexed form:
  \[
  \{Q_\alpha^I,\bar Q_{\dot\beta J}\}
  =2\,\delta^I{}_J\,\sigma^\mu_{\alpha\dot\beta}P_\mu,
  \qquad
  \{Q_\alpha^I,Q_\beta^J\}
  =\{\bar Q_{\dot\alpha I},\bar Q_{\dot\beta J}\}=0;
  \]
\item the internal HLS symmetry sector, when present, acts on the odd
  generators through endomorphisms of \(\mathcal V_Q\); equivalently, in
  any chosen basis it is represented by matrices acting on the index
  \(I\).
\end{enumerate}
\end{theorem}

\begin{theorem}[Indexed super-Poincar\'e seed at the UV boundary]
\label{thm:Main-HLS}
Assume $\mathbf{2T}_{\mathrm{Main}}^{\mathrm{spine}}$, Theorem~\ref{thm:subprogram1b-Wightman},
and Theorem~\ref{thm:subprogram3b-HLS-hypotheses}. Then, by
Theorem~\ref{thm:HLS}, the Hamiltonized odd derivation \(\delta_{\Lambda_k}\) extends
to a finite-dimensional odd charge module \(\mathcal V_Q\) on the
Minkowski UV boundary.  Choosing a basis \(\{e_I\}_{I=1}^N\) of
\(\mathcal V_Q\) and writing \(Q_\alpha^I:=Q_\alpha(e_I)\), the UV
boundary carries densely defined odd generators
\[
Q_\alpha^I,\qquad \bar Q_{\dot\alpha I},
\]
on a common invariant core \(\mathcal D_Q\) such that
\[
\{Q_\alpha^I,\bar Q_{\dot\beta J}\}
=2\,\delta^I{}_J\,\sigma^\mu_{\alpha\dot\beta}P_\mu,
\qquad
\{Q_\alpha^I,Q_\beta^J\}
=\{\bar Q_{\dot\alpha I},\bar Q_{\dot\beta J}\}=0,
\qquad
[P_\mu,Q_\alpha^I]=[P_\mu,\bar Q_{\dot\alpha I}]=0,
\]
with \(Q_\alpha^I\) and \(\bar Q_{\dot\alpha I}\) transforming in the
two Weyl spinor representations under \(M_{\mu\nu}\).  The internal HLS
symmetry sector acts on this odd
charge multiplet through endomorphisms of \(\mathcal V_Q\). Every module
element \(Q(v)\) acts on the same invariant core \(\mathcal D_Q\) as a local odd
symmetry generator.
\end{theorem}

\begin{proof}
Theorem~\ref{thm:subprogram3b-HLS-hypotheses} proves that the Hamiltonized
derivation \(\delta_{\Lambda_k}\) satisfies the interface
\(\mathbf{2T}_{\mathrm{HLS}}\). Hence the invoked HLS classification theorem
\textnormal{(Theorem~\ref{thm:HLS})} applies. Choosing an orthonormal basis
\(\{e_I\}_{I=1}^N\) of the resulting odd charge module \(\mathcal V_Q\) and writing
\(Q_\alpha^I:=Q_\alpha(e_I)\) gives the indexed generators on the common invariant
core \(\mathcal D_Q\), together with the displayed translation anticommutators,
vanishing same-chirality anticommutators, and Weyl-spinor Lorentz transformation
law. The final sentence is the corresponding module statement from
Theorem~\ref{thm:HLS}.
\end{proof}

\begin{lemma}[Injectivity of the HLS module realization]
\label{lem:HLS-module-injective}
Assume Theorem~\ref{thm:Main-HLS}. If \(Q(v)=0\) as a Weyl pair, then \(v=0\).
\end{lemma}

\begin{proof}
By the invoked HLS theorem used in Theorem~\ref{thm:Main-HLS}, the original
nonzero odd derivation is implemented by some nonzero \(v_0\in\mathcal V_Q\); in
particular \(Q(v_0)\neq 0\) on \(\mathcal D_Q\).

We first show that the translation representation is nontrivial. If all
translation generators vanished, then \(P_\mu=0\) for all \(\mu\), hence also
\(P_{\alpha\dot\beta}=0\) for all spinor indices. The HLS anticommutator would
then give
\[
0
=
\{Q_\alpha(v_0),Q_\alpha(v_0)^\dagger\}
\]
on \(\mathcal D_Q\). Therefore, for every \(\psi\in\mathcal D_Q\),
\[
0
=
\langle\psi,\{Q_\alpha(v_0),Q_\alpha(v_0)^\dagger\}\psi\rangle
=
\|Q_\alpha(v_0)\psi\|^2+\|Q_\alpha(v_0)^\dagger\psi\|^2.
\]
Hence \(Q_\alpha(v_0)=0\) for both \(\alpha=1,2\), contradicting that \(Q(v_0)\)
implements a nontrivial derivation. So not all \(P_\mu\) vanish. Since the
matrices \(\sigma^\mu_{\alpha\dot\beta}\) span the Hermitian \(2\times2\)
matrices, some \(P_{\alpha\dot\beta}\) is nonzero.

Now let \(Q(v)=0\). For every \(w\in\mathcal V_Q\), the indexed HLS
anticommutator gives
\[
0
=
\{Q_\alpha(v),\bar Q_{\dot\beta}(w)\}
=
2\,\langle v,w\rangle_Q\,P_{\alpha\dot\beta}
\]
on \(\mathcal D_Q\). Choose \(\alpha,\dot\beta\) such that
\(P_{\alpha\dot\beta}\neq 0\). Then \(\langle v,w\rangle_Q=0\) for every \(w\),
and therefore \(v=0\).
\end{proof}

\begin{theorem}[HLS spinor completeness on the UV vacuum representation]
\label{thm:HLS-spinor-completeness}
Assume Theorem~\ref{thm:Main-HLS}, and assume the invoked HLS classification step
identifies the full translation-invariant local odd Weyl-spinor sector of the
given chirality with the HLS charge module \(\mathcal V_Q\), in the standard
indexed supercharge-module presentation and classification picture recorded in
\cite{Sohnius:1985qm,Nahm1978}. Let
\(\mathscr Q_+\) be the space of all translation-invariant local odd Weyl-spinor
symmetry generators of the given chirality on the vacuum representation
\((\mathcal H,\pi,U,\Omega)\), and define \(\mathscr Q_-\) similarly for the
conjugate chirality. Then the HLS charge module \(\mathcal V_Q\) is complete in the
sense that for every \(X_\alpha\in\mathscr Q_+\) there exists a unique
\(v\in\mathcal V_Q\) with
\[
X_\alpha = Q_\alpha(v),
\]
and similarly every barred generator in \(\mathscr Q_-\) is a unique linear
combination of the \(\bar Q_{\dot\alpha}(v)\).
\end{theorem}

\begin{proof}
Theorem~\ref{thm:Main-HLS} already identifies a finite-dimensional HLS odd charge
module \(\mathcal V_Q\) whose elements act as translation-invariant local odd Weyl
spinor symmetry generators on the same vacuum representation. The invoked HLS
classification input states that this finite-dimensional odd module
exhausts the corresponding translation-invariant local odd spinor sector. Hence
every such generator \(X_\alpha\) lies in the span of the \(Q_\alpha(v)\).
Uniqueness is not merely basis independence: if \(Q(v)=Q(v')\), then
\(Q(v-v')=0\), so Lemma~\ref{lem:HLS-module-injective} gives \(v=v'\). The barred
statement is identical.
\end{proof}

\begin{theorem}[Optional indexed current reconstruction under separate current data]
\label{thm:indexed-current-reconstruction}
Assume Theorem~\ref{thm:Main-HLS}. Choose a basis \(\{e_I\}_{I=1}^N\) of
\(\mathcal V_Q\) and write
\(
Q_\alpha^I:=Q_\alpha(e_I)
\),
\(
\bar Q_{\dot\alpha I}:=(Q_\alpha^I)^\dagger
\).
Assume moreover that for every \(I\) and \(\alpha\) there is given a conserved local
current datum
\[
\mathcal J^\mu{}_\alpha{}^I(x)
\]
whose equal-time cutoff charges stabilize to \(Q_\alpha^I\) on a dense invariant core
\(\mathcal D_{\mathrm{curr}}^{I,\alpha}\subset\mathcal D_Q\), and whose ambiguity is only
by local improvements and derivation-trivial terms. Then, after replacing
\(\mathcal D_{\mathrm{curr}}^{I,\alpha}\) by the finite intersection over \(I,\alpha\), the
whole indexed family may be realized on one dense common current domain.
\end{theorem}

\begin{proof}
This is exactly the additional datum assumed in the statement. Because the family
\(\{(I,\alpha)\}\) is finite, intersecting the dense invariant cores
\(\mathcal D_{\mathrm{curr}}^{I,\alpha}\) again yields a dense invariant domain, and the
stated ambiguity for each current remains the standard one by improvements and
derivation-trivial terms.
\end{proof}

\begin{remark}[Optional indexed current reconstruction]
Theorem~\ref{thm:indexed-current-reconstruction} is an optional AQFT
supplement. If separate current data are chosen for the HLS generators, it
yields a common dense current domain for the indexed family. A separate indexed
current theorem therefore requires its own current data for those generators.
\end{remark}

\begin{corollary}[Same-\(\mathcal H\) superconformal closure]
\label{cor:Main-postHLS-superconf}
Assume Theorem~\ref{thm:Main-HLS},
Theorem~\ref{thm:HLS-spinor-completeness},
the UV conformal common-sector interface \(\mathbf{2T}_{\mathrm{conf}}\), and the
same-\(\mathcal H\) conformal-closure results of Subprograms~4a/4b. Then on the same UV Hilbert space
\(\mathcal H\) the generators
\[
P_\mu,\ M_{\mu\nu},\ D,\ K_\mu,\ Q_\alpha^I,\ \bar Q_{\dot\alpha I},\
S^\alpha{}_I,\ \bar S^{\dot\alpha I},\ R^I{}_J
\]
close on a dense common invariant core to the full \(4\)D superconformal algebra.
More precisely, the resulting operator Lie superalgebra is
\(\mathfrak{su}(2,2|N)\) for \(N\neq 4\), while in the \(N=4\) case the standard
simple form is obtained after quotienting the central \(\mathfrak u(1)\), giving
\(\mathfrak{psu}(2,2|4)\).
\end{corollary}

\begin{proof}
Combine Theorem~\ref{thm:Main-HLS},
Theorem~\ref{thm:HLS-spinor-completeness},
Theorem~\ref{thm:4b-narrow-common-sector}, and
Theorem~\ref{thm:4d-SConf}.
\end{proof}
